% L’Etat et le pouvoir — Philosophie 1re Tech — Cours signé OnBuch+
% Programme officiel MINESEC. Compiler avec : tectonic N10-etat-pouvoir.tex
\newcommand{\DOCMATIERE}{Philosophie}
\newcommand{\DOCNIVEAU}{1re Tech}
\newcommand{\DOCMODULE}{Philosophie – classe de Première technique}
\newcommand{\DOCLECON}{N10}
\newcommand{\DOCTITRE}{L’Etat et le pouvoir}
% OnBuch+ — préambule « cours signés » ENRICHI (compilé avec tectonic / XeLaTeX)
% Système visuel « Pop » d'OnBuch+ : fond crème (jamais blanc), bordures encre
% épaisses, ombres « dures » décalées, titres Archivo Black en capitales.
% Version MATHÉMATIQUES : graphes (pgfplots/tikz), boîtes pédagogiques
% (définition, propriété, méthode, attention, le savais-tu, exercice résolu,
% exercice, TP, bilan), page de garde et signature OnBuch+ ancrée en bas.
%
% Macros attendues dans main.tex AVANT \input{preamble} :
%   \DOCMATIERE  \DOCNIVEAU  \DOCMODULE  \DOCLECON (numéro)  \DOCTITRE
\documentclass[11pt]{article}

\usepackage{fontspec}
\usepackage[french]{babel}
\usepackage[a4paper,margin=2cm,top=3.4cm,bottom=2.8cm,headheight=1.4cm,footskip=1.3cm]{geometry}
\usepackage{amsmath,amssymb,amsfonts}
\usepackage{mathtools} % \xmapsto, \coloneqq... souvent utilisés par les modèles
\usepackage{cancel} % simplifications barrées (bilans de forces)
% \abs{z} (module, valeur absolue) et \norm{u} (norme), utilisés par les modèles
\providecommand{\abs}{}\let\abs\relax\DeclarePairedDelimiter{\abs}{\lvert}{\rvert}
\providecommand{\norm}{}\let\norm\relax\DeclarePairedDelimiter{\norm}{\lVert}{\rVert}
\usepackage[table]{xcolor} % [table] : fournit aussi \rowcolors (alternance de lignes)
\usepackage{pagecolor}
\usepackage{tikz}
\usetikzlibrary{arrows.meta,positioning,calc,shapes.geometric,shapes.misc,shapes.arrows,decorations.pathmorphing,decorations.pathreplacing,decorations.markings,patterns,shadows,mindmap,trees,fit,backgrounds,matrix,intersections,angles,quotes}
\usepackage{pgfplots}
\usepackage[version=4]{mhchem} % équations nucléaires \ce{^{235}_{92}U}
\usepackage[european,siunitx]{circuitikz} % schémas électriques
\pgfplotsset{compat=1.18}
\usepgfplotslibrary{fillbetween,groupplots} % aires entre courbes (calcul intégral)
\usepackage{enumitem}
\usepackage{fancyhdr}
% [most] charge toutes les bibliothèques tcolorbox (listings, minted, raster,
% documentation...) et gonfle inutilement la mémoire TeX sur un document long
% (~300 boîtes) : on ne charge que ce qui est réellement utilisé.
\usepackage[skins,breakable]{tcolorbox}
\usepackage{graphicx}
\usepackage{colortbl}
\usepackage{booktabs}
\usepackage{tabularx}
\usepackage{multirow}
\usepackage{diagbox} % case d'en-tête diagonale des tableaux à double entrée
% Colonnes extensibles centrée / gauche / droite (C, L, R), souvent utilisées par les modèles
\newcolumntype{C}{>{\centering\arraybackslash}X}
\newcolumntype{L}{>{\raggedright\arraybackslash}X}
\newcolumntype{R}{>{\raggedleft\arraybackslash}X}
\usepackage{array}
\usepackage{multicol}
\usepackage{siunitx}
% parse-numbers=false : accepte aussi \SI{2,5 \times 10^{-3}}{...}
\sisetup{locale=FR,output-decimal-marker={,},group-minimum-digits=4,parse-numbers=false}
\DeclareSIUnit\ohm{\ensuremath{\symup{\Omega}}} % U+2126 absent de la police
\DeclareSIUnit\division{div} % réglages d'oscilloscope (ms/div, V/div)
\DeclareSIUnit\tr{tr} % tours (vitesse de rotation, ex. \SI{1500}{\tr\per\minute})
\DeclareSIUnit\molar{M} % concentration molaire (ex. \SI{3}{\molar}, \SI{1}{\micro\molar})
\DeclareSIUnit\an{an} % année (le modèle confond parfois avec \year, un compteur TeX) — ex. \SI{5}{\kilo\meter\per\an}
\DeclareSIUnit\FCFA{FCFA} % franc CFA (ex. \SI{500}{\FCFA\per\kilo\gram})
\DeclareSIUnit\degreeBrix{\textdegree Brix} % teneur en sucre d'un sirop/jus (réfractométrie)
\DeclareSIUnit\month{mois} % le modèle utilise parfois \month, un compteur TeX, comme unité de durée
\DeclareSIUnit\microliter{\micro\liter} % le modèle écrit parfois \microliter en un seul token
\DeclareSIUnit\million{millions} % dénombrement (ex. \SI{15}{\million\per\mL} spermatozoïdes)
\usepackage{qrcode}
% \degree : commande fréquemment utilisée par les modèles pour un angle ou une
% température en dehors de \SI{}{} (non fournie par les paquets chargés).
\providecommand{\degree}{^{\circ}}
% \qed (fin de démonstration), utilisé par les modèles sans amsthm
\providecommand{\qed}{\hfill$\square$}
% \barre : double barre des valeurs interdites dans un tableau de variations (tkz-tab)
\providecommand{\barre}{\ensuremath{\|}}
% environnement proof (amsthm non chargé) utilisé par les modèles
\ifdefined\proof\else\newenvironment{proof}[1][Démonstration]{\par\noindent\textit{#1.}\ }{\qed\par}\fi
\usepackage{lastpage}
\usepackage{hyperref}
\hypersetup{hidelinks}

% ============================================================
% Palette « Pop » OnBuch+ — mêmes hex que la classe P (plus_ui.dart)
% ============================================================
\definecolor{poporange}{HTML}{F59321}
\definecolor{popdark}{HTML}{E07A0C}
\definecolor{popcream}{HTML}{FFF6E8}
\definecolor{popink}{HTML}{1C1714}
\definecolor{poppurple}{HTML}{7A5AE0}
\definecolor{popgreen}{HTML}{2E9E6A}
\definecolor{poppink}{HTML}{E0457B}
\definecolor{popblue}{HTML}{2F7DE1}
\definecolor{popgold}{HTML}{C98A12}
\definecolor{popmuted}{HTML}{8A7F76}
\definecolor{popline}{HTML}{EADFD2}
\definecolor{popgray}{HTML}{8A7F76}
\colorlet{popgrey}{popgray}
% Alias tolérants (le modèle nomme parfois les couleurs différemment)
\colorlet{popwhite}{white}
\colorlet{popblack}{popink}
\colorlet{popred}{poppink}
\colorlet{popyellow}{popgold}
% Teintes claires (fonds de tableaux/encadrés) — le modèle nomme parfois
% les couleurs avec un suffixe L (ex. popblueL) sans les avoir définies.
\colorlet{poporangeL}{poporange!20}
\colorlet{popdarkL}{popdark!20}
\colorlet{poppurpleL}{poppurple!20}
\colorlet{popgreenL}{popgreen!20}
\colorlet{poppinkL}{poppink!20}
\colorlet{popblueL}{popblue!20}
\colorlet{popgoldL}{popgold!20}
\colorlet{popredL}{popred!20}
\colorlet{popgrayL}{popgray!20}
% Teintes claires pour fonds de tableaux / zones
\colorlet{popblueL}{popblue!10!white}
\colorlet{poporangeL}{poporange!14!white}
\colorlet{popgreenL}{popgreen!12!white}
\colorlet{poppurpleL}{poppurple!12!white}
\colorlet{poppinkL}{poppink!12!white}
\colorlet{popgoldL}{popgold!14!white}

\pagecolor{popcream}

% ============================================================
% Typographie
% ============================================================
\defaultfontfeatures{Ligatures=TeX}
\setmainfont{PlusJakartaSans-Medium.ttf}[Path=../../../fonts/,BoldFont=PlusJakartaSans-Bold.ttf]
% Maths en sans-serif (Fira Math) avec chiffres et lettres droites de Plus
% Jakarta, pour que les formules se fondent dans le texte.
\usepackage[math-style=ISO,bold-style=ISO]{unicode-math}
\setmathfont{FiraMath-Regular.otf}[Path=../../../fonts/]
\setmathfont{PlusJakartaSans-Medium.ttf}[Path=../../../fonts/,range={up/num,up/{Latin,latin},\mathup}]
\usepackage{icomma} % virgule décimale sans espace en mode math
% Caractères mathématiques Unicode absents des polices de texte (Plus Jakarta,
% Archivo Black) : rendus par leur équivalent mathématique, en texte comme en
% formule — sinon ils s'affichent en carré vide (ex. « Arithmétique dans ℤ »).
\usepackage{newunicodechar}
\newunicodechar{ℝ}{\ensuremath{\mathbb{R}}}
\newunicodechar{ℤ}{\ensuremath{\mathbb{Z}}}
\newunicodechar{ℂ}{\ensuremath{\mathbb{C}}}
\newunicodechar{ℕ}{\ensuremath{\mathbb{N}}}
\newunicodechar{ℚ}{\ensuremath{\mathbb{Q}}}
\newunicodechar{𝔻}{\ensuremath{\mathbb{D}}}
\newunicodechar{α}{\ensuremath{\alpha}}
\newunicodechar{β}{\ensuremath{\beta}}
\newunicodechar{γ}{\ensuremath{\gamma}}
\newunicodechar{δ}{\ensuremath{\delta}}
\newunicodechar{ε}{\ensuremath{\varepsilon}}
\newunicodechar{θ}{\ensuremath{\theta}}
\newunicodechar{λ}{\ensuremath{\lambda}}
\newunicodechar{μ}{\ensuremath{\mu}}
\newunicodechar{σ}{\ensuremath{\sigma}}
\newunicodechar{τ}{\ensuremath{\tau}}
\newunicodechar{φ}{\ensuremath{\varphi}}
\newunicodechar{ψ}{\ensuremath{\psi}}
\newunicodechar{ω}{\ensuremath{\omega}}
\newunicodechar{Σ}{\ensuremath{\Sigma}}
% majuscules produites par \MakeUppercase dans les titres (α-aminés -> Α)
\newunicodechar{Α}{\ensuremath{\alpha}}
\newunicodechar{Β}{\ensuremath{\beta}}
\newunicodechar{Γ}{\ensuremath{\Gamma}}
\newunicodechar{Δ}{\ensuremath{\Delta}}
\newunicodechar{Ε}{\ensuremath{\varepsilon}}
\newunicodechar{Θ}{\ensuremath{\Theta}}
\newunicodechar{Λ}{\ensuremath{\Lambda}}
\newunicodechar{Μ}{\ensuremath{\mu}}
\newunicodechar{Τ}{\ensuremath{\tau}}
\newunicodechar{Φ}{\ensuremath{\Phi}}
\newunicodechar{Ψ}{\ensuremath{\Psi}}
\newunicodechar{Ω}{\ensuremath{\Omega}}
\newunicodechar{↦}{\ensuremath{\mapsto}}
\newunicodechar{∈}{\ensuremath{\in}}
\newunicodechar{∉}{\ensuremath{\notin}}
\newunicodechar{∧}{\ensuremath{\wedge}}
\newunicodechar{⇔}{\ensuremath{\Leftrightarrow}}
\newunicodechar{⟺}{\ensuremath{\Longleftrightarrow}}
\newunicodechar{⇒}{\ensuremath{\Rightarrow}}
\newunicodechar{⟹}{\ensuremath{\Longrightarrow}}
\newunicodechar{∘}{\ensuremath{\circ}}
\newunicodechar{∪}{\ensuremath{\cup}}
\newunicodechar{∩}{\ensuremath{\cap}}
\newunicodechar{≡}{\ensuremath{\equiv}}
\newunicodechar{⊂}{\ensuremath{\subset}}
\newunicodechar{⊥}{\ensuremath{\perp}}
\newunicodechar{∃}{\ensuremath{\exists}}
\newunicodechar{∀}{\ensuremath{\forall}}
\newunicodechar{∛}{\ensuremath{\sqrt[3]{\,}}}
\newunicodechar{∜}{\ensuremath{\sqrt[4]{\,}}}
\newunicodechar{⃗}{\ensuremath{{}^{\to}}}
\newunicodechar{ⁿ}{\ifmmode^{n}\else\textsuperscript{\ensuremath{n}}\fi}
\newunicodechar{ˣ}{\ifmmode^{x}\else\textsuperscript{\ensuremath{x}}\fi}
\newunicodechar{ᵇ}{\ifmmode^{b}\else\textsuperscript{\ensuremath{b}}\fi}
\newunicodechar{ʳ}{\ifmmode^{r}\else\textsuperscript{\ensuremath{r}}\fi}
\newunicodechar{ᵘ}{\ifmmode^{u}\else\textsuperscript{\ensuremath{u}}\fi}
\newunicodechar{ʸ}{\ifmmode^{y}\else\textsuperscript{\ensuremath{y}}\fi}
\newunicodechar{ᵃ}{\ifmmode^{a}\else\textsuperscript{\ensuremath{a}}\fi}
\newunicodechar{ᵉ}{\ifmmode^{e}\else\textsuperscript{\ensuremath{e}}\fi}
\newunicodechar{ᵏ}{\ifmmode^{k}\else\textsuperscript{\ensuremath{k}}\fi}
\newunicodechar{ᵖ}{\ifmmode^{p}\else\textsuperscript{\ensuremath{p}}\fi}
\newunicodechar{ᴺ}{\ifmmode^{N}\else\textsuperscript{\ensuremath{N}}\fi}
\newunicodechar{ᶜ}{\ifmmode^{c}\else\textsuperscript{\ensuremath{c}}\fi}
\newunicodechar{ʰ}{\ifmmode^{h}\else\textsuperscript{\ensuremath{h}}\fi}
\newunicodechar{ᵠ}{\ifmmode^{q}\else\textsuperscript{\ensuremath{q}}\fi}
\newunicodechar{ₙ}{\ifmmode_{n}\else\textsubscript{\ensuremath{n}}\fi}
\newunicodechar{ₐ}{\ifmmode_{a}\else\textsubscript{\ensuremath{a}}\fi}
\newunicodechar{ᵢ}{\ifmmode_{i}\else\textsubscript{\ensuremath{i}}\fi}
\newunicodechar{ᵣ}{\ifmmode_{r}\else\textsubscript{\ensuremath{r}}\fi}
\newunicodechar{ₖ}{\ifmmode_{k}\else\textsubscript{\ensuremath{k}}\fi}
\newunicodechar{ᵦ}{\ifmmode_{\beta}\else\textsubscript{\ensuremath{\beta}}\fi}
\newunicodechar{ₓ}{\ifmmode_{x}\else\textsubscript{\ensuremath{x}}\fi}
\newfontfamily\popdisplay{ArchivoBlack-Regular.ttf}[Path=../../../fonts/]
\newfontfamily\poplogo{SpaceGrotesk-Bold.ttf}[Path=../../../fonts/]
\newfontfamily\popsemi{PlusJakartaSans-SemiBold.ttf}[Path=../../../fonts/]
\newfontfamily\popxbold{PlusJakartaSans-ExtraBold.ttf}[Path=../../../fonts/]
\newfontfamily\popxboldls{PlusJakartaSans-ExtraBold.ttf}[Path=../../../fonts/,LetterSpace=3.0]

\setlist[enumerate]{leftmargin=1.7em,itemsep=5pt,topsep=4pt}
\setlist[itemize]{leftmargin=1.7em,itemsep=4pt,topsep=4pt,label={\color{poporange}\textbullet}}
\setlength{\parindent}{0pt}
\setlength{\parskip}{5pt}
\renewcommand{\arraystretch}{1.25}

% pgfplots : style Pop par défaut
\pgfplotsset{
  every axis/.append style={axis line style={popink,line width=1pt},tick style={popink},
    grid style={popline,line width=0.5pt},label style={font=\small},tick label style={font=\footnotesize}},
  popcurve/.style={poporange,line width=1.8pt,no markers,smooth},
}
% Même style disponible dans un \draw TikZ simple (hors pgfplots)
\tikzset{popcurve/.style={poporange,line width=1.8pt}}
\tikzset{popgrid/.style={popline,very thin}}
\colorlet{popcurve}{poporange} % le nom du style est parfois utilisé comme couleur

% ============================================================
% Pill / squiggle
% ============================================================
\newcommand{\poppill}[3]{%
  \tikz[baseline=-0.55ex]\node[draw=popink,line width=1.2pt,rounded corners=2.6pt,%
    fill=#2,inner xsep=6.5pt,inner ysep=3pt]{%
    \popxboldls\fontsize{7}{7}\selectfont\color{#3}\MakeUppercase{#1}};%
}
\newcommand{\popsquiggle}[1]{%
  \tikz[baseline=-0.4ex]{
    \draw[#1,line width=2.6pt,line cap=round]
      plot [smooth] coordinates {(0,0.09) (0.34,0.01) (0.68,0.21) (1.02,0.01) (1.36,0.10)};
  }%
}
% Mot-clé surligné (marqueur orange sous le texte)
\newcommand{\cle}[1]{{\popxbold\color{popdark}#1}}

% ============================================================
% En-tête / pied
% ============================================================
\pagestyle{fancy}
\fancyhf{}
\renewcommand{\headrulewidth}{0pt}
\renewcommand{\footrulewidth}{0pt}
\lhead{\raisebox{-0.15ex}{\poplogo\fontsize{13}{13}\selectfont\color{popink}OnBuch\color{poporange}+}}
\rhead{\poppill{\DOCMATIERE\ \textbullet\ \DOCNIVEAU\ \textbullet\ Leçon \DOCLECON}{white}{popink}}
\lfoot{\popxbold\fontsize{7.6}{7.6}\selectfont\color{popmuted}\MakeUppercase{Cours signé OnBuch+}\ \textbullet\ {\popsemi\DOCTITRE}}
\rfoot{\tikz[baseline=-0.6ex]\node[rounded corners=8.5pt,fill=popink,minimum height=17pt,minimum width=17pt,inner xsep=5pt,inner sep=0]{\popxbold\fontsize{8}{8}\selectfont\color{popcream}\thepage\,/\,\pageref*{LastPage}};}

% ============================================================
% Titres
% ============================================================
\newcommand{\chaptitre}[2]{%
  \begin{tcolorbox}[enhanced,colback=poporange,colframe=popink,boxrule=2.6pt,arc=11pt,
      drop shadow=popink,left=15pt,right=15pt,top=12pt,bottom=11pt,before skip=3pt,after skip=18pt]
    \poppill{#2}{popink}{popcream}\par\vspace{9pt}
    {\raggedright\hyphenpenalty=10000\exhyphenpenalty=10000\popdisplay\fontsize{22}{24}\selectfont\color{popcream}\MakeUppercase{#1}\par}
  \end{tcolorbox}
}
% Section numérotée : pastille orange avec le numéro + titre Archivo
\newcounter{popsec}
\newcommand{\coursec}[1]{%
  \stepcounter{popsec}\setcounter{popsub}{0}%
  \par\vspace{16pt}\noindent
  \begin{minipage}{\linewidth}
  \tikz[baseline=-0.7ex]\node[draw=popink,line width=1.6pt,fill=poporange,rounded corners=4pt,minimum size=22pt,inner sep=0]{\popdisplay\fontsize{12}{12}\selectfont\color{popcream}\thepopsec};%
  \hspace{8pt}{\popdisplay\fontsize{15}{17}\selectfont\color{popink}\MakeUppercase{#1}}\par
  \vspace{4pt}\noindent{\color{poporange}\rule{36pt}{3.6pt}}
  \end{minipage}\par\nopagebreak\vspace{8pt}
}
\newcounter{popsub}
\newcommand{\courssub}[1]{%
  \stepcounter{popsub}%
  \par\vspace{9pt}\noindent{\popxbold\fontsize{11.5}{13}\selectfont\color{popink}{\color{poporange}\thepopsec.\thepopsub}\hspace{6pt}#1}\par\nopagebreak\vspace{4pt}
}

% ============================================================
% Boîtes pédagogiques — même recette : fond blanc, bordure épaisse colorée,
% ombre dure encre, badge plein. Titre optionnel [#1].
% ============================================================
\tcbset{popbase/.style={breakable,enhanced,colback=white,boxrule=2.3pt,arc=9pt,
  shadow={3pt}{-3pt}{0pt}{popink},left=14pt,right=14pt,top=11pt,bottom=11pt,before skip=11pt,after skip=15pt,
  fontupper=\popsemi\fontsize{10.6}{14.5}\selectfont\color{popink},
  fontlower=\popsemi\fontsize{10.6}{14.5}\selectfont\color{popink}}}
% Coche dessinée (le glyphe ✓ n'existe pas dans les polices du document)
\newcommand{\popcheck}[1]{\tikz[baseline=-0.5ex]\draw[#1,line width=1.6pt,line cap=round,line join=round](-0.13,0.0)--(-0.04,-0.09)--(0.13,0.1);}
\providecommand{\coche}{\popcheck{popgreen}}
\providecommand{\resultat}[1]{\par\smallskip\noindent\fcolorbox{popgreen}{popgreenL}{\parbox{\dimexpr\linewidth-2\fboxsep-2\fboxrule\relax}{\textbf{Résultat.} #1}}\par\smallskip}
\newcommand{\popboxhead}[3]{\poppill{#1}{#2}{white}\ifx\relax#3\relax\else\hspace{6pt}{\popxbold\fontsize{10.6}{12}\selectfont\color{popink}#3}\fi\par\vspace{7pt}}

\newtcolorbox{aretenir}[1][]{popbase,colframe=popgreen,before upper={\popboxhead{À retenir}{popgreen}{#1}}}
\newtcolorbox{exemplebox}[1][]{popbase,colframe=poporange,before upper={\popboxhead{Exemple}{poporange}{#1}}}
\newtcolorbox{definition}[1][]{popbase,colframe=popblue,before upper={\popboxhead{Définition}{popblue}{#1}}}
\newtcolorbox{propriete}[1][]{popbase,colframe=poppurple,before upper={\popboxhead{Propriété}{poppurple}{#1}}}
\newtcolorbox{methode}[1][]{popbase,colframe=popgold,before upper={\popboxhead{Méthode}{popgold}{#1}}}
\newtcolorbox{attention}[1][]{popbase,colframe=poppink,before upper={\popboxhead{Attention}{poppink}{#1}}}
\newtcolorbox{experience}[1][]{popbase,colframe=popblue,colback=popblueL,before upper={\popboxhead{Expérience}{popblue}{#1}}}
% « Le savais-tu ? » : carte violette pleine (contexte, culture, Cameroun)
\newtcolorbox{savaistu}[1][]{popbase,colback=poppurple,colframe=popink,
  fontupper=\popsemi\fontsize{10.4}{14.2}\selectfont\color{white},
  before upper={\poppill{Le savais-tu\,?}{white}{poppurple}\ifx\relax#1\relax\else\hspace{6pt}{\popxbold\color{white}#1}\fi\par\vspace{7pt}}}
% Exercice résolu : énoncé en haut, \tcblower puis la solution sur fond crème
\newcounter{popexo}\newcounter{popexr}
\newtcolorbox{exoresolu}[1][]{popbase,colframe=poporange,colbacklower=popcream,
  segmentation style={popink,line width=1.2pt,dashed},
  before upper={\stepcounter{popexr}\popboxhead{Exercice résolu \thepopexr}{poporange}{#1}},
  before lower={\poppill{Solution}{popink}{popcream}\par\vspace{6pt}}}
% Exercice (non résolu en place) — la correction est dans la partie Corrigés
\newtcolorbox{exercice}[1][]{popbase,colframe=popink,
  before upper={\stepcounter{popexo}\popboxhead{Exercice \thepopexo}{popink}{#1}}}
% Corrigé
\newtcolorbox{corrige}[1][]{popbase,colframe=popgreen,colback=popgreenL,
  before upper={\popboxhead{Corrigé}{popgreen}{#1}}}
% Objectifs / prérequis / bilan
\newtcolorbox{objectifs}{popbase,colframe=popink,colback=poporangeL,
  before upper={\popboxhead{Objectifs de la leçon}{popink}{}}}
\newtcolorbox{prerequis}{popbase,colframe=popink,colback=popblueL,
  before upper={\popboxhead{Prérequis}{popblue}{}}}
\newtcolorbox{bilan}{popbase,colframe=popink,colback=white,boxrule=2.8pt,
  before upper={\popboxhead{Fiche bilan}{poporange}{L'essentiel à connaître pour l'examen}}}
% Figure encadrée avec légende
\newtcolorbox{popfigure}[1][]{enhanced,breakable=false,colback=white,colframe=popline,boxrule=1.4pt,arc=8pt,
  left=10pt,right=10pt,top=10pt,bottom=8pt,before skip=10pt,after skip=12pt,halign=center,
  lower separated=false,halign lower=center,
  fontlower=\popsemi\fontsize{9}{11.5}\selectfont\color{popmuted},
  #1}
\newcommand{\legende}[1]{\par\vspace{6pt}{\popsemi\fontsize{9}{11.5}\selectfont\color{popmuted}#1}}

% ============================================================
% Signature OnBuch+
% ============================================================
\newcommand{\poplogoinline}[1]{{\poplogo\fontsize{#1}{#1}\selectfont\color{popink}OnBuch\color{poporange}+}}
\newcommand{\signatureonbuch}{%
  \begin{tcolorbox}[enhanced,colback=popink,colframe=popink,boxrule=2.2pt,arc=10pt,
      drop shadow=poporange,left=14pt,right=14pt,top=10pt,bottom=10pt,before skip=0pt,after skip=0pt]
    \tikz[baseline=-0.7ex]\node[circle,fill=popgreen,draw=popcream,line width=1.3pt,minimum size=22pt,inner sep=0]{\popcheck{white}};%
    \hspace{9pt}\begin{minipage}[c]{0.62\linewidth}
      {\popxbold\fontsize{12}{13}\selectfont\color{popcream}Signé\ \textbullet\ {\poplogo OnBuch\color{poporange}+}}\par\vspace{2pt}
      {\raggedright\popsemi\fontsize{8}{10.5}\selectfont\color{popline}\MakeUppercase{Équipe pédagogique OnBuch+}\\\MakeUppercase{Conforme au programme officiel MINESEC}\par}
    \end{minipage}\hfill
    {\poplogo\fontsize{22}{22}\selectfont\color{popcream}OnBuch\color{poporange}+}
  \end{tcolorbox}%
}

% ============================================================
% Page de garde
% #1 titre · #2 matière · #3 niveau · #4 module · #5 n° leçon · #6 durée ·
% #7 description / objectifs (texte ou itemize)
% ============================================================
\newcommand{\pagegarde}[7]{%
  \thispagestyle{empty}
  \newgeometry{a4paper,margin=1.7cm,top=1.6cm,bottom=1.4cm}%
  \begin{tikzpicture}[remember picture,overlay]
    \fill[poporange!18] ([xshift=-3.2cm,yshift=-3.6cm]current page.north east) circle (4.6cm);
    \fill[poppurple!14] ([xshift=2.4cm,yshift=7.6cm]current page.south west) circle (3.8cm);
  \end{tikzpicture}%
  {\poplogo\fontsize{30}{30}\selectfont\color{popink}OnBuch\color{poporange}+}\hfill
  \raisebox{0.6ex}{\poppill{Cours signé \textbullet\ Premium}{popink}{popcream}}\par
  \vspace{1pt}\popsquiggle{poppurple}\par
  \vspace{8pt}
  \begin{tcolorbox}[enhanced,colback=poporange,colframe=popink,boxrule=2.8pt,arc=13pt,
      drop shadow=popink,left=18pt,right=18pt,top=12pt,bottom=13pt,after skip=10pt]
    \poppill{#2 \textbullet\ #3}{popink}{popcream}\hspace{5pt}\poppill{#4}{white}{popink}\par\vspace{8pt}
    {\popdisplay\fontsize{14}{14}\selectfont\color{popink}LEÇON #5}\par\vspace{4pt}
    {\raggedright\hyphenpenalty=10000\exhyphenpenalty=10000\popdisplay\fontsize{25}{27}\selectfont\color{popcream}\MakeUppercase{#1}\par}
    \vspace{8pt}
    {\popxbold\fontsize{9.5}{9.5}\selectfont\color{popink}\MakeUppercase{Durée conseillée : #6}}
  \end{tcolorbox}
  \begin{tcolorbox}[enhanced,colback=white,colframe=popink,boxrule=2.2pt,arc=10pt,
      drop shadow=popink,left=16pt,right=16pt,top=10pt,bottom=10pt,after skip=8pt]
    \poppill{Dans cette leçon}{poppurple}{white}\par\vspace{5pt}
    \popsemi\fontsize{9.3}{12.6}\selectfont\color{popink}#7
  \end{tcolorbox}
  \noindent\begin{minipage}[t]{0.487\linewidth}
    \begin{tcolorbox}[enhanced,colback=popgreen,colframe=popink,boxrule=2.2pt,arc=10pt,
        drop shadow=popink,left=11pt,right=11pt,top=10pt,bottom=10pt,height=3.1cm,valign=center]
      \tcbox[colback=white,colframe=popink,boxrule=1.3pt,arc=5pt,boxsep=0pt,nobeforeafter,
        left=3pt,right=3pt,top=3pt,bottom=3pt,valign=center]{\qrcode[height=1.9cm]{https://play.google.com/store/apps/details?id=com.luvvix.onbuch&hl=fr}}%
      \hspace{8pt}\begin{minipage}[c]{0.5\linewidth}
        {\popdisplay\fontsize{11}{12.5}\selectfont\color{white}\MakeUppercase{Télécharge OnBuch}}\par\vspace{4pt}
        {\raggedright\popsemi\fontsize{8.4}{11}\selectfont\color{white}Gratuit \textbullet\ Google Play\\Tuteur IA, annales, concours, résultats\par}
      \end{minipage}
    \end{tcolorbox}
  \end{minipage}\hfill
  \begin{minipage}[t]{0.487\linewidth}
    \begin{tcolorbox}[enhanced,colback=poppurple,colframe=popink,boxrule=2.2pt,arc=10pt,
        drop shadow=popink,left=11pt,right=11pt,top=10pt,bottom=10pt,height=3.1cm,valign=center]
      \tikz[baseline=-0.7ex]\node[circle,fill=white,draw=popink,line width=1.3pt,minimum size=1.6cm,inner sep=0]{\popdisplay\fontsize{24}{24}\selectfont\color{poppurple}+};%
      \hspace{8pt}\begin{minipage}[c]{0.6\linewidth}
        {\popdisplay\fontsize{11}{12.5}\selectfont\color{white}\MakeUppercase{Passe à OnBuch+}}\par\vspace{4pt}
        {\raggedright\popsemi\fontsize{8.4}{11}\selectfont\color{white}Tous les cours signés, Tuteur IA illimité, fiches PDF — onglet OnBuch+ de l'app\par}
      \end{minipage}
    \end{tcolorbox}
  \end{minipage}
  \vspace{8pt}
  \popcontenu
  \vfill
  \signatureonbuch
  \restoregeometry
  \newpage
}

% Rangée « Ce cours contient » (page de garde)
\newcommand{\poptile}[3]{% #1 couleur #2 gros texte #3 légende
  \begin{tcolorbox}[enhanced,colback=#1,colframe=popink,boxrule=2pt,arc=9pt,shadow={2.5pt}{-2.5pt}{0pt}{popink},
      left=6pt,right=6pt,top=9pt,bottom=9pt,halign=center,valign=center,height=1.75cm,width=\linewidth,nobeforeafter]
    {\popdisplay\fontsize{12.5}{13}\selectfont\color{white}\MakeUppercase{#2}}\par\vspace{3pt}
    {\centering\popsemi\fontsize{7.8}{9.5}\selectfont\color{white}#3\par}
  \end{tcolorbox}}
\newcommand{\popcontenu}{%
  \par\noindent{\popxboldls\fontsize{7.5}{7.5}\selectfont\color{popmuted}CE COURS SIGNÉ CONTIENT}\par\vspace{7pt}
  \noindent\begin{minipage}[t]{0.235\linewidth}\poptile{popblue}{Cours}{complet, illustré, pas à pas}\end{minipage}\hfill
  \begin{minipage}[t]{0.235\linewidth}\poptile{poporange}{Méthodes}{et exercices résolus}\end{minipage}\hfill
  \begin{minipage}[t]{0.235\linewidth}\poptile{poppink}{Exercices}{type examen + corrigés}\end{minipage}\hfill
  \begin{minipage}[t]{0.235\linewidth}\poptile{popgreen}{Fiche bilan}{l'essentiel à retenir}\end{minipage}\par}

% Sommaire
\newtcolorbox{sommaire}{enhanced,colback=white,colframe=popink,boxrule=2.2pt,arc=10pt,
  drop shadow=popink,left=18pt,right=18pt,top=13pt,bottom=13pt,before skip=6pt,after skip=20pt,
  before upper={\poppill{Sommaire}{poppurple}{white}\par\vspace{9pt}\popsemi\fontsize{10.6}{14.5}\selectfont\color{popink}}}

% Signature de fin de cours — poussée tout en bas de la dernière page
\newcommand{\findecours}{%
  \par\vfill\nopagebreak
  \begin{center}\popsquiggle{poporange}\end{center}\vspace{4pt}
  \signatureonbuch
}
\AtBeginDocument{\def\llbracket{\mathopen{[\mkern-3.5mu[}}\def\rrbracket{\mathclose{]\mkern-3.5mu]}}} % intervalles d'entiers (glyphe ⟦ absent de la police)
% Glyphes absents de Fira Math : redéfinis à partir de glyphes disponibles
\AtBeginDocument{%
  \def\setminus{\mathbin{\backslash}}\let\smallsetminus\setminus
  \def\vdots{\mathord{\vbox{\baselineskip3.2pt\lineskiplimit0pt\kern4pt\hbox{.}\hbox{.}\hbox{.}}}}%
  \def\ddots{\mathinner{\mkern1mu\raise6.5pt\hbox{.}\mkern2mu\raise3.6pt\hbox{.}\mkern2mu\raise0.7pt\hbox{.}\mkern1mu}}%
  \def\triangle{\mathord{\scalebox{0.9}{$\symup{\Delta}$}}}%
  \def\nabla{\mathord{\raisebox{\depth}{\scalebox{1}[-1]{$\symup{\Delta}$}}}}%
  \def\checkmark{\popcheck{popdark}}%
  \def\star{\mathbin{\ast}}\def\bigstar{\mathord{\ast}}%
  \def\diamond{\mathbin{\scalebox{0.6}{$\Diamond$}}}%
  \def\bigcup{\mathop{\mathchoice{\raisebox{-0.3em}{\scalebox{1.7}{$\cup$}}}{\scalebox{1.25}{$\cup$}}{\cup}{\cup}}\displaylimits}%
  \def\bigcap{\mathop{\mathchoice{\raisebox{-0.3em}{\scalebox{1.7}{$\cap$}}}{\scalebox{1.25}{$\cap$}}{\cap}{\cap}}\displaylimits}%
  \def\lesseqqgtr{\mathrel{\lessgtr}}\def\bigoplus{\mathop{\oplus}\displaylimits}%
}
\providecommand{\ssi}{si et seulement si} % abréviation fréquente
\AtBeginDocument{\providecommand{\wideparen}{\overparen}} % arc de cercle

%% FIGFIX-BEGIN v3 — correctifs globaux des figures (docs/onbuch-generation/tools/figfix)
\usepackage{adjustbox}\usepackage{etoolbox}
% toute figure TikZ/pgfplots plus large que la ligne est réduite à la largeur disponible
\BeforeBeginEnvironment{tikzpicture}{\begin{adjustbox}{max width=\linewidth}}
\AfterEndEnvironment{tikzpicture}{\end{adjustbox}}
% légendes pgfplots lisibles par-dessus les courbes
\pgfplotsset{every axis/.append style={legend style={fill=white,fill opacity=0.92,text opacity=1,draw=black!15,inner sep=3pt}}}
% étiquettes d'arêtes (syntaxe edge["..."]) avec fond blanc
\tikzset{every edge quotes/.append style={fill=white,inner sep=1.2pt}}
% bulles de cartes mentales : texte à la ligne dans la bulle, bulle un peu plus grande
\tikzset{every concept/.append style={text width=2.0cm,align=center,minimum size=2.3cm,inner sep=2pt}}
%% FIGFIX-END
\begin{document}

\pagegarde{L’Etat et le pouvoir}{Philosophie}{1re Tech}{Philosophie – classe de Première technique}{N10}{3 séances (fiche de progression harmonisée)}{Cette leçon interroge la nature de l'État : ce qui le définit, ce qui fonde son autorité et la légitimité de son pouvoir. Elle examine ensuite les rôles concrets de l'État dans la vie des Camerounais et sa finalité : l'État est-il une fin en soi ou un instrument au service de l'homme et du bien commun ?
\vspace{4pt}
\begin{multicols}{2}\small
\begin{itemize}
\item À la fin de cette leçon, je sais définir l'État et identifier ses éléments constitutifs (territoire, population, pouvoir souverain) à partir d'exemples camerounais.
\item Je sais distinguer l'État des autres formes de pouvoir (famille, chefferie traditionnelle, église, entreprise).
\item Je sais expliquer les différents fondements du pouvoir étatique (force, tradition, contrat social, droit) en mobilisant les thèses de Hobbes, Locke, Rousseau et Weber.
\item Je sais analyser un document (texte de philosophe, article de la Constitution camerounaise, donnée chiffrée) pour dégager une conception de l'État.
\item Je sais décrire les rôles de l'État (régalien, économique, social) et les illustrer par des situations concrètes du Cameroun.
\item Je sais discuter la finalité de l'État : instrument du bien commun ou fin en soi, en confrontant plusieurs positions philosophiques.
\end{itemize}\end{multicols}}

\begin{sommaire}
\begin{multicols}{2}
\begin{enumerate}
\item Qu'est-ce que l'État ? Définition et éléments constitutifs
\item Les fondements de l'État : pourquoi obéit-on au pouvoir ?
\item Les rôles de l'État : que fait concrètement l'État ?
\item La finalité de l'État : pour quoi l'État existe-t-il ?
\item TP guidé : construire des outils d'analyse de l'État camerounais
\item Méthode et entraînement : rédiger et justifier une réponse philosophique
\item Activité d'intégration
\item Méthodes et exercices résolus
\item Exercices
\item Corrigés des exercices
\item Fiche bilan
\end{enumerate}
\end{multicols}
\end{sommaire}

\begin{objectifs}
\begin{itemize}
\item À la fin de cette leçon, je sais définir l'État et identifier ses éléments constitutifs (territoire, population, pouvoir souverain) à partir d'exemples camerounais.
\item Je sais distinguer l'État des autres formes de pouvoir (famille, chefferie traditionnelle, église, entreprise).
\item Je sais expliquer les différents fondements du pouvoir étatique (force, tradition, contrat social, droit) en mobilisant les thèses de Hobbes, Locke, Rousseau et Weber.
\item Je sais analyser un document (texte de philosophe, article de la Constitution camerounaise, donnée chiffrée) pour dégager une conception de l'État.
\item Je sais décrire les rôles de l'État (régalien, économique, social) et les illustrer par des situations concrètes du Cameroun.
\item Je sais discuter la finalité de l'État : instrument du bien commun ou fin en soi, en confrontant plusieurs positions philosophiques.
\item Je sais construire un schéma ou un tableau comparatif synthétisant les conceptions de l'État.
\item Je sais rédiger et justifier une réponse argumentée (dissertation ou explication de texte) sur un problème relatif à l'État et au pouvoir.
\end{itemize}
\end{objectifs}

\begin{prerequis}
\begin{itemize}
\item Notion de société et de vivre-ensemble (cours sur l'homme et la société, début d'année).
\item Notions de liberté et d'obligation (leçon sur le devoir).
\item Connaissance générale des institutions camerounaises (Président, Assemblée, tribunaux) vue en Éducation à la citoyenneté au collège.
\item Distinction entre règle morale, règle religieuse et règle juridique.
\item Notion d'autorité et d'obéissance dans la vie quotidienne (famille, école, village).
\end{itemize}
\end{prerequis}

\newpage

\coursec{Qu'est-ce que l'État ? Définition et éléments constitutifs}

\textbf{Situation d'accroche.} En janvier 2024, lors d'une panne prolongée d'électricité à Douala, les habitants d'un quartier se sont organisés eux-mêmes pour sécuriser les rues et répartir l'eau d'un forage. Très vite, des conflits ont éclaté : qui décide ? qui sanctionne ? qui a le droit d'imposer une règle à tous ? Certains jeunes voulaient punir eux-mêmes un voleur présumé ; le chef de quartier réclamait l'autorité ; un notable invoquait la tradition ; finalement, c'est l'arrivée de la police, représentante de l'État, qui a rétabli l'ordre. Cette situation montre que dès que l'autorité publique fait défaut, la question du pouvoir devient vitale. Elle pose notre problématique : qu'est-ce qui distingue le chef de quartier, le chef traditionnel et l'État camerounais ? Pourquoi obéit-on à l'État, et jusqu'où doit aller son pouvoir ? Pour répondre, il faut d'abord savoir ce qu'est l'État : c'est l'objet de cette première section.

\courssub{Le mot « État » : du sens commun au sens philosophique}

Dans le langage courant, le mot « État » a plusieurs sens. On parle de « l'état de santé » d'un malade, de « l'état des routes » après la saison des pluies, de « faire état de » quelque chose. Ici, « état » signifie simplement la manière d'être d'une chose à un moment donné. Mais en philosophie politique, le mot prend un sens précis et technique : il désigne une \cle{institution}, c'est-à-dire une organisation durable qui structure la vie collective.

\begin{savaistu}[D'où vient le mot « État » ?]
Le mot « État » vient du latin \emph{status}, qui signifie « manière d'être », « position stable ». C'est le penseur italien Nicolas Machiavel qui l'impose au sens politique moderne dans son ouvrage \emph{Le Prince} (1513) : « lo Stato » désigne alors l'ensemble du pouvoir politique organisé sur un territoire. Le mot exprime donc l'idée de \cle{stabilité} : l'État est ce qui demeure, par-delà les personnes qui gouvernent.
\end{savaistu}

Pour bien penser l'État, il faut le distinguer de deux notions voisines avec lesquelles on le confond souvent : le \cle{gouvernement} et la \cle{nation}.

\begin{itemize}
\item Le \textbf{gouvernement} est l'\emph{équipe dirigeante} qui administre l'État à un moment donné : le Président, le Premier ministre, les ministres. Au Cameroun, les remaniements ministériels changent les visages du gouvernement, mais l'État, lui, ne change pas : les impôts continuent d'être collectés, les écoles continuent de fonctionner, les tribunaux continuent de juger.
\item La \textbf{nation} est une \emph{communauté culturelle et humaine} : un peuple uni par une histoire, une langue, des traditions, une volonté de vivre ensemble. Le Cameroun rassemble plus de 250 groupes ethniques qui forment une seule nation, mais une nation peut exister sans État propre (comme les Kurdes), et un État peut rassembler plusieurs nations culturelles.
\item L'\textbf{État} est le \emph{cadre juridique et politique permanent} : il possède des institutions (Constitution, administration, armée, justice) qui survivent aux hommes.
\end{itemize}

\begin{attention}[Erreur fréquente : confondre État et gouvernement]
Dans une dissertation, ne dites jamais « l'État a été remanié » ou « l'État a perdu les élections ». Ce sont les \textbf{gouvernements} qui changent (remaniements, élections, alternances) ; l'\textbf{État demeure}. Quand un nouveau Président entre au Palais de l'Unité à Yaoundé, l'État camerounais continue d'exister : ses lois, ses frontières et ses administrations restent les mêmes. Retenez la formule : \emph{les gouvernements passent, l'État reste}.
\end{attention}

\courssub{Définition philosophique de l'État}

Partons d'une intuition simple. Imaginez un groupe de personnes installées quelque part, sans aucune règle commune ni aucune autorité reconnue : c'est la situation du quartier de Douala pendant la panne. Chacun fait ce qu'il veut, les conflits se règlent par la force du plus fort. Pour que la vie collective soit possible, il faut trois choses : un \emph{lieu} où vivre, des \emph{personnes} qui y vivent, et une \emph{autorité} capable d'imposer des règles à tous. Ces trois éléments réunis, c'est exactement l'État.

\begin{definition}[L'État]
L'\cle{État} est une organisation politique \cle{souveraine} qui exerce un pouvoir de contrainte légitime sur une \cle{population} installée sur un \cle{territoire} déterminé.
\end{definition}

Cette définition contient les trois éléments constitutifs de l'État, qu'aucun État au monde ne peut perdre sans cesser d'exister :

\begin{enumerate}
\item \textbf{Un territoire délimité} : un espace géographique précis, borné par des frontières reconnues. Le territoire comprend le sol, le sous-sol (minerais, pétrole), les eaux intérieures (fleuves, lacs) et l'espace aérien.
\item \textbf{Une population} : l'ensemble des personnes qui vivent sur ce territoire et qui sont soumises aux lois de l'État, qu'elles soient citoyennes ou étrangères résidentes.
\item \textbf{Un pouvoir souverain} : une autorité politique organisée qui s'exerce sur la population et le territoire, capable de faire des lois, de les appliquer et de sanctionner ceux qui les violent.
\end{enumerate}

\begin{popfigure}
\begin{center}
\begin{tikzpicture}[scale=1]
% Cercles concentriques : Territoire (ext) > Population (mid) > Pouvoir (centre)
\fill[popblueL] (0,0) circle (3.4);
\fill[popgreenL] (0,0) circle (2.3);
\fill[poporangeL] (0,0) circle (1.2);
% Labels
\node[popdark,font=\bfseries, align=center] at (0,0){\shortstack{Pouvoir\\souverain}};
\node[popdark,font=\bfseries] at (0,1.75) {Population};
\node[popdark,font=\bfseries] at (0,2.95) {Territoire};
% Exemples camerounais (annotations)
\draw[->,popblue,thick] (3.4,2.6) -- (5.0,2.6) node[right,align=left,font=\small,text=popdark] {\SI{475000}{km^2}, frontières\\avec Nigeria, Tchad, RCA,\\Gabon, Guinée-Équatoriale, Congo};
\draw[->,popgreen!60!black,thick] (2.3,-1.4) -- (5.0,-1.4) node[right,align=left,font=\small,text=popdark] {Plus de \num{28000000}\\d'habitants (est. 2024)};
\draw[->,poporange!80!black,thick] (1.2,-0.6) -- (5.0,-3.2) node[right,align=left,font=\small,text=popdark] {Constitution 18 fév. 1996,\\administration, armée, justice};
\end{tikzpicture}
\end{center}
\legende{Les trois éléments constitutifs de l'État : le pouvoir souverain (centre) s'exerce sur la population, elle-même installée sur un territoire délimité. Données camerounaises indicatives.}
\end{popfigure}

\begin{exemplebox}[L'État camerounais en chiffres]
Le Cameroun compte environ \textbf{\num{28000000} d'habitants}, répartis dans \textbf{10 régions}, \textbf{58 départements} et \textbf{360 communes}. Chacun de ces échelons est un niveau où s'exerce concrètement le pouvoir étatique : le gouverneur de région, le préfet de département, le sous-préfet d'arrondissement et le maire de commune sont les représentants ou les relais de l'État auprès des citoyens. Ainsi, un élève de Ngaoundéré comme un élève de Bafoussam relève du même État, des mêmes lois et de la même Constitution.
\end{exemplebox}

\begin{popfigure}
\begin{center}
\begin{tikzpicture}[scale=0.9]
% Contour très simplifié du Cameroun (polygone approximatif)
\draw[thick,popdark,fill=popcream] 
  (0,0) -- (1.6,-0.3) -- (3.2,0.1) -- (4.4,0.8) -- (4.8,1.8) -- (4.2,2.8) 
  -- (4.6,3.8) -- (3.8,4.6) -- (3.4,5.6) -- (2.6,6.4) -- (1.8,7.4) -- (1.2,8.2) 
  -- (0.6,7.2) -- (0.2,6.0) -- (-0.4,5.0) -- (-0.2,3.8) -- (-0.8,2.8) -- (-0.6,1.4) -- cycle;
% Capitale
\fill[popink] (1.6,1.6) circle (0.09);
\node[right,font=\small\bfseries,text=popink] at (1.75,1.6) {Yaoundé (capitale)};
% Quelques villes repères
\fill[popdark] (0.4,0.7) circle (0.07); \node[left,font=\footnotesize] at (0.3,0.7) {Douala};
\fill[popdark] (2.8,4.6) circle (0.07); \node[right,font=\footnotesize] at (2.9,4.6) {Ngaoundéré};
\fill[popdark] (1.6,7.0) circle (0.07); \node[right,font=\footnotesize] at (1.7,7.0) {Maroua};
\fill[popdark] (0.9,2.6) circle (0.07); \node[left,font=\footnotesize] at (0.8,2.6) {Bafoussam};
% Flèche Nord
\draw[->,thick,popblue] (-1.2,6.2) -- (-1.2,7.2);
\node[font=\small\bfseries,text=popblue] at (-1.2,7.5) {N};
% Légende intégrée
\node[align=center,font=\small,text=popdark,fill=popblueL,rounded corners,inner sep=4pt] at (6.8,5.5) {\textbf{10 régions} : Adamaoua, Centre,\\Est, Extrême-Nord, Littoral, Nord,\\Nord-Ouest, Ouest, Sud, Sud-Ouest};
\node[align=left,font=\footnotesize,text=popmuted] at (6.8,3.6) {Carte schématique simplifiée,\\sans précision cartographique.};
\end{tikzpicture}
\end{center}
\legende{Carte stylisée du Cameroun : le territoire national, délimité par des frontières héritées de la colonisation, s'étend sur environ \SI{475000}{km^2} et comprend 10 régions. Yaoundé est la capitale politique, siège des institutions de la République.}
\end{popfigure}

\courssub{La souveraineté : le cœur du pouvoir étatique}

Le troisième élément, le pouvoir souverain, mérite une attention particulière, car c'est lui qui distingue l'État de toute autre association. La \cle{souveraineté} est le caractère suprême du pouvoir de l'État. Elle a deux visages :

\begin{itemize}
\item \textbf{À l'intérieur}, la souveraineté signifie que l'État est la \emph{puissance suprême} : aucune autre autorité sur le territoire ne peut lui contester le dernier mot. L'État décide des lois, lève les impôts, organise la justice, et peut recourir à la force pour faire respecter ses décisions. C'est ce qu'on appelle le \cle{monopole de la contrainte légitime}.
\item \textbf{À l'extérieur}, la souveraineté signifie l'\emph{indépendance} : l'État camerounais ne reçoit d'ordres d'aucun autre État. Il signe lui-même ses traités, entretient des relations diplomatiques, est membre de l'Organisation des Nations unies et de l'Union africaine. Depuis les indépendances de 1960 et la réunification de 1961, le Cameroun est un État souverain reconnu par la communauté internationale.
\end{itemize}

\begin{definition}[Le monopole de la violence légitime (Max Weber)]
Le sociologue allemand Max Weber (1864-1920) donne de l'État la définition suivante : « L'État est une communauté humaine qui revendique, avec succès, le \cle{monopole de la violence physique légitime} sur un territoire donné. » Cela signifie que seul l'État a le droit d'user légalement de la force : la police peut arrêter un suspect, l'armée peut défendre le territoire, la justice peut condamner un coupable à la prison. Aucun particulier, aucun groupe privé ne dispose de ce droit.
\end{definition}

Reprenons notre situation d'accroche : à Douala, les jeunes du quartier qui voulaient punir eux-mêmes le voleur présumé commettaient une \emph{justice privée}, ce que le droit interdit formellement. Ni le chef de quartier, ni un groupe d'autodéfense, ni une entreprise de sécurité privée ne peuvent emprisonner ou infliger une sanction physique : seuls les organes de l'État (police, gendarmerie, tribunaux) le peuvent, et selon des procédures fixées par la loi. Si chacun se faisait justice, ce serait la loi du plus fort, c'est-à-dire la guerre de tous contre tous que décrivait le philosophe anglais Thomas Hobbes dans \emph{Le Léviathan} (1651).

\courssub{L'État face aux autres formes de pouvoir}

Dans la vie quotidienne camerounaise, plusieurs autorités coexistent avec l'État. Il est essentiel de les distinguer pour ne pas leur attribuer les prérogatives de l'État.

\begin{center}
\begin{tabularx}{\linewidth}{|l|X|X|}
\hline
\rowcolor{popblueL} \textbf{Type de pouvoir} & \textbf{Exemple camerounais} & \textbf{Limite face à l'État} \\
\hline
Pouvoir paternel / familial & Un père éduque et corrige ses enfants & Ne peut violer les lois (interdiction des châtiments graves, scolarisation obligatoire) \\
\hline
Autorité traditionnelle & Lamibé du Nord, fons de l'Ouest, sultans & Reconnue par l'État comme auxiliaire d'administration, mais ne peut ni juger les crimes ni lever d'impôt \\
\hline
Autorité religieuse & Pasteurs, imams, évêques & S'exerce sur les consciences, pas sur les corps : aucune contrainte légale \\
\hline
Pouvoir économique & Grandes entreprises, banques & Dispose de l'argent, non de la force légitime ; doit payer les impôts et respecter le code du travail \\
\hline
\end{tabularx}
\end{center}

\begin{aretenir}[L'essentiel]
L'État se définit par trois éléments indissociables : un \cle{territoire}, une \cle{population}, un \cle{pouvoir souverain}. Sa marque propre est le monopole de la violence physique légitime (Weber) : seul l'État peut légalement contraindre. Toutes les autres autorités (familiale, traditionnelle, religieuse, économique) s'exercent \emph{à l'intérieur du cadre} que l'État fixe et sous le contrôle de ses lois.
\end{aretenir}

\courssub{L'État de droit : l'État soumis à ses propres lois}

Une dernière précision est indispensable. Dire que l'État détient le monopole de la force ne signifie pas qu'il peut tout faire. Un État qui use de la force sans règles, selon le bon plaisir de ses dirigeants, n'est pas un État de droit : c'est un \cle{État de fait}, où règne l'\emph{arbitraire} — les arrestations sans motif, les décisions imprévisibles, la loi du plus puissant.

\begin{definition}[L'État de droit]
Un \cle{État de droit} est un État dont le pouvoir lui-même est soumis au \cle{droit} : la Constitution est la loi suprême, les lois sont votées selon des procédures définies, les citoyens peuvent contester les actes de l'administration devant des juges indépendants. Au Cameroun, la Constitution du 18 février 1996 (révisée en 2008) affirme dans son préambule l'attachement du peuple camerounais à l'État de droit et aux droits de l'homme.
\end{definition}

La différence est fondamentale : dans l'État de fait, la force crée le droit ; dans l'État de droit, c'est le droit qui encadre la force. Le policier qui arrête un suspect doit respecter une procédure ; le préfet qui prend un arrêté doit le fonder sur une loi ; le citoyen lésé peut saisir un tribunal. Ainsi, la puissance de l'État est à la fois \emph{forte} (elle peut contraindre) et \emph{limitée} (elle ne peut contraindre que selon la loi).

\begin{exoresolu}[Distinguer État, gouvernement et nation]
\textbf{Énoncé.} Un élève déclare : « Après les élections, si le parti au pouvoir perd, l'État camerounais disparaîtra et la nation devra en construire un autre. » Relevez les deux confusions contenues dans cette phrase et corrigez-les en vous appuyant sur la définition de l'État.
\tcblower
\textbf{Solution détaillée.}
Première confusion : l'élève confond l'\textbf{État} et le \textbf{gouvernement}. Une élection fait changer l'équipe dirigeante (le gouvernement), mais l'État demeure : le territoire de \SI{475000}{km^2}, la population de plus de \num{28000000} d'habitants et le pouvoir souverain organisé par la Constitution de 1996 subsistent, quel que soit le vainqueur du scrutin. Les gouvernements passent, l'État reste.

Deuxième confusion : l'élève confond l'\textbf{État} et la \textbf{nation}. La nation camerounaise, communauté humaine unie par une histoire et une volonté de vivre ensemble, n'a pas besoin de « construire » un État : elle en possède déjà un, doté d'institutions permanentes (administration, justice, armée). La nation est le fondement humain, l'État est le cadre juridique ; ni l'un ni l'autre ne disparaissent à chaque alternance politique.

\textbf{Conclusion corrigée :} « Après les élections, si le parti au pouvoir perd, le \emph{gouvernement} change, mais l'État camerounais et la nation camerounaise demeurent. »
\end{exoresolu}

En résumé, l'État n'est ni une équipe au pouvoir ni un simple peuple : c'est l'organisation politique souveraine qui, sur un territoire délimité, exerce sur une population un pouvoir de contrainte légitime, encadré par le droit. Mais une question demeure : d'où vient ce pouvoir et pourquoi les hommes lui obéissent-ils ? C'est le problème des \emph{fondements} de l'État, que nous examinerons dans la section suivante.

\coursec{Les fondements de l'État : pourquoi obéit-on au pouvoir ?}

Nous avons vu dans la section précédente que l'État se définit par trois éléments constitutifs : un \cle{territoire}, une \cle{population} et un \cle{pouvoir} organisé qui exerce la souveraineté. Mais cette définition, pour nécessaire qu'elle soit, laisse une question entière : d'où vient ce pouvoir ? Pourquoi des millions d'individus acceptent-ils d'obéir à quelques-uns ? Au Cameroun, un automobiliste s'arrête au feu rouge même lorsqu'aucun gendarme n'est visible ; un contribuable paie ses impôts ; un écolier respecte le règlement de son lycée. Qu'est-ce qui rend cette obéissance possible et durable ? C'est le problème des \cle{fondements} de l'État : sur quoi repose l'autorité politique ?

\courssub{Le problème : la force seule suffit-elle ?}

Commençons par une intuition simple. On pourrait croire que l'État tient par la force : il dispose de la police, de l'armée, des tribunaux, des prisons. L'État détient ce que Max Weber appelle le \emph{monopole de la violence physique légitime} : lui seul peut user légalement de la contrainte. Pourtant, un raisonnement rapide montre que la force seule ne peut pas suffire. Les gendarmes et policiers camerounais sont quelques dizaines de milliers face à plus de vingt-cinq millions d'habitants. Si chacun refusait d'obéir, aucune force ne pourrait contraindre tout le monde. Il faut donc que l'obéissance repose sur autre chose que la peur : sur une forme d'\emph{acceptation}.

C'est exactement la thèse d'Étienne de \textbf{La Boétie} dans son \emph{Discours de la servitude volontaire} (1549). Il s'étonne que « un seul homme » puisse commander à des multitudes, et répond : ce n'est pas la force du tyran qui l'explique, mais le \emph{consentement} des sujets. On n'obéit pas seulement par peur ; on obéit aussi par habitude, par intérêt, par croyance. Le tyran n'a de puissance que celle qu'on lui donne : « Soyez résolus de ne servir plus, et vous voilà libres. » La Boétie renverse donc le problème : le mystère n'est pas la désobéissance, mais l'\cle{obéissance volontaire}.

\begin{definition}[La légitimité]
La \cle{légitimité} est le caractère d'un pouvoir reconnu comme juste et acceptable par ceux qui lui obéissent. Elle fonde l'obéissance durable : on se soumet au pouvoir non seulement parce qu'il peut contraindre, mais parce qu'on juge qu'il a le \emph{droit} de commander. Un pouvoir illégitime n'est que domination ; un pouvoir légitime est autorité.
\end{definition}

Toute la réflexion philosophique sur les fondements de l'État peut alors s'organiser comme une enquête sur les sources possibles de cette légitimité. Nous en distinguerons quatre : la force, la tradition et le sacré, le contrat social, et enfin le droit rationnel.

\courssub{Premier fondement : la force et la contrainte}

L'État est d'abord un pouvoir de \cle{coercition} : il peut punir, emprisonner, saisir. Ce fondement est réel : sans force, les lois resteraient des vœux pieux. Un État qui ne peut faire respecter ses décisions s'effondre, comme on l'observe dans les situations de guerre civile où l'autorité centrale disparaît.

Mais ce fondement a deux limites décisives. D'une part, comme l'a montré La Boétie, la force est numériquement insuffisante : elle suppose que la majorité accepte déjà d'obéir. D'autre part, un pouvoir purement fondé sur la terreur est instable : il engendre la haine, la révolte, et doit sans cesse augmenter la répression pour se maintenir. La force peut donc \emph{accompagner} l'État, mais elle ne peut pas le \emph{fonder} durablement. Il faut que la contrainte soit perçue comme légitime.

\courssub{Deuxième fondement : la tradition et le sacré}

Un deuxième fondement, très ancien, est la \cle{tradition} : on obéit parce que « cela a toujours été ainsi ». Le pouvoir est héréditaire, transmis de père en fils, et souvent sacralisé : le roi règne « par la grâce de Dieu », il est le représentant des ancêtres ou des dieux sur terre. La légitimité repose alors sur le respect des coutumes et des croyances religieuses.

\begin{exemplebox}[Les chefferies traditionnelles au Cameroun]
Le Cameroun offre des exemples vivants de légitimité traditionnelle. Dans les sultanats du Nord (Foumban, Garoua, Maroua), le lamido ou le sultan est à la fois chef politique et religieux. Dans les fondoms de l'Ouest (Bamoun, Bamiléké), le \emph{fon} ou le \emph{fo} est le gardien des coutumes, entouré de notables et de sociétés secrètes. Lors de funérailles traditionnelles ou de grandes cérémonies, l'autorité du chef s'impose naturellement, non par la force, mais parce qu'elle est enracinée dans l'histoire et le sacré. Nul ne discute la place du chef lors de la danse rituelle : sa légitimité est héritée et vécue comme évidente.
\end{exemplebox}

Ce fondement est puissant mais fragile à l'époque moderne : dès que les croyances partagées s'affaiblissent, dès que les individus demandent des \emph{raisons} d'obéir et non plus seulement des habitudes, la tradition ne suffit plus à justifier le pouvoir. C'est pourquoi la philosophie moderne a cherché un fondement rationnel : le contrat social.

\courssub{Troisième fondement : le contrat social}

L'idée du \cle{contrat social} est une hypothèse philosophique : imaginons les hommes avant toute société politique, dans ce qu'on appelle l'\cle{état de nature}. Pourquoi et comment auraient-ils accepté de former un État ? La réponse : par un \emph{pacte}, un accord par lequel chacun renonce à une part de sa liberté naturelle en échange des avantages de la vie commune. Trois grands philosophes ont développé cette hypothèse de manières différentes.

\textbf{Hobbes : sortir de la guerre de tous contre tous.} Dans le \emph{Léviathan} (1651), Thomas \textbf{Hobbes} décrit l'état de nature comme une situation terrible : les hommes, égaux en force et en désirs, sans autorité commune pour les arbitrer, sont en « guerre de tous contre tous ». Chacun craint pour sa vie, nul ne peut travailler, commercer, élever ses enfants en sécurité.

\begin{aretenir}[Citation à commenter : Hobbes]
« Hors de l'État, la vie de l'homme est solitaire, misérable, dangereuse, animale et brève » (\emph{Léviathan}, chap. XIII). \par
\emph{Commentaire} : Hobbes affirme que sans pouvoir commun capable de faire peur à tous, la vie humaine est invivable. L'argument repose sur la peur de la mort violente, passion universelle. Pour échapper à cette guerre, les hommes concluent un pacte : ils \emph{aliènent} (transfèrent) leur liberté et leur droit de se défendre eux-mêmes à un \cle{souverain} unique, le Léviathan, en échange de la \emph{sécurité}. Chez Hobbes, ce souverain est absolu : le contrat se fait entre les sujets, non avec le souverain, qui n'est donc jamais tenu de rendre des comptes.
\end{aretenir}

\textbf{Locke : protéger les droits naturels.} John \textbf{Locke}, dans le \emph{Traité du gouvernement civil} (1690), donne de l'état de nature une image moins sombre : les hommes y possèdent déjà des \cle{droits naturels} — la vie, la liberté, la propriété — mais ces droits sont mal protégés, car chacun est juge de ses propres conflits. Le contrat social a donc pour but de créer un juge impartial et des lois connues qui \emph{garantissent} ces droits. Conséquence révolutionnaire : le pouvoir est un \emph{dépôt} confié sous condition. Si l'État viole les droits qu'il devait protéger, le contrat est rompu et le peuple a le \cle{droit de résistance} : il peut renverser le gouvernement. Locke fonde ainsi le libéralisme politique et limite strictement le pouvoir.

\textbf{Rousseau : la souveraineté du peuple.} Jean-Jacques \textbf{Rousseau}, dans \emph{Du contrat social} (1762), pose le problème en des termes célèbres.

\begin{aretenir}[Citation à commenter : Rousseau]
« L'homme est né libre, et partout il est dans les fers » (\emph{Du contrat social}, livre I). \par
\emph{Commentaire} : paradoxe apparent — si l'homme est naturellement libre, comment se fait-il qu'il vive partout soumis à des lois et à des maîtres ? Rousseau veut rendre cette chaîne \emph{légitime}. Sa solution : le pacte par lequel chacun s'engage envers tous et se donne tout entier à la communauté. En obéissant à la loi, le citoyen n'obéit pas à un maître étranger mais à la \cle{volonté générale}, c'est-à-dire à la volonté du peuple souverain, donc en un sens à lui-même. La souveraineté appartient au peuple, elle est inaliénable et indivisible : aucun représentant, aucun roi ne peut s'en emparer. Rousseau fonde ainsi la légitimité démocratique.
\end{aretenir}

\begin{attention}[Erreur fréquente]
Ne croyez pas que le « contrat social » désigne un document historique réellement signé par nos ancêtres. Aucun historien n'a retrouvé un tel contrat. Il s'agit d'une \emph{hypothèse philosophique}, d'un outil de pensée : il permet de juger si un pouvoir est légitime en se demandant si des individus libres et rationnels auraient pu y consentir. Confondre l'hypothèse avec un fait historique est une faute grave en dissertation.
\end{attention}

\courssub{Quatrième fondement : le droit et la légitimité rationnelle-légale}

Le sociologue Max \textbf{Weber} complète l'analyse en distinguant trois types de domination légitime : la domination \emph{traditionnelle} (on obéit aux coutumes), la domination \emph{charismatique} (on obéit à un leader exceptionnel) et la domination \emph{rationnelle-légale}, propre aux États modernes. Dans celle-ci, on n'obéit ni à une tradition ni à une personne, mais à des \cle{règles impersonnelles} : des lois écrites, des procédures, des fonctions. Le président, le préfet ou le percepteur n'ont d'autorité que parce qu'une loi définit leur fonction et ses limites. Quand ils quittent leur poste, l'autorité passe à leur successeur, non à leur famille. C'est le fondement de l'État de droit : la loi est au-dessus de tous, y compris des gouvernants.

\begin{popfigure}
\begin{center}
\begin{tikzpicture}[
  box/.style={draw=popblue, fill=popblueL, rounded corners=2pt, align=center, font=\small, minimum height=1.15cm, text width=3.1cm},
  head/.style={draw=popdark, fill=popdark, text=white, rounded corners=2pt, align=center, font=\small\bfseries, minimum height=0.8cm, text width=3.1cm},
  lab/.style={font=\small\bfseries, text=popdark, align=left, text width=2.4cm}
]
\node[lab] at (0,4.4) {Philosophe};
\node[head] at (3.2,4.4) {Hobbes (1651)};
\node[head] at (6.9,4.4) {Locke (1690)};
\node[head] at (10.6,4.4) {Rousseau (1762)};
\node[head] at (14.3,4.4) {Weber (XX\textsuperscript{e} s.)};
\node[lab] at (0,3.0) {État de nature};
\node[box] at (3.2,3.0) {Guerre de tous contre tous ; vie brève et misérable};
\node[box] at (6.9,3.0) {Libre mais droits naturels mal protégés};
\node[box] at (10.6,3.0) {Homme né libre, bon, sans maître};
\node[box] at (14.3,3.0) {(Non retenu : analyse des types de domination)};
\node[lab] at (0,1.6) {But du pacte};
\node[box] at (3.2,1.6) {Obtenir la sécurité en échange de la liberté};
\node[box] at (6.9,1.6) {Protéger vie, liberté, propriété};
\node[box] at (10.6,1.6) {Rendre l'obéissance légitime : obéir à la volonté générale};
\node[box] at (14.3,1.6) {Obéir à des règles impersonnelles, non à des personnes};
\node[lab] at (0,0.2) {Souveraineté et conséquence};
\node[box] at (3.2,0.2) {Souverain absolu ; le citoyen doit obéir};
\node[box] at (6.9,0.2) {Pouvoir limité ; droit de résistance si les droits sont violés};
\node[box] at (10.6,0.2) {Peuple souverain ; citoyen législateur};
\node[box] at (14.3,0.2) {État de droit ; fonctionnaires soumis à la loi};
\end{tikzpicture}
\end{center}
\legende{Tableau comparatif des fondements de l'État : état de nature, but du pacte, souveraineté et conséquence pour le citoyen chez Hobbes, Locke, Rousseau et Weber.}
\end{popfigure}

\begin{popfigure}
\begin{center}
\begin{tikzpicture}[
  etape/.style={draw=popdark, fill=popcream, rounded corners=3pt, align=center, font=\small, minimum height=1cm, text width=2.9cm},
  fleche/.style={-{Stealth[length=3mm]}, very thick, poporange},
  variante/.style={font=\scriptsize, text=popmuted, align=center, text width=2.9cm}
]
\node[etape] (nature) at (0,0) {\textbf{État de nature}\par liberté sans loi, insécurité};
\node[etape] (pacte) at (4.6,0) {\textbf{Pacte / contrat}\par accord libre des individus};
\node[etape] (etat) at (9.2,0) {\textbf{État}\par pouvoir commun institué};
\node[etape] (fin) at (13.8,0) {\textbf{Sécurité et droits}\par paix, justice, liberté civile};
\draw[fleche] (nature) -- (pacte);
\draw[fleche] (pacte) -- (etat);
\draw[fleche] (etat) -- (fin);
\node[variante] at (0,-1.7) {Hobbes : guerre de tous contre tous};
\node[variante] at (4.6,-1.7) {Hobbes : aliénation totale ; Locke : dépôt sous condition ; Rousseau : don de soi à tous};
\node[variante] at (9.2,-1.7) {Hobbes : souverain absolu ; Locke : pouvoir limité ; Rousseau : peuple souverain};
\node[variante] at (13.8,-1.7) {Hobbes : sécurité ; Locke : droits naturels garantis ; Rousseau : liberté civile};
\end{tikzpicture}
\end{center}
\legende{Schéma de l'argument contractualiste : de l'état de nature à l'État, avec les trois variantes de Hobbes, Locke et Rousseau.}
\end{popfigure}

\courssub{Synthèse : légitimité contre domination}

L'enquête est terminée : ni la force seule, ni la tradition seule ne fondent durablement l'État moderne. Ce qui distingue un véritable État d'une bande armée qui rackette une population, c'est la \cle{légitimité}. Le pouvoir \emph{légal} est accepté parce qu'il est reconnu comme juste : il protège la sécurité (Hobbes), garantit les droits (Locke), exprime la volonté générale (Rousseau) et respecte des règles impersonnelles (Weber). La pure \emph{domination}, au contraire, est imposée par la seule contrainte : elle peut durer un temps, mais elle est instable et injuste. Retenons la formule : la force fait obéir, mais seule la légitimité fait consentir.

\begin{savaistu}[Deux légitimités au Cameroun]
L'article 65 de la Constitution camerounaise du 18 janvier 1996 reconnaît les \emph{chefferies traditionnelles} comme auxiliaires de l'administration. Ainsi coexistent officiellement deux légitimités : la légitimité rationnelle-légale de l'État moderne (préfets, sous-préfets, maires élus) et la légitimité traditionnelle des sultans, fons et chefs de village. Le chef traditionnel participe à la collecte de l'impôt, au règlement des conflits de voisinage et à l'organisation des cérémonies. Cette coexistence montre concrètement que les fondements du pouvoir étudiés dans cette leçon ne sont pas de simples idées abstraites : ils structurent encore la vie politique camerounaise aujourd'hui.
\end{savaistu}

\begin{methode}[Analyser un texte de philosophe sur l'État]
Pour le commentaire de texte au Baccalauréat technique, procédez en quatre étapes :
\begin{enumerate}
\item \textbf{Repérer la thèse} : quelle affirmation l'auteur défend-il ? (Ex. Hobbes : sans État, la vie est une guerre permanente.)
\item \textbf{Dégager l'argument principal} : quel raisonnement soutient la thèse ? (Ex. les hommes sont égaux en force, donc rivaux, donc ennemis sans pouvoir commun.)
\item \textbf{Identifier l'exemple ou l'image} : comment l'auteur illustre-t-il son idée ? (Ex. le Léviathan, « dieu mortel » duquel naît la paix.)
\item \textbf{Formuler l'objection possible} : quelle critique un adversaire ferait-il ? (Ex. Locke objecterait qu'un souverain absolu est plus dangereux que l'état de nature.)
\end{enumerate}
Cette démarche vous permet de rédiger une analyse structurée et de justifier votre conclusion, conformément aux savoir-faire exigés.
\end{methode}

\begin{exoresolu}[La Boétie et l'obéissance]
\emph{Énoncé.} Un élève affirme : « Si les Camerounais respectent le code de la route, c'est uniquement par peur des contrôles de gendarmerie. » À l'aide de La Boétie et de la notion de légitimité, montrer en quelques lignes que cette explication est insuffisante.
\tcblower
\emph{Solution détaillée.} L'explication par la peur est insuffisante pour deux raisons. Premièrement, la force est numériquement incapable de tout surveiller : les forces de l'ordre ne peuvent pas être présentes à chaque carrefour, à chaque heure, sur chaque route du pays. Or la plupart des conducteurs s'arrêtent au feu rouge même la nuit, en l'absence de tout contrôle. Deuxièmement, comme le montre La Boétie dans le \emph{Discours de la servitude volontaire}, l'obéissance repose toujours en partie sur le consentement : on obéit parce qu'on juge la règle utile et acceptable — ici, parce que le code de la route protège la vie de tous, y compris la nôtre. C'est donc la \emph{légitimité} de la règle, et non la seule crainte de la sanction, qui fonde l'obéissance durable. Un pouvoir qui ne reposerait que sur la peur s'effondrerait dès que la surveillance faiblirait.
\end{exoresolu}

\begin{exercice}[Pour s'entraîner]
\begin{enumerate}
\item Définir la légitimité et distinguer pouvoir légal et domination.
\item Pourquoi, selon Hobbes, les hommes acceptent-ils de renoncer à leur liberté naturelle ? Qu'obtiennent-ils en échange ?
\item En quoi le contrat social de Locke fonde-t-il un droit de résistance que Hobbes refuse ?
\item Expliquer la formule de Rousseau : « obéir à la volonté générale, c'est obéir à soi-même ».
\item Identifier, dans l'organisation politique de votre village ou de votre quartier, un exemple de légitimité traditionnelle et un exemple de légitimité rationnelle-légale. Justifier.
\end{enumerate}
\end{exercice}

\coursec{Les rôles de l'État : que fait concrètement l'État ?}

Dans la section précédente, nous avons cherché les \cle{fondements} de l'État : pourquoi obéit-on au pouvoir ? Nous avons vu que l'obéissance peut reposer sur la force, sur la tradition, sur la loi ou sur un contrat consenti. Mais une question demeure : une fois fondé et reconnu, que \emph{fait} réellement l'État ? Quelles tâches accomplit-il chaque jour, et jusqu'où doit-il intervenir ? Un élève de Yaoundé qui entre dans son lycée public, un commerçant de Douala qui décharge ses marchandises au port, un malade soigné à l'hôpital de district : tous rencontrent l'État sans toujours le savoir. Cette section décrit les \cle{rôles} concrets de l'État, puis ouvre le grand débat philosophique : l'État doit-il se contenter d'être un simple « gendarme », ou doit-il intervenir dans l'économie et la vie sociale ?

\courssub{Les fonctions régaliennes : le cœur du métier d'État}

Commençons par une intuition simple. Imagine qu'un conflit éclate entre deux voisins à Bafoussam au sujet d'une parcelle de terrain. Qui tranchera ? Ni les deux voisins eux-mêmes (chacun serait juge et partie), ni une entreprise privée (elle n'a aucune autorité reconnue), mais un \emph{tribunal}, c'est-à-dire une institution de l'État. De même, si un pays étranger menaçait le Cameroun, seule une armée nationale pourrait défendre le territoire. Il existe donc des missions que \emph{seul} l'État peut légitimement accomplir : ce sont les \cle{fonctions régaliennes} (du latin \emph{rex}, \emph{regis} : le roi, car elles appartenaient autrefois au souverain).

\begin{definition}[Les fonctions régaliennes]
On appelle \cle{fonctions régaliennes} les missions que seul l'État peut assumer : la \textbf{justice}, la \textbf{police}, l'\textbf{armée}, la \textbf{monnaie} et la \textbf{diplomatie}. Elles expriment la souveraineté de l'État et ne peuvent être abandonnées à des particuliers sans que l'État cesse d'être un État.
\end{definition}

Détaillons ces fonctions une à une, avec leurs réalisations camerounaises.

\begin{itemize}
\item \textbf{La sécurité intérieure.} La \emph{police} et la \emph{gendarmerie} maintiennent l'ordre public, protègent les personnes et les biens, enquêtent sur les infractions. Sans elles, la loi du plus fort remplacerait la loi tout court. Max Weber l'avait bien vu : l'État détient le \emph{monopole de la violence légitime} — seuls ses agents peuvent user de la force, et seulement dans le cadre de la loi.
\item \textbf{La défense.} L'\emph{armée} protège le territoire contre les menaces extérieures. Au Cameroun, les forces de défense sont engagées notamment dans la lutte contre Boko Haram dans l'Extrême-Nord : exemple tragique mais clair d'une mission qu'aucun acteur privé ne pourrait assumer.
\item \textbf{La justice.} Les \emph{tribunaux} règlent les conflits et punissent les délits selon des règles communes. La justice est la condition de la confiance : sans juge impartial, chacun se ferait justice soi-même, et la société retomberait dans la vengeance privée.
\item \textbf{La diplomatie.} Les \emph{ambassades} représentent le pays à l'étranger, négocient les traités, protègent les ressortissants. Le Cameroun entretient des relations diplomatiques avec de nombreux États et organisations internationales.
\item \textbf{La monnaie.} Battre monnaie est un attribut classique de la souveraineté. Le franc CFA, utilisé au Cameroun dans le cadre de la CEMAC (Communauté économique et monétaire de l'Afrique centrale), illustre que cette fonction peut aussi être exercée en commun par plusieurs États, qui la délèguent alors à une banque centrale commune (la BEAC).
\end{itemize}

\begin{popfigure}
\begin{center}
\begin{tikzpicture}[scale=1]
\node[circle, draw=popblue, fill=popblueL, thick, minimum size=2.2cm, align=center, font=\bfseries] (etat) at (0,0) {L'ÉTAT};
\node[draw=poporange, fill=poporangeL, rounded corners, align=center, font=\small] (sec) at (90:3.4) {\textbf{Sécurité}\\ Police, gendarmerie\\ (commissariats)};
\node[draw=popgreen, fill=popgreenL, rounded corners, align=center, font=\small] (jus) at (162:3.6) {\textbf{Justice}\\ Tribunaux\\ (TPI, cours)};
\node[draw=poppurple, fill=poppurpleL, rounded corners, align=center, font=\small] (eco) at (234:3.6) {\textbf{Économie}\\ Impôts, routes, monnaie\\ (port de Douala)};
\node[draw=poppink, fill=poppinkL, rounded corners, align=center, font=\small] (soc) at (306:3.6) {\textbf{Social}\\ Écoles, hôpitaux\\ (lycées publics)};
\node[draw=popgold, fill=popgoldL, rounded corners, align=center, font=\small] (cul) at (18:3.6) {\textbf{Culture}\\ Drapeau, hymne\\ (fête du 20 mai)};
\draw[->, thick, popblue] (etat) -- (sec);
\draw[->, thick, popblue] (etat) -- (jus);
\draw[->, thick, popblue] (etat) -- (eco);
\draw[->, thick, popblue] (etat) -- (soc);
\draw[->, thick, popblue] (etat) -- (cul);
\end{tikzpicture}
\end{center}
\legende{Schéma en étoile des cinq grands rôles de l'État, avec un exemple camerounais par branche. Les fonctions régaliennes (sécurité, justice, défense) forment le noyau historique ; les rôles économique, social et culturel se sont ajoutés progressivement.}
\end{popfigure}

\courssub{Le rôle économique : l'État acteur de la vie matérielle}

Au-delà de ses fonctions régaliennes, l'État intervient dans l'\cle{économie}. Observons concrètement : comment l'État se procure-t-il des ressources, et qu'en fait-il ?

\textbf{Les impôts et le budget.} L'État prélève des \emph{impôts} et des taxes (sur les revenus, sur les marchandises, sur les entreprises). Ces recettes alimentent le \emph{budget de l'État}, voté chaque année dans la loi de finances. Le budget est le miroir fidèle des priorités politiques d'un pays : dire ce qu'un État finance, c'est dire ce qu'il juge important.

\textbf{Les infrastructures.} L'État construit et entretient les équipements collectifs que nul particulier ne financerait seul : les \emph{routes} (par exemple l'axe Yaoundé--Douala), les \emph{ports} de Douala et de Kribi (port en eau profonde, moteur du commerce extérieur), les \emph{barrages hydroélectriques} (Song Loulou et Édéa sur la Sanaga, Memve'ele sur la Ntem) qui fournissent l'électricité. Ces infrastructures profitent à tous et rendent possible l'activité économique privée : sans route, pas de transport des marchandises ; sans électricité, pas d'usine.

\textbf{La régulation.} L'État fixe les règles du jeu économique : droit du travail, normes d'hygiène, lutte contre les monopoles abusifs, protection des consommateurs. Il peut aussi soutenir des secteurs (agriculture, petites entreprises) par des subventions ou des prix garantis.

\begin{exemplebox}[Un budget parlant : les priorités chiffrées de l'État camerounais]
Le budget de l'État du Cameroun dépasse \textbf{6 000 milliards de FCFA} (lois de finances récentes). L'\emph{éducation} et la \emph{défense} figurent parmi les premiers postes de dépense, devant la santé et les infrastructures. Une part importante sert aussi au remboursement de la \emph{dette publique}. Lire un budget, c'est donc lire les choix politiques d'un État : investir dans les écoles, c'est parier sur la jeunesse ; financer la défense, c'est répondre à des menaces réelles.
\end{exemplebox}

Pour objectiver ces priorités, analysons une répartition indicative d'un budget de l'État.

\begin{popfigure}
\begin{center}
\begin{tikzpicture}[scale=1.1]
\fill[popblue] (0,0) -- (90:2.6) arc (90:18:2.6) -- cycle;
\fill[poporange] (0,0) -- (18:2.6) arc (18:-36:2.6) -- cycle;
\fill[popgreen] (0,0) -- (-36:2.6) arc (-36:-72:2.6) -- cycle;
\fill[poppurple] (0,0) -- (-72:2.6) arc (-72:-126:2.6) -- cycle;
\fill[popgold] (0,0) -- (-126:2.6) arc (-126:-162:2.6) -- cycle;
\fill[popmuted] (0,0) -- (-162:2.6) arc (-162:-270:2.6) -- cycle;
\draw[thick] (0,0) circle (2.6);
\node[font=\small\bfseries, white] at (54:1.6) {20\,\%};
\node[font=\small\bfseries, white] at (-9:1.6) {15\,\%};
\node[font=\small\bfseries, white] at (-54:1.6) {10\,\%};
\node[font=\small\bfseries, white] at (-99:1.6) {15\,\%};
\node[font=\small\bfseries, white] at (-144:1.6) {10\,\%};
\node[font=\small\bfseries, white] at (144:1.6) {30\,\%};
\node[right, font=\small] at (3.1,1.8) {\textcolor{popblue}{\rule{0.35cm}{0.35cm}} Éducation : 20\,\%};
\node[right, font=\small] at (3.1,1.2) {\textcolor{poporange}{\rule{0.35cm}{0.35cm}} Défense et sécurité : 15\,\%};
\node[right, font=\small] at (3.1,0.6) {\textcolor{popgreen}{\rule{0.35cm}{0.35cm}} Santé : 10\,\%};
\node[right, font=\small] at (3.1,0.0) {\textcolor{poppurple}{\rule{0.35cm}{0.35cm}} Infrastructures : 15\,\%};
\node[right, font=\small] at (3.1,-0.6) {\textcolor{popgold}{\rule{0.35cm}{0.35cm}} Dette : 10\,\%};
\node[right, font=\small] at (3.1,-1.2) {\textcolor{popmuted}{\rule{0.35cm}{0.35cm}} Autres (administration, culture\dots) : 30\,\%};
\end{tikzpicture}
\end{center}
\legende{Répartition indicative (fictive mais réaliste) d'un budget de l'État. Vérification : $20 + 15 + 10 + 15 + 10 + 30 = 100\,\%$. L'éducation et la défense occupent les premières places ; la dette absorbe une part non négligeable, ce qui réduit les marges de manœuvre pour le social.}
\end{popfigure}

\begin{methode}[Analyser un document chiffré sur le budget de l'État]
\begin{enumerate}
\item \textbf{Lire} les données : identifier les grandeurs (milliards de FCFA, pourcentages) et les postes de dépense.
\item \textbf{Comparer} : classer les postes du plus important au moins important ; calculer des rapports (l'éducation représente ici le double de la santé : $20\,\%$ contre $10\,\%$).
\item \textbf{Interpréter} : en déduire les priorités politiques de l'État (investissement dans la jeunesse, réponse aux menaces sécuritaires).
\item \textbf{Nuancer} : rappeler qu'un budget est un choix discutable, pas une fatalité — d'autres répartitions sont possibles.
\end{enumerate}
\end{methode}

\courssub{Le rôle social et éducatif : vers l'État-providence}

L'État ne se contente plus, depuis le XX\textsuperscript{e} siècle, de maintenir l'ordre. Il prend en charge le \emph{bien-être} de la population.

\begin{definition}[L'État-providence]
L'\cle{État-providence} (ou État social) est un État qui intervient pour garantir le bien-être social — \textbf{éducation}, \textbf{santé}, \textbf{aide aux démunis} — au-delà de ses fonctions régaliennes. Il part de l'idée que l'égalité formelle (mêmes droits pour tous) ne suffit pas : il faut aussi donner à chacun des moyens réels de vivre dignement.
\end{definition}

Concrètement, au Cameroun : les \emph{écoles et universités publiques} (écoles primaires publiques, lycées, universités de Yaoundé, Douala, Dschang\dots) rendent l'instruction accessible à tous, y compris aux familles pauvres ; les \emph{hôpitaux publics} et centres de santé intégrés soignent à des tarifs modérés ; des programmes d'aide soutiennent les plus démunis. La justification philosophique est forte : comme le souligne Rousseau, la liberté véritable suppose des conditions matérielles minimales. Un enfant qui n'a pas accès à l'école est formellement libre, mais réellement empêché de développer ses capacités. L'État-providence transforme ainsi les droits proclamés en droits effectifs.

\courssub{Le rôle culturel et symbolique : fabriquer une nation}

Un État ne gouverne pas seulement des corps, il unit des esprits. Il entretient les \emph{symboles} qui font d'une population une \emph{nation} : le \emph{drapeau} (vert, rouge, jaune avec l'étoile), l'\emph{hymne national} (« Ô Cameroun, berceau de nos ancêtres »), la \emph{fête nationale du 20 mai}, qui célèbre l'unité du pays. Au Cameroun, la promotion du \emph{bilinguisme français--anglais} répond à la même finalité : rassembler deux héritages linguistiques dans une même citoyenneté. L'école publique joue ici un rôle décisif : en enseignant une histoire et des valeurs communes, elle forme des citoyens et pas seulement des élèves. Ce rôle symbolique n'est pas un luxe : sans sentiment d'appartenance commune, l'obéissance aux lois redeviendrait pure contrainte.

\courssub{Le grand débat : État minimal ou État interventionniste ?}

Maintenant que nous avons décrit ce que fait l'État, la question philosophique peut être posée dans toute sa force : tout cela \emph{doit}-il relever de l'État ? Deux thèses s'affrontent.

\textbf{La thèse libérale : l'État minimal.} Pour Adam Smith (XVIII\textsuperscript{e} siècle), l'économie fonctionne mieux lorsque le marché est libre : la « main invisible » des intérêts individuels produit l'harmonie générale mieux que toute planification étatique. L'État doit donc se limiter à trois tâches : la défense, la justice, et quelques travaux publics indispensables. Au XX\textsuperscript{e} siècle, Robert Nozick radicalise cette position : toute intervention de l'État au-delà de la protection des personnes et des contrats viole les droits individuels ; l'impôt redistributif serait une forme de travail forcé. C'est la thèse de l'\emph{État minimal}, dit « État gendarme ».

\begin{savaistu}[D'où vient l'expression « État gendarme » ?]
L'expression « État gendarme » vient des libéraux du XIX\textsuperscript{e} siècle : l'État devrait se contenter de faire respecter l'ordre, comme un \emph{gendarme} de service, sans intervenir dans l'économie. Ses adversaires socialistes utilisaient d'ailleurs l'expression de manière critique : un État qui ne fait que surveiller et punir abandonne les plus faibles à la misère.
\end{savaistu}

\textbf{La thèse interventionniste.} Face aux crises du capitalisme (notamment celle de 1929), John Maynard Keynes montre que le marché livré à lui-même peut produire chômage de masse et misère : l'État doit alors intervenir par la dépense publique pour relancer l'activité et soutenir l'emploi. La \emph{social-démocratie} prolonge cette idée : l'État doit corriger les inégalités que le marché engendre, par l'impôt progressif, l'éducation et la santé publiques. Sans intervention, arguent ses partisans, la liberté des plus pauvres reste un mot vide.

\textbf{Discussion.} Les deux thèses ont des arguments sérieux. Le libéral a raison de rappeler que l'État peut être inefficace, lourd, parfois prédateur ; l'interventionniste a raison de rappeler que le marché seul ne garantit ni l'éducation des pauvres ni la santé de tous. Le débat n'est pas tranché une fois pour toutes : il se rejoue dans chaque pays, à chaque élection, à chaque loi de finances. Au Cameroun, il prend la forme de questions concrètes : faut-il subventionner les carburants ? Privatiser certaines entreprises publiques ? Augmenter les bourses des étudiants ?

\begin{exoresolu}[Disserter sur les rôles de l'État]
\textbf{Sujet :} L'État doit-il se contenter d'assurer l'ordre ? Dégagez l'intérêt philosophique de cette question et proposez un plan de réponse.
\tcblower
\textbf{Intérêt philosophique.} La question oppose deux conceptions de l'État : l'État minimal (se limiter à l'ordre) et l'État interventionniste (agir aussi pour le bien-être). Elle touche au sens même de la liberté : est-ce être libre que d'être simplement protégé, ou faut-il des conditions matérielles (école, santé) pour exercer sa liberté ?

\textbf{Plan possible.}
\begin{enumerate}
\item \emph{Thèse :} l'État gendarme suffit. Argument de Smith : la « main invisible » du marché ; argument de Nozick : toute redistribution viole les droits individuels. Exemple : l'efficacité des entreprises privées.
\item \emph{Antithèse :} l'ordre sans justice sociale est une paix injuste. Argument de Keynes : le marché livré à lui-même produit crises et chômage ; argument de Rousseau : la liberté réelle suppose l'éducation et des conditions de vie dignes. Exemple : sans écoles publiques, les enfants pauvres restent exclus.
\item \emph{Synthèse :} l'État doit garantir l'ordre (c'est son cœur régalien) mais aussi créer les conditions de la liberté effective ; l'intervention doit cependant rester contrôlée pour éviter l'arbitraire et l'inefficacité.
\end{enumerate}
\end{exoresolu}

\courssub{Limites et dérives : pourquoi l'État doit être contrôlé}

Décrire les rôles de l'État ne doit pas faire oublier ses \cle{dérives} possibles. L'État dispose d'une puissance immense — la force légitime, l'impôt, le droit — et toute puissance peut être mal employée. Trois dérives classiques :

\begin{itemize}
\item La \textbf{bureaucratie} : lourdeurs administratives, lenteur des procédures, multiplication des formulaires qui découragent les citoyens et les entrepreneurs.
\item La \textbf{corruption} : détournement des deniers publics, favoritisme dans les marchés publics, « petits arrangements » qui volent à l'école ou à l'hôpital les moyens qui leur étaient destinés.
\item L'\textbf{abus de pouvoir} : lorsque l'État utilise sa force contre les libertés au lieu de les protéger (arrestations arbitraires, censure).
\end{itemize}

D'où la nécessité du \emph{contrôle de l'État} : la \emph{séparation des pouvoirs} (exécutif, législatif, judiciaire), théorisée par Montesquieu dans \emph{De l'esprit des lois}, fait que « le pouvoir arrête le pouvoir » ; la \emph{presse libre} informe et dénonce ; la \emph{société civile} (associations, syndicats, citoyens vigilants) surveille l'action publique. Un État fort n'est pas un État sans contrôle : c'est un État capable d'accomplir ses rôles \emph{sous le regard} des citoyens.

\begin{attention}[Ne pas tout attribuer à l'État]
Erreur fréquente dans les copies : attribuer à l'État ce qui relève des \textbf{communes} (marchés locaux, voirie de quartier, état civil de proximité), des \textbf{entreprises privées} (production de biens, emplois marchands) ou des \textbf{ONG} (actions humanitaires ponctuelles). Avant d'écrire « l'État fait\dots », vérifie toujours \emph{qui agit réellement}. Citer un exemple précis et correct (le port de Kribi, un lycée public, un tribunal) vaut mieux qu'une liste vague.
\end{attention}

\begin{aretenir}[L'essentiel de la section]
\begin{itemize}
\item Les \cle{fonctions régaliennes} (justice, police, armée, monnaie, diplomatie) sont le cœur irremplaçable de l'État.
\item L'État joue aussi un rôle \cle{économique} (impôts, budget, infrastructures comme les ports de Douala et Kribi), \cle{social} (écoles, hôpitaux : l'État-providence) et \cle{culturel} (drapeau, hymne, fête du 20 mai, bilinguisme).
\item Le débat philosophique oppose l'\cle{État minimal} (Smith, Nozick) à l'\cle{État interventionniste} (Keynes, social-démocratie).
\item L'analyse du budget (plus de 6 000 milliards de FCFA au Cameroun) permet d'objectiver les priorités de l'État.
\item Les dérives (bureaucratie, corruption, abus de pouvoir) rendent nécessaire le \cle{contrôle} de l'État : séparation des pouvoirs, presse, société civile.
\end{itemize}
\end{aretenir}

\begin{exercice}[Pour s'entraîner]
\begin{enumerate}
\item Définis les fonctions régaliennes et donne, pour chacune, un exemple camerounais précis.
\item À partir du diagramme circulaire de cette section, rédige un paragraphe de cinq lignes qui dégage les priorités de l'État et les discute.
\item Sujet de dissertation : « Plus d'État, est-ce plus de liberté ? » Rédige l'introduction (accroche, définitions, problématique, annonce du plan).
\end{enumerate}
\end{exercice}

\coursec{La finalité de l'État : pour quoi l'État existe-t-il ?}

Nous venons de voir ce que l'État \emph{fait} concrètement : il assure la sécurité, rend la justice, organise l'éducation, construit les routes, collecte les impôts. Mais une question plus profonde demeure : tout cela, \emph{pour quoi} ? À quoi sert finalement l'État ? Est-il lui-même la fin suprême, auquel cas l'individu n'existerait que pour le servir ? Ou bien n'est-il qu'un instrument, un moyen au service des hommes ? C'est le problème de la \cle{finalité} de l'État : non plus \og comment agit-il ? \fg{}, mais \og dans quel but existe-t-il ? \fg{}. Cette question est décisive, car de la réponse dépend la place de l'individu dans la cité : sujet obéissant ou citoyen libre ?

\courssub{Problématique : l'État, fin ou moyen ?}

\begin{definition}[Finalité de l'État]
La \cle{finalité} de l'État désigne le but ultime que poursuit l'organisation politique. Deux réponses s'opposent : ou bien l'État est une \cle{fin en soi} (l'individu sert l'État), ou bien il est un \cle{moyen} au service des individus (l'État sert l'homme).
\end{definition}

Prenons une image simple. Une machette est un \emph{moyen} : elle sert à couper, elle n'a de valeur que par l'usage qu'on en fait. En revanche, la santé est une \emph{fin} : on ne cherche pas la santé \og pour autre chose \fg{}, elle est désirable en elle-même. L'État est-il comme la machette (un outil) ou comme la santé (un bien absolu) ? Si l'État est une fin, alors sacrifier des individus pour lui peut se justifier. Si l'État est un moyen, alors il doit être jugé, contrôlé, voire changé dès qu'il ne sert plus les hommes. Toute la philosophie politique se joue dans cette alternative.

\courssub{Thèse 1 : l'État comme fin en soi (Hegel)}

Pour le philosophe allemand \cle{Hegel} (1770-1831), l'État n'est pas un simple instrument : il est \og la réalisation de la liberté concrète \fg{}. L'idée est subtile. L'individu isolé, livré à ses désirs immédiats, n'est pas vraiment libre : il est esclave de ses passions. C'est en appartenant à une communauté organisée par des lois rationnelles que l'homme devient pleinement lui-même. La famille, la société civile et l'État sont les étapes par lesquelles l'esprit humain s'accomplit. L'État est donc, chez Hegel, \og la marche de Dieu dans le monde \fg{} : il réalise l'Idée rationnelle dans l'histoire.

Conséquence : l'individu trouve sa dignité en servant l'État, et l'État a une valeur supérieure à celle de chaque citoyen pris isolément. L'État est une \emph{fin}.

\begin{attention}[La dérive totalitaire]
Poussée à l'extrême, la thèse de l'État-fin conduit au \cle{totalitarisme} : si l'État vaut plus que les individus, alors on peut sacrifier les individus à l'État. Les régimes qui ont érigé l'État (ou la race, ou le parti) en absolu ont justifié les pires crimes. Il faut donc manier la thèse hégélienne avec prudence : Hegel pensait un État \emph{rationnel et de droit}, non une dictature.
\end{attention}

\begin{savaistu}[Les totalitarismes du XX\textsuperscript{e} siècle]
Le nazisme en Allemagne (1933-1945) et le stalinisme en URSS illustrent la dérive de l'État érigé en fin absolue : au nom de la race ou de la révolution, des millions d'êtres humains furent déportés et exterminés. L'individu n'était plus qu'un moyen, un rouage sacrifiable. C'est en réaction à ces catastrophes que furent proclamées les grandes déclarations des droits de l'homme après 1945, dont la Déclaration universelle des droits de l'homme de 1948.
\end{savaistu}

\courssub{Thèse 2 : l'État comme moyen (Locke, Kant)}

À l'opposé, la tradition \cle{contractualiste} et \cle{libérale} affirme que l'État n'est qu'un moyen. Pour \cle{Locke} (1632-1704), les hommes, libres et égaux à l'état de nature, concluent un \emph{contrat} pour créer un pouvoir chargé d'une seule mission : protéger leurs droits naturels — la vie, la liberté, la propriété. L'État est donc un instrument de garantie. S'il trahit sa mission, le peuple a le droit de le dissoudre : c'est le fondement du droit de résistance à l'oppression.

\cle{Kant} (1724-1804) radicalise cette idée par son principe moral célèbre : il faut traiter l'humanité, en soi et en autrui, \og toujours en même temps comme une fin, et jamais simplement comme un moyen \fg{}.

\begin{aretenir}[Citation à commenter : Kant]
\og Agis de telle sorte que tu traites l'humanité, aussi bien dans ta personne que dans toute autre, toujours en même temps comme une fin, et jamais simplement comme un moyen. \fg{} (Kant, \emph{Fondements de la métaphysique des mœurs}, 1785). Appliqué à la politique : un État légitime ne peut jamais utiliser les citoyens comme de simples instruments (chair à canon, main-d'œuvre forcée, boucs émissaires). Chaque homme a une \cle{dignité} absolue, et l'État existe pour la respecter et la protéger — non l'inverse.
\end{aretenir}

Pour le libéralisme, l'État est donc comparable à un gardien ou à un arbitre : nécessaire, mais subordonné à ceux qu'il sert. C'est la conception qui fonde les États de droit modernes, les constitutions et les déclarations des droits.

\courssub{Thèse 3 : l'État au service du bien commun (Aristote)}

Entre les deux, une troisième voie, plus ancienne : celle d'\cle{Aristote} (384-322 av. J.-C.). Pour lui, l'homme est un \og animal politique \fg{} : il ne peut vivre pleinement qu'en société. La cité (l'État) ne vise pas seulement la survie ou la sécurité : elle vise la \emph{vie bonne}, c'est-à-dire le bonheur (l'\emph{eudaimonia}) accompli par la vertu. L'État a donc une mission positive : rendre les citoyens meilleurs et heureux, non pas seulement les protéger.

\begin{definition}[Le bien commun]
Le \cle{bien commun} est l'ensemble des conditions sociales qui permettent à tous les membres de la société d'atteindre leur épanouissement : la paix, la justice, la sécurité, l'éducation, la santé, le travail. Il ne se réduit ni à la somme des intérêts privés, ni à l'intérêt de l'État lui-même : c'est le bien de \emph{tous et de chacun}.
\end{definition}

Cette tradition, reprise par la doctrine sociale, affirme que l'État doit promouvoir la justice et le bonheur de tous, en particulier des plus faibles. L'État n'est ni un maître à adorer (contre la dérive totalitaire), ni un simple gendarme minimal (contre le libéralisme strict) : il est le \emph{serviteur du bien commun}.

\courssub{Position critique : les anarchistes}

Reste une objection radicale : et si l'État, quel qu'il soit, était par nature oppressif ? C'est la thèse des \cle{anarchistes}. Pour \cle{Bakounine} (1814-1876), \og l'État est la plus flagrante, la plus cynique et la plus complète négation de l'humanité \fg{}. Pour \cle{Proudhon} (1809-1865), être gouverné, c'est être \og noté, enregistré, recensé, taxé, commandé \fg{}. Tout pouvoir, disent-ils, corrompt ; toute autorité étatique finit par écraser la liberté. Il faut donc supprimer l'État et le remplacer par l'entraide et l'association libre.

\textbf{Objection (Hobbes).} Mais que se passerait-il sans État ? \cle{Hobbes} répond : l'état de nature est \og la guerre de tous contre tous \fg{}, où la vie de l'homme est \og solitaire, misérable, dangereuse, animale et brève \fg{}. Sans pouvoir commun, c'est la loi du plus fort. L'anarchie risque donc de produire non l'harmonie, mais le chaos. Le débat reste ouvert : les anarchistes répliquent que la violence de l'État a historiquement fait plus de victimes que la violence privée.

\begin{attention}[Erreur fréquente]
Ne pas opposer \emph{totalement} liberté individuelle et État, comme si toute loi était une atteinte à la liberté. Pour les contractualistes (Locke, Rousseau, Kant), l'État bien fondé \emph{protège} la liberté au lieu de la détruire : la loi m'empêche de nuire, mais elle empêche aussi autrui de me nuire. La vraie opposition n'est pas \og liberté contre État \fg{}, mais \og État légitime (qui sert la liberté) contre État arbitraire (qui l'écrase) \fg{}.
\end{attention}

\begin{popfigure}
\begin{center}
\begin{tikzpicture}[scale=1, every node/.style={font=\small}]
% pivot
\fill[popdark] (0,-0.3) -- (-0.35,-1.6) -- (0.35,-1.6) -- cycle;
\fill[popdark] (0,0) circle (0.09);
% fléau
\draw[line width=2pt,popdark] (-4.2,0.35) -- (4.2,0.35);
% fils
\draw[popmuted] (-4.2,0.35) -- (-4.2,-0.55);
\draw[popmuted] (4.2,0.35) -- (4.2,-0.55);
% plateaux
\draw[fill=poppinkL,draw=poppink] (-5.6,-0.55) arc (180:360:1.4 and 0.45) -- cycle;
\draw[fill=popblueL,draw=popblue] (2.8,-0.55) arc (180:360:1.4 and 0.45) -- cycle;
\node[align=center] at (-4.2,-1.45) {\textbf{État fin en soi}\\ Hegel\\ dérive : totalitarismes};
\node[align=center] at (4.2,-1.45) {\textbf{État moyen}\\ Locke, Kant\\ droits de l'homme};
% curseur bien commun
\draw[fill=popgold,draw=popgold] (0.9,0.35) circle (0.14);
\node[align=center,popgold!60!black] at (0.9,1.05) {\textbf{Curseur :}\\ \textbf{le bien commun}};
\draw[->,popgold!60!black,line width=1pt] (0.9,0.85) -- (0.9,0.52);
\node[align=center,popmuted] at (0,-2.15) {Un État légitime équilibre autorité et liberté par le bien commun};
\end{tikzpicture}
\end{center}
\legende{La balance des finalités : l'État comme fin en soi (Hegel, dérive totalitaire) contre l'État comme moyen (Locke, Kant, droits de l'homme). Le curseur du bien commun indique la position d'équilibre : l'État sert la dignité humaine sous le contrôle du droit.}
\end{popfigure}

\begin{popfigure}
\begin{center}
\begin{tikzpicture}[scale=1, every node/.style={font=\footnotesize}]
\draw[->,line width=1.2pt,popdark] (-6.4,0) -- (6.6,0) node[right]{temps};
\foreach \x/\nom/\date/\pos in {
 -5.6/{Aristote}/{IV\textsuperscript{e} s. av. J.-C.}/above,
 -3.9/{Machiavel}/{1532}/below,
 -2.2/{Hobbes}/{1651}/above,
 -0.5/{Locke}/{1690}/below,
 1.2/{Rousseau}/{1762}/above,
 2.9/{Hegel}/{1821}/below,
 4.6/{Anarchistes}/{XIX\textsuperscript{e} s.}/above,
 6.0/{État de droit}/{XX\textsuperscript{e} s.}/below}{
 \fill[popblue] (\x,0) circle (0.07);
 \node[align=center,\pos=2pt] at (\x,0) {\textbf{\nom}\\ \date};
}
\end{tikzpicture}
\end{center}
\legende{Frise des grandes conceptions de la finalité de l'État : de la cité visant la vie bonne (Aristote) à l'État de droit contemporain, en passant par la raison d'État (Machiavel), la sécurité (Hobbes), les droits naturels (Locke), la volonté générale (Rousseau), l'État rationnel (Hegel) et la critique anarchiste.}
\end{popfigure}

\courssub{Application au Cameroun : une finalité de bien commun}

Le Cameroun inscrit explicitement sa vie politique dans une finalité de bien commun. Sa devise, \og \cle{Paix-Travail-Patrie} \fg{}, affiche trois biens collectifs : la paix (condition de toute vie commune), le travail (chemin de développement et de dignité), la patrie (communauté à laquelle chacun appartient). Le préambule de la Constitution de 1996 affirme les droits de l'homme et proclame que le peuple est la source du pouvoir : l'État y apparaît comme un \emph{moyen} au service de la nation, non comme une fin en soi.

De même, la vision \og Cameroun émergent à l'horizon 2035 \fg{} définit des objectifs de développement — croissance, industrialisation, lutte contre la pauvreté — qui donnent à l'action de l'État un but mesurable : améliorer les conditions de vie de tous les Camerounais. Construire une route entre Yaoundé et Douala, ouvrir un lycée technique à Maroua ou un hôpital à Bamenda, ce n'est pas glorifier l'État : c'est créer les conditions du bien commun. Inversement, lorsque des deniers publics sont détournés, c'est la finalité même de l'État qui est trahie : l'outil cesse de servir tous pour servir quelques-uns.

\begin{exoresolu}[Analyser la finalité d'une action de l'État]
Énoncé : l'État camerounais construit un lycée technique dans une localité rurale et y affecte des enseignants. Montrez que cette action peut être analysée selon les trois conceptions de la finalité de l'État, puis dites laquelle rend le mieux compte de sa légitimité.
\tcblower
\textbf{Solution détaillée.}
\begin{enumerate}
\item \emph{Selon Hegel} (État fin) : le lycée forme des citoyens qui s'accomplissent dans la communauté nationale ; instruire la jeunesse, c'est réaliser l'esprit du peuple. L'individu s'élève en servant la patrie éduquée.
\item \emph{Selon Locke et Kant} (État moyen) : le lycée est un service rendu aux individus ; il garantit un droit (l'éducation) et traite chaque enfant comme une fin, en développant ses capacités, non comme un simple instrument.
\item \emph{Selon Aristote} (bien commun) : le lycée crée une condition d'épanouissement pour tous — instruction, insertion professionnelle, paix sociale — et vise la \og vie bonne \fg{} de la communauté.
\end{enumerate}
\textbf{Conclusion.} C'est la conception du bien commun qui rend le mieux compte de la légitimité de cette action : l'État y apparaît comme un moyen ordonné à l'épanouissement de tous, ce qui rejoint la thèse libérale (il sert les individus) sans tomber dans la sacralisation de l'État (qui justifierait tout, même l'injuste).
\end{exoresolu}

\courssub{Bilan : vers un État légitime}

\begin{methode}[Construire une discussion philosophique : thèse, antithèse, synthèse]
Pour traiter une dissertation sur la finalité de l'État :
\begin{enumerate}
\item \textbf{Thèse} : l'État est une fin en soi — l'individu s'accomplit dans la cité (Hegel). Argument : l'isolement est impuissance ; la loi réalise la liberté.
\item \textbf{Antithèse} : l'État n'est qu'un moyen — il existe pour garantir sécurité, liberté et droits (Locke, Kant). Argument : l'homme a une dignité ; nul ne peut être traité comme simple instrument.
\item \textbf{Discussion} : la thèse 1 risque le totalitarisme ; la thèse 2 risque de réduire l'État à un gendarme minimal incapable de justice sociale ; la position anarchiste (supprimer l'État) se heurte au chaos de l'état de nature (Hobbes).
\item \textbf{Synthèse argumentée} : l'État légitime est un \emph{moyen au service du bien commun}, c'est-à-dire de la dignité de tous, et il doit être \emph{contrôlé par le droit et par les citoyens} (constitution, séparation des pouvoirs, élections, liberté de la presse).
\end{enumerate}
\end{methode}

\begin{aretenir}[L'essentiel]
La finalité de l'État fait débat : fin en soi (Hegel, avec le risque totalitaire), simple moyen de protection des droits (Locke, Kant), ou serviteur du bien commun (Aristote). Les anarchistes (Bakounine, Proudhon) récusent l'État lui-même, mais Hobbes objecte le chaos de l'état de nature. Position de synthèse : un État est \cle{légitime} lorsqu'il demeure un \emph{moyen} au service de la dignité humaine et du bien commun, limité par le droit et contrôlé par les citoyens. La devise camerounaise \og Paix-Travail-Patrie \fg{} et l'objectif d'émergence en 2035 expriment précisément une telle finalité de bien commun.
\end{aretenir}

\coursec{TP guidé : construire des outils d'analyse de l'État camerounais}

Dans les sections précédentes, nous avons montré que l'État n'existe pas pour lui-même : il a une \cle{finalité}, c'est-à-dire une fin, un but à atteindre — la sécurité pour Hobbes, la protection des droits naturels pour Locke, la liberté et l'égalité civiles pour Rousseau. Mais comment vérifier concrètement qu'un État réel, comme l'État camerounais, remplit bien ses rôles et sert bien le bien commun ? Il ne suffit pas de répéter des définitions : il faut \emph{observer}, \emph{classer}, \emph{mesurer}. C'est précisément l'objet de ce travail pratique : fabriquer, en groupes, des \cle{outils d'analyse} (carte mentale, tableau comparatif, diagramme statistique, croquis institutionnel) qui permettront de répondre avec des preuves à la question : l'État camerounais fait-il ce pour quoi il existe ?

\courssub{TP 1 — Construire une carte mentale de l'État camerounais}

\textbf{Objectif.} Organiser visuellement tout ce que nous savons de l'État camerounais : ses éléments constitutifs, ses institutions, ses rôles et ses symboles.

\textbf{Matériel et documents.} Une feuille A3, des feutres de couleurs, un extrait du préambule de la Constitution de 1996 (loi n° 96-06 du 18 janvier 1996), des photos d'institutions (le Palais de l'Unité, l'Assemblée nationale, la Cour suprême, un commissariat, une école publique).

\textbf{Protocole.} Chaque groupe part du nœud central « L'État » et fait pousser quatre branches principales : \emph{définition}, \emph{fondements}, \emph{rôles}, \emph{finalité}. Chaque document doit être « accroché » à une branche : la photo de l'Assemblée nationale illustre le pouvoir législatif, l'extrait constitutionnel illustre la finalité.

\begin{methode}[Construire une carte mentale efficace]
\begin{enumerate}
\item Placer le \cle{mot-clé} central au milieu de la feuille, dans un cadre visible.
\item Tracer une branche par grande idée (ici : définition, fondements, rôles, finalité) avec \textbf{un seul mot-clé par branche}.
\item Respecter une hiérarchie claire : aller \textbf{du général au particulier} (branche principale, sous-branches, puis feuilles).
\item Ajouter \textbf{un exemple concret par notion} : pour « souveraineté », écrire « le Cameroun décide seul de ses lois » ; pour « territoire », écrire « les frontières reconnues internationalement » ; pour « population », écrire « les citoyens camerounais ».
\item Relier chaque document fourni à la branche qu'il justifie, et non le recopier en entier.
\end{enumerate}
\end{methode}

\begin{popfigure}
\begin{center}
\begin{tikzpicture}[mindmap, grow cyclic, every node/.style={concept}, text width=2.6cm, align=center,
root concept/.append style={concept color=popblue, font=\bfseries},
level 1/.append style={level distance=3.6cm, sibling angle=90},
level 1 concept/.append style={concept color=poporange, font=\small\bfseries},
level 2/.append style={level distance=2.6cm},
level 2 concept/.append style={concept color=popgreenL, font=\footnotesize}]
\node[root concept]{L'État camerounais}
  child { node {Définition et éléments}
      child { node {Territoire} }
      child { node {Population} }
      child { node {Pouvoir souverain} } }
  child { node {Fondements}
      child { node {Force ?} }
      child { node {Contrat ?} }
      child { node {Légitimité ?} } }
  child { node {Rôles}
      child { node {Sécurité (régalien)} }
      child { node {Justice (régalien)} }
      child { node {Éducation, santé (social)} }
      child { node {Infrastructures (économique)} } }
  child { node {Finalité}
      child { node {Bien commun} }
      child { node {Droits de l'homme} }
      child { node {Intérêt général} } };
\end{tikzpicture}
\end{center}
\legende{Modèle de carte mentale vierge à compléter : le nœud central et les quatre branches principales sont donnés ; chaque groupe ajoute les sous-branches, les mots-clés et les exemples camerounais tirés des documents.}
\end{popfigure}

\begin{attention}[Erreur fréquente]
Recopier les documents sans les analyser. Une carte mentale n'est pas un résumé : chaque extrait de la Constitution ou chaque photo doit \emph{servir à justifier une réponse} à la question posée (« qu'est-ce que l'État ? », « pourquoi obéit-on ? »). Un document non relié à une idée est un document perdu.
\end{attention}

\courssub{TP 2 — Tableau comparatif : Hobbes, Locke, Rousseau face aux situations camerounaises}

\textbf{Objectif.} Croiser les trois grandes conceptions du fondement de l'État avec des situations concrètes de la vie politique camerounaise, afin de comprendre que les théories ne sont pas des idées abstraites mais des grilles de lecture du réel.

\textbf{Protocole.} Compléter le tableau suivant en groupes : pour chaque situation, indiquer quel auteur permet de l'expliquer et pourquoi.

\begin{center}
\begin{tabularx}{\linewidth}{|l|X|X|X|}
\hline
\rowcolor{popblueL} \textbf{Situation camerounaise} & \textbf{Hobbes (sécurité)} & \textbf{Locke (droits)} & \textbf{Rousseau (volonté générale)} \\
\hline
Opération de maintien de l'ordre lors de troubles & L'État garantit la paix contre la guerre de tous & L'ordre ne doit pas violer les droits naturels & L'ordre est légitime s'il exprime la volonté générale \\
\hline
Organisation d'élections & Le souverain tranche les conflits pour la stabilité & Le consentement fonde le pouvoir (mandat) & Le peuple souverain choisit ses représentants \\
\hline
Débat sur une révision constitutionnelle & \ldots & \ldots & \ldots \\
\hline
\end{tabularx}
\end{center}

La dernière ligne est laissée à compléter : c'est l'exercice le plus formateur, car un débat constitutionnel peut être lu à la fois comme un exercice de souveraineté (Hobbes), comme une discussion sur les limites du pouvoir (Locke) et comme une expression de la volonté générale (Rousseau). Le groupe doit choisir une lecture dominante et la \emph{justifier}.

\begin{exemplebox}[Exemple de justification attendue]
« Nous lisons les élections avec Rousseau : le vote est le moment où le peuple exprime la volonté générale. Mais nous ajoutons Locke : si les élections ne sont pas libres, le consentement disparaît et le pouvoir perd sa légitimité. » Ce type de réponse croise les auteurs au lieu de les juxtaposer : c'est cela, analyser.
\end{exemplebox}

\courssub{TP 3 — Diagramme statistique : lire les priorités de l'État dans son budget}

\textbf{Objectif.} Passer du discours aux chiffres : un budget est une décision politique chiffrée. Ce que l'État finance le plus révèle ce qu'il juge le plus important.

\textbf{Données fournies (simplifiées, fictives, à titre pédagogique uniquement).} Répartition du budget de l'État en pourcentage : Éducation 15 \%, Défense et sécurité 12 \%, Santé 8 \%, Infrastructures 20 \%, Administration 25 \%, Dette et autres 20 \%.

\textbf{Protocole.} Reporter ces valeurs dans le gabarit ci-dessous (ou construire un diagramme circulaire), puis formuler une conclusion argumentée.

\begin{popfigure}
\begin{center}
\begin{tikzpicture}
\begin{axis}[
  ybar, bar width=0.55cm,
  width=0.95\linewidth, height=6cm,
  ymin=0, ymax=30,
  ylabel={Part du budget (en \%)},
  symbolic x coords={Éducation, Défense, Santé, Infrastructures, Administration, Dette},
  xtick=data,
  x tick label style={rotate=25, anchor=east, font=\small},
  nodes near coords, nodes near coords style={font=\small},
  axis lines=left,
  ymajorgrids=true, grid style={popline!50}
]
\addplot[fill=popblue, draw=popdark] coordinates {(Éducation,15) (Défense,12) (Santé,8) (Infrastructures,20) (Administration,25) (Dette,20)};
\end{axis}
\end{tikzpicture}
\end{center}
\legende{Gabarit de diagramme en barres à remplir avec les données budgétaires fournies en classe (données simplifiées et fictives, à titre d'entraînement). Source : simulation pédagogique.}
\end{popfigure}

\begin{methode}[Lire un diagramme statistique]
\begin{enumerate}
\item Lire le \textbf{titre} : que mesure le diagramme ?
\item Lire les \textbf{axes} (diagramme en barres) ou les \textbf{parts} (diagramme circulaire) : quelles grandeurs, quelles unités ?
\item Repérer la \textbf{valeur maximale} et la \textbf{valeur minimale} : ici, l'administration (25 \%) domine, la santé (8 \%) est la plus faible.
\item Formuler une \textbf{conclusion en une phrase} : « La part la plus importante du budget est consacrée à \ldots, ce qui montre que la priorité de l'État est \ldots ».
\end{enumerate}
\end{methode}

\begin{exemplebox}[Discussion chiffrée]
Si l'éducation représente environ 15 \% du budget et la défense 12 \%, on peut discuter : ces choix traduisent-ils une finalité de \cle{bien commun} ? Un élève peut soutenir que oui (instruire et protéger servent tous les citoyens), un autre que la faiblesse de la santé (8 \%) montre un écart entre la finalité proclamée et les moyens réels. Les deux réponses sont recevables \emph{si elles s'appuient sur les chiffres}.
\end{exemplebox}

\courssub{TP 4 — Croquis d'analyse : la séparation des pouvoirs au Cameroun}

\textbf{Objectif.} Schématiser les trois pouvoirs et identifier les organes camerounais qui les incarnent, afin de vérifier comment la théorie de la séparation des pouvoirs (Montesquieu, au service de la liberté) s'incarne dans les institutions réelles.

\textbf{Protocole.} Reproduire le schéma ci-dessous et compléter les cases vides avec les organes étudiés : la présidence de la République et le gouvernement (exécutif), l'Assemblée nationale et le Sénat (législatif), la Cour suprême et les juridictions (judiciaire). Flécher ensuite les relations de contrôle entre pouvoirs (ex. : promulgation des lois, nomination aux hautes fonctions, contrôle de constitutionnalité, responsabilité du gouvernement).

\begin{popfigure}
\begin{center}
\begin{tikzpicture}[node distance=1.2cm,
box/.style={draw=popdark, fill=popblueL, rounded corners, minimum width=3.4cm, minimum height=1.2cm, align=center, font=\small\bfseries},
ctrl/.style={-{Stealth[length=3mm]}, thick, poporange}]
\node[box, align=center] (exec){Pouvoir exécutif\\ \footnotesize (organes : \ldots)};
\node[box, right=2.2cm of exec, align=center] (leg){Pouvoir législatif\\ \footnotesize (organes : \ldots)};
\node[box, right=2.2cm of leg, align=center] (jud){Pouvoir judiciaire\\ \footnotesize (organes : \ldots)};
\draw[ctrl] (exec) to[bend left=20] node[above, font=\footnotesize, align=center]{promulgation\\ des lois} (leg);
\draw[ctrl] (leg) to[bend left=20] node[above, font=\footnotesize, align=center]{vote de la loi\\ contrôle action} (exec);
\draw[ctrl] (jud) to[bend right=25] node[below, font=\footnotesize, align=center]{contrôle de\\ constitutionnalité} (exec);
\draw[ctrl] (jud) to[bend left=25] node[below, font=\footnotesize, align=center]{contrôle de\\ constitutionnalité} (leg);
\end{tikzpicture}
\end{center}
\legende{Croquis à compléter : la séparation des pouvoirs au Cameroun. Le groupe nomme les organes dans chaque case et discute les flèches de contrôle mutuel.}
\end{popfigure}

La question de fond à discuter en groupe : les trois pouvoirs sont-ils réellement \emph{équilibrés} au Cameroun, ou l'exécutif domine-t-il ? La réponse doit mobiliser le vocabulaire de la leçon : \cle{séparation des pouvoirs}, \cle{souveraineté}, \cle{légitimité}.

\begin{savaistu}[Un indice textuel de la finalité de l'État]
La Constitution camerounaise de 1996 (loi n° 96-06 du 18 janvier 1996) proclame, dans son préambule, l'attachement du peuple camerounais aux droits fondamentaux de l'homme tels qu'énoncés dans la Déclaration universelle des droits de l'homme de 1948 et la Charte africaine des droits de l'homme et des peuples de 1981. Ce texte n'est pas un détail : il constitue un \emph{indice textuel} de la finalité de l'État camerounais, qui se donne officiellement pour but la protection de la dignité et des libertés des citoyens. Dans une copie, citer ce préambule est une preuve solide.
\end{savaistu}

\courssub{Dossier guidé : étude d'un cas concret}

\textbf{Objectif.} Appliquer les notions de l'unité à une situation de la vie courante et justifier la démarche, conformément aux savoir-faire du programme.

\begin{exoresolu}[La construction d'une route nationale : quel rôle de l'État ?]
Énoncé. L'État camerounais décide de construire une route nationale reliant une zone de production agricole à un grand marché urbain. Identifier le rôle de l'État mobilisé, puis discuter la finalité de cette action en mobilisant au moins un auteur de la leçon.
\tcblower
Solution détaillée.
\textbf{Étape 1 — Identifier le rôle.} Construire une route relève du rôle \emph{économique et social} de l'État : il fournit un bien collectif (une infrastructure) qu'aucun particulier ne financerait seul, et il organise la vie de la cité.
\textbf{Étape 2 — Identifier le fondement.} Cette action est légitime parce que l'État est souverain sur son territoire et agit au nom de la population ; elle suppose l'impôt, c'est-à-dire la contribution de tous.
\textbf{Étape 3 — Discuter la finalité.} Avec Rousseau, on dira que la route sert la volonté générale si elle profite à tous (transport des récoltes, accès aux soins, désenclavement). Avec Locke, on vérifiera qu'elle protège la propriété et la liberté de circulation sans exproprier injustement. Avec Hobbes, on ajoutera qu'en reliant les régions, elle renforce l'unité et la paix du corps politique.
\textbf{Conclusion.} La construction d'une route nationale illustre le rôle de service public de l'État et sa finalité de bien commun — à condition que l'ouvrage serve réellement la population et non quelques intérêts particuliers.
\end{exoresolu}

\begin{exercice}[À vous de jouer]
Choisir l'un des cas suivants : la création d'une université d'État dans une région, ou une opération de maintien de l'ordre lors d'une manifestation. En une quinzaine de lignes : (1) identifier le rôle de l'État mobilisé ; (2) nommer le fondement de son action ; (3) discuter sa finalité en citant un auteur. Rédiger et justifier la conclusion.
\end{exercice}

\begin{corrige}[Exercice : indications de correction]
Pour l'université : rôle éducatif et social (service public de l'enseignement supérieur) ; fondement : la souveraineté et le consentement (l'impôt finance le service) ; finalité : l'égalité des chances et le bien commun (Rousseau, Locke). Pour le maintien de l'ordre : rôle régalien de sécurité ; fondement : chez Hobbes, la protection contre la guerre de tous ; finalité : la paix, mais à discuter avec Locke — la force publique est légitime seulement si elle respecte les droits des citoyens. Toute copie qui identifie correctement rôle, fondement et finalité, et qui justifie par un auteur ou un document, est valorisée.
\end{corrige}

\courssub{Restitution : présenter et justifier oralement}

Chaque groupe présente son support (carte mentale, tableau comparatif, diagramme, croquis) en trois minutes. La présentation doit suivre une règle simple : \emph{montrer, expliquer, justifier}. Montrer le support, expliquer les choix de classement, justifier chaque choix par le vocabulaire de la leçon (\cle{État}, \cle{pouvoir}, \cle{souveraineté}, \cle{légitimité}, \cle{bien commun}, \cle{séparation des pouvoirs}).

\begin{aretenir}[Ce que ce TP doit vous laisser]
Analyser l'État, c'est croiser quatre regards : ce qu'il \emph{est} (définition et éléments), ce qui le \emph{fonde} (force, contrat, volonté générale), ce qu'il \emph{fait} (rôles régaliens, économiques, sociaux) et ce pour quoi il \emph{existe} (finalité du bien commun). Un document, un chiffre ou un schéma ne vaut que s'il justifie une réponse argumentée : c'est cette exigence qui transforme un travail de groupe en raisonnement philosophique.
\end{aretenir}

\coursec{Méthode et entraînement : rédiger et justifier une réponse philosophique}

Dans les sections précédentes, nous avons construit ensemble des outils d'analyse de l'État camerounais : fiches de notions, repères historiques, exemples concrets tirés de notre vie quotidienne (la carte nationale d'identité, les impôts, les tribunaux, les élections). Mais posséder des connaissances ne suffit pas : à l'examen du Probatoire (Première technique) comme au Baccalauréat technique (Terminale), l'épreuve de philosophie demande de \cle{rédiger} et de \cle{justifier} une réponse. Autrement dit, il faut transformer ce que l'on sait en un raisonnement écrit, structuré et convaincant. Cette dernière section est donc un entraînement méthodique : nous allons apprendre à construire une dissertation et une explication de texte sur le thème de l'État et du pouvoir, pas à pas, comme un artisan apprend son geste.

\courssub{La méthode de la dissertation philosophique}

\begin{definition}[La dissertation philosophique]
La \cle{dissertation} est un exercice écrit qui consiste à répondre à une question philosophique par un raisonnement rigoureux, organisé en trois moments : une \cle{introduction} qui pose le problème, un \cle{développement} qui examine les réponses possibles, et une \cle{conclusion} qui tranche et ouvre. Elle ne rapporte pas des connaissances brutes : elle résout un problème.
\end{definition}

\textbf{L'introduction : quatre étapes obligatoires.} Une bonne introduction comporte, dans l'ordre :

\begin{enumerate}
\item \textbf{L'accroche} : une phrase qui situe le sujet dans l'expérience commune. Elle part du concret pour montrer que la question concerne tout le monde. Exemple : « Chaque matin, en passant devant le commissariat ou en payant une taxe au marché, le citoyen camerounais rencontre l'État sans même y penser. »
\item \textbf{Le problème} : on dégage la difficulté cachée dans le sujet. Pourquoi la question se pose-t-elle ? Quelle tension contient-elle ?
\item \textbf{La problématique} : la question précise que le devoir va traiter, formulée en une phrase interrogative.
\item \textbf{L'annonce du plan} : on indique brièvement les étapes du raisonnement (« Nous verrons d'abord que..., puis que..., enfin que... »).
\end{enumerate}

\begin{methode}[Construire une bonne problématique]
Une problématique n'est pas une question quelconque : c'est une question \textbf{précise} qui naît d'une \textbf{tension} entre deux idées. Procédez en trois étapes :
\begin{enumerate}
\item Repérez dans le sujet deux termes qui semblent s'opposer (par exemple : \emph{contrainte} et \emph{protection}).
\item Formulez la tension : « l'État nous contraint ; or il prétend nous protéger ».
\item Transformez la tension en question : « comment la contrainte de l'État peut-elle servir la liberté des citoyens ? »
\end{enumerate}
Une problématique réussie est celle dont on ne peut pas répondre simplement par « oui » ou par « non ».
\end{methode}

\textbf{Le développement : trois parties, trois idées.} Au niveau de la Première technique, le plan le plus sûr est le plan dialectique en trois parties :

\begin{itemize}
\item \textbf{Première partie (thèse)} : on défend une première réponse, la plus spontanée ou la plus classique.
\item \textbf{Deuxième partie (antithèse)} : on montre les limites de cette réponse et on défend la réponse opposée.
\item \textbf{Troisième partie (synthèse)} : on dépasse l'opposition en proposant une position nuancée et argumentée.
\end{itemize}

\begin{methode}[Construire un paragraphe de développement]
Chaque paragraphe du développement suit la même recette, en quatre ingrédients :
\begin{enumerate}
\item \textbf{Une idée} clairement énoncée dans la première phrase ;
\item \textbf{Une référence philosophique} expliquée (ne pas se contenter de nommer l'auteur : dire ce qu'il affirme et pourquoi) ;
\item \textbf{Un exemple concret}, camerounais si possible, qui illustre l'idée ;
\item \textbf{Une transition} qui prépare le paragraphe suivant.
\end{enumerate}
Un paragraphe sans référence est une opinion ; une référence sans exemple est de l'érudition vide. Il faut les deux.
\end{methode}

\textbf{La conclusion.} Elle répond explicitement à la problématique, résume le chemin parcouru en deux ou trois phrases, et peut s'achever par une ouverture (une nouvelle question qui prolonge la réflexion, sans la traiter).

\begin{popfigure}
\begin{center}
\begin{tikzpicture}[
box/.style={draw=popblue, fill=popblueL, rounded corners=2pt, text width=3.1cm, minimum height=0.9cm, align=center, font=\small},
fleche/.style={-{Stealth[length=2.5mm]}, thick, popdark}]
\node[box, align=center] (a) at (0,4.5){\textbf{Accroche}\\ (expérience commune)};
\node[box, fill=poporangeL, draw=poporange, align=center] (b) at (3.8,3.4){\textbf{Problème}\\ (la tension)};
\node[box, fill=poppurpleL, draw=poppurple, align=center] (c) at (7.6,2.3){\textbf{Problématique}\\ (la question)};
\node[box, fill=popgreenL, draw=popgreen, align=center] (d) at (3.8,1.2){\textbf{Plan}\\ thèse / antithèse / synthèse};
\node[box, fill=popgoldL, draw=popgold, align=center] (e) at (0,0.1){\textbf{Conclusion}\\ (réponse + ouverture)};
\draw[fleche] (a) -- (b);
\draw[fleche] (b) -- (c);
\draw[fleche] (c) -- (d);
\draw[fleche] (d) -- (e);
\end{tikzpicture}
\end{center}
\legende{Schéma en escalier de la dissertation philosophique : chaque étape descend logiquement de la précédente, de l'accroche jusqu'à la conclusion.}
\end{popfigure}

\begin{popfigure}
\begin{center}
\begin{tikzpicture}
\node[draw=popdark, fill=popblueL, font=\bfseries\small, minimum width=3.4cm, minimum height=0.7cm] (t1) at (0,3.2) {Annoncer une étape};
\node[draw=popdark, fill=poporangeL, font=\bfseries\small, minimum width=3.4cm, minimum height=0.7cm] (t2) at (4.2,3.2) {Expliquer, justifier};
\node[draw=popdark, fill=poppurpleL, font=\bfseries\small, minimum width=3.4cm, minimum height=0.7cm] (t3) at (8.4,3.2) {Opposer, nuancer};
\node[draw=popdark, fill=popgreenL, font=\bfseries\small, minimum width=3.4cm, minimum height=0.7cm] (t4) at (2.1,0) {Conclure une étape};
\node[draw=popdark, fill=popgoldL, font=\bfseries\small, minimum width=3.4cm, minimum height=0.7cm] (t5) at (6.3,0) {Conclure le devoir};
\node[below=2pt of t1, text width=3.4cm, align=center, font=\small] {d'abord, premièrement, en premier lieu, ensuite};
\node[below=2pt of t2, text width=3.4cm, align=center, font=\small] {en effet, c'est-à-dire, autrement dit, car, puisque};
\node[below=2pt of t3, text width=3.4cm, align=center, font=\small] {cependant, néanmoins, mais, or, en revanche};
\node[below=2pt of t4, text width=3.4cm, align=center, font=\small] {par conséquent, ainsi, donc, c'est pourquoi};
\node[below=2pt of t5, text width=3.4cm, align=center, font=\small] {en définitive, en somme, finalement, pour conclure};
\end{tikzpicture}
\end{center}
\legende{Les connecteurs logiques classés par fonction : ils sont les articulations du raisonnement et guident le correcteur dans votre devoir.}
\end{popfigure}

\courssub{Application à un sujet type : « L'État est-il un mal nécessaire ? »}

Appliquons immédiatement la méthode. Le sujet contient une tension : l'expression « mal nécessaire » suppose que l'État fait du mal (il contraint, il punit, il prélève l'impôt), mais qu'on ne peut pas s'en passer. D'où la problématique possible : \emph{la contrainte étatique est-elle le prix à payer pour la sécurité et la liberté, ou un obstacle à supprimer ?}

\begin{exemplebox}[Plan possible pour « L'État est-il un mal nécessaire ? »]
\begin{itemize}
\item \textbf{I. L'État comme mal nécessaire (thèse hobbesienne).} L'État contraint, punit, prélève l'impôt : c'est un « mal » (restriction de la liberté naturelle). Mais sans lui, c'est l'état de nature, la guerre de tous contre tous (Hobbes), qui est un mal bien plus grand. L'État est le moindre mal, nécessaire à la sécurité. Exemple : la police qui réprime un vol au marché Mokolo (mal pour le voleur, nécessaire pour la paix des commerçants).
\item \textbf{II. L'État comme institution positive (antithèse lockéenne et rousseauiste).} L'État n'est pas un mal qu'on subit, mais une institution voulue par les hommes pour protéger leurs droits naturels (Locke : vie, liberté, propriété) ou réaliser la liberté civile (Rousseau : passage de la liberté naturelle à la liberté morale). La contrainte n'est pas un mal, c'est la forme de la loi qu'on se donne. Exemple : l'impôt payé pour financer l'école publique ou l'hôpital de district (service du bien commun).
\item \textbf{III. Synthèse : la nature de l'État dépend de sa légitimité (État de droit).} L'État n'est un « mal nécessaire » que s'il est arbitraire ou tyrannique. S'il est un État de droit (séparation des pouvoirs, respect des droits fondamentaux, consentement des gouvernés), la contrainte devient service de la liberté. Le « mal » disparaît dans la légitimité démocratique.
\end{itemize}
\end{exemplebox}

\courssub{La méthode de l'explication de texte}

L'explication de texte est le deuxième exercice possible à l'examen. Elle consiste à commenter un court passage d'un philosophe. Sa méthode comporte quatre gestes :

\begin{enumerate}
\item \textbf{Dégager la thèse} : quelle est l'idée principale défendue par l'auteur ? On la formule en une phrase, avec ses propres mots.
\item \textbf{Suivre l'argumentation} : comment l'auteur prouve-t-il sa thèse ? On repère les étapes du raisonnement, les exemples, les distinctions.
\item \textbf{Expliquer les concepts} : on définit les termes importants du texte (par exemple « État de nature », « volonté générale », « contrat »).
\item \textbf{Discuter la thèse} : en conclusion, on peut approuver, nuancer ou contester la position de l'auteur, à condition d'argumenter.
\end{enumerate}

\begin{attention}[Opinion ou argument ?]
L'erreur la plus fréquente est de confondre \cle{opinion} et \cle{argument}. Une opinion s'\emph{affirme} : « je pense que l'État est trop présent ». Un argument se \emph{justifie} par une raison et une référence : « une intervention excessive de l'État peut menacer les libertés individuelles, car, comme le montre Tocqueville, le pouvoir centralisé tend à étouffer les initiatives des citoyens ». Au Probatoire technique, seul l'argument est noté.
\end{attention}

\courssub{Entraînement guidé : rédiger une introduction}

Travail collectif proposé : rédiger l'introduction d'une dissertation sur le sujet « Pourquoi obéit-on à l'État ? ». Suivons les quatre étapes.

\begin{exoresolu}[Introduction rédigée : « Pourquoi obéit-on à l'État ? »]
Rédiger l'introduction complète de ce sujet, avec accroche, problème, problématique et annonce du plan.
\tcblower
\textbf{Accroche} : Chaque jour, le citoyen camerounais obéit à des règles qu'il n'a pas choisies : il paie ses impôts, respecte le code de la route, présente sa carte nationale d'identité lors d'un contrôle. Cette obéissance semble aller de soi, au point qu'on ne la remarque plus.
\textbf{Problème} : Pourtant, obéir, c'est renoncer à une part de sa liberté. Pourquoi un être libre accepte-t-il de se soumettre à un pouvoir ? Est-ce par peur de la sanction, par habitude, ou parce qu'il reconnaît à l'État une légitimité ?
\textbf{Problématique} : L'obéissance à l'État repose-t-elle sur la contrainte ou sur le consentement raisonné des citoyens ?
\textbf{Annonce du plan} : Nous verrons d'abord que l'obéissance peut s'expliquer par la force et la crainte ; nous montrerons ensuite qu'elle se fonde plus profondément sur un consentement rationnel ; nous soutiendrons enfin qu'une obéissance légitime suppose un État de droit et un citoyen vigilant.
\end{exoresolu}

\begin{exercice}[À vous de jouer]
Rédigez l'introduction complète du sujet « L'État sert-il la liberté des citoyens ? », puis construisez le plan en trois parties. Vous utiliserez au moins deux références étudiées dans la leçon (Hobbes, Locke, Rousseau) et un exemple camerounais.
\end{exercice}

\begin{corrige}[Exercice 1]
\textbf{Introduction possible.} Accroche : lors d'une élection ou d'un contrôle d'identité à Yaoundé, le citoyen découvre que l'État organise sa vie. Problème : l'État impose des lois qui limitent nos actions ; or il prétend aussi garantir nos droits. La contrainte et la liberté semblent s'opposer. Problématique : la contrainte étatique est-elle un obstacle à la liberté ou la condition de son exercice ? Annonce : nous verrons d'abord que l'État limite la liberté, puis qu'il la rend possible, enfin qu'il ne la sert que s'il est un État de droit.
\textbf{Plan possible.} I. L'État contraint : lois, impôts, sanctions (Hobbes : la peur de la punition maintient l'ordre). II. Mais sans État, pas de liberté réelle : dans l'état de nature, le plus fort écrase le plus faible ; la loi protège (Locke : l'État garantit les droits naturels ; exemple : le tribunal qui protège le petit commerçant contre l'expropriation abusive). III. Synthèse : l'État ne sert la liberté que s'il respecte le droit et la séparation des pouvoirs (Montesquieu) ; un pouvoir arbitraire la détruit.
\end{corrige}

\courssub{Les critères d'évaluation au Probatoire technique}

Pour réussir, il faut savoir comment le correcteur lit votre copie. Quatre critères principaux guident la notation :

\begin{center}
\begin{tabularx}{\linewidth}{|l|X|}
\hline
\rowcolor{popblueL} \textbf{Critère} & \textbf{Ce que le correcteur attend} \\
\hline
Exactitude des notions & Les termes (État, souveraineté, légitimité, contrat) sont définis correctement et employés à bon escient. \\
\hline
Pertinence des références & Les auteurs cités (Hobbes, Rousseau, Locke...) correspondent réellement à l'idée défendue ; leurs thèses sont expliquées, pas seulement nommées. \\
\hline
Qualité de l'argumentation & Chaque idée est justifiée par une raison, illustrée par un exemple, reliée à la suivante par des connecteurs logiques. \\
\hline
Clarté de l'expression & Phrases simples, orthographe soignée, plan visible (alinéas, transitions). \\
\hline
\end{tabularx}
\end{center}

\begin{attention}[Les trois erreurs qui coûtent le plus de points]
\begin{enumerate}
\item \textbf{Paraphraser le sujet} : recopier la question en la reformulant pendant toute la copie, sans jamais y répondre. Le sujet est un point de départ, pas un contenu.
\item \textbf{Citer sans expliquer} : écrire « comme le dit Rousseau » sans dire ce que Rousseau affirme ni en quoi cela sert votre raisonnement. Une référence non expliquée vaut zéro.
\item \textbf{Donner une opinion personnelle non argumentée} : « à mon avis, l'État est corrompu » n'est pas une thèse philosophique. Toute position doit être soutenue par une raison et une référence.
\end{enumerate}
\end{attention}

\begin{savaistu}[La justification de la démarche est valorisée]
Au Probatoire technique, l'épreuve de philosophie valorise explicitement la \emph{justification de la démarche} : un plan clair, annoncé et suivi, accompagné de références exactes, rapporte autant de points que l'érudition. Un candidat qui connaît peu d'auteurs mais raisonne rigoureusement obtient souvent une meilleure note qu'un candidat qui accumule les citations sans les organiser. La philosophie est un art du raisonnement avant d'être un art de la mémoire.
\end{savaistu}

\begin{aretenir}[L'essentiel de la méthode]
\begin{itemize}
\item Introduction : accroche, problème, problématique (une question née d'une tension), annonce du plan.
\item Développement : trois parties (thèse, antithèse, synthèse) ; chaque paragraphe = une idée + une référence expliquée + un exemple concret + une transition.
\item Conclusion : réponse explicite à la problématique, bilan, ouverture.
\item Explication de texte : dégager la thèse, suivre l'argumentation, expliquer les concepts, discuter.
\item Un argument se justifie ; une opinion s'affirme. Seul l'argument est noté.
\end{itemize}
\end{aretenir}

Vous disposez désormais de tous les outils : les notions de la leçon (définition, fondements, rôles et finalité de l'État) et la méthode pour les mobiliser. La philosophie, comme un métier technique, s'apprend par la pratique : c'est en rédigeant, en corrigeant et en rédigeant encore que le raisonnement devient un réflexe. À vos copies !

\newpage

\coursec{Activité d'intégration}

\courssub{Objectifs de l'activité}

Cette activité d'intégration clôt la leçon « L'État et le pouvoir ». Elle vise à faire travailler ensemble l'ensemble des savoirs étudiés : la \cle{définition} de l'État et ses trois éléments constitutifs (population, territoire, pouvoir souverain), les \cle{fondements} du pouvoir (la force, la tradition, l'accord ou consentement, le droit), les \cle{rôles} de l'État (sécurité, justice, éducation, santé, infrastructures, redistribution) et sa \cle{finalité} (l'ordre, la justice, le bien commun).

À l'issue de l'activité, l'élève doit être capable de :
\begin{itemize}
\item appliquer une grille d'analyse philosophique à une situation concrète de la vie courante (savoir-faire 1) ;
\item distinguer un pouvoir de fait d'un pouvoir légitime ;
\item rédiger et justifier une réponse argumentée, en mobilisant au moins deux auteurs du programme (savoir-faire 2) ;
\item défendre oralement une position et écouter celle des autres dans un débat réglé.
\end{itemize}

\textbf{Organisation matérielle :} classe répartie en groupes de 4 à 6 élèves ; durée : 2 heures (1 h de travail de groupe, 30 min de restitution, 30 min de débat et synthèse) ; support : le dossier documentaire ci-dessous, reproduit pour chaque groupe.

\courssub{Situation}

\textbf{Dossier documentaire : le quartier « Montée des Aigles », périphérie de Yaoundé.}

Au-delà du quartier Nkolbisson, à l'ouest de Yaoundé, s'est développé depuis 2018 un nouveau quartier périurbain appelé « Montée des Aigles ». Né de lotissements informels, il n'est rattaché à aucune commune et ne figure sur aucun plan administratif officiel. Voici les données recueillies par une équipe d'enquête :

\begin{center}
\begin{tabularx}{\linewidth}{|l|X|}\hline
\rowcolor{popblueL} \textbf{Rubrique} & \textbf{Données de l'enquête} \\ \hline
Population & Environ \num{12000} habitants, soit environ \num{2100} ménages \\ \hline
Territoire & Environ 3 km$^2$, sans limites officielles reconnues \\ \hline
Sécurité & Aucun commissariat ni brigade de gendarmerie ; le plus proche est à 14 km \\ \hline
Éducation & Aucune école publique ; deux écoles privées coûteuses ; 62 \% des enfants de 6 à 11 ans scolarisés \\ \hline
Santé & Aucun centre de santé ; l'hôpital le plus proche est à 17 km \\ \hline
Voirie et eau & Pistes non entretenues, impraticables en saison des pluies ; forages privés \\ \hline
Conflits & En 2024 : 38 litiges fonciers déclarés, 12 vols signalés, 3 affaires de violence \\ \hline
\end{tabularx}
\end{center}

Face à ce vide, les habitants ont créé en 2020 un « \textbf{Comité de quartier} » dirigé par un président élu à main levée lors d'une assemblée générale, assisté de six « chefs de bloc » désignés par les anciens du quartier. Le comité :
\begin{itemize}
\item fixe des règles écrites (interdiction de jeter les ordures sur la piste, couvre-feu informel à 22 h, obligation de participer aux travaux communautaires du samedi) ;
\item lève une « taxe de développement » de 500 FCFA par mois et par ménage, qui rapporte environ $2100 \times 500 = \num{1050000}$ FCFA par mois, soit environ 12,6 millions de FCFA par an, utilisés pour réparer la piste principale et payer deux veilleurs de nuit ;
\item sanctionne les récalcitrants : amende de \num{2500} FCFA, puis exclusion des travaux communautaires, puis, dans deux cas en 2024, expulsion forcée du quartier par un groupe de jeunes ;
\item arbitre les petits conflits de voisinage, mais renvoie les litiges fonciers graves « devant les chefs traditionnels » du village voisin.
\end{itemize}

Certains habitants se félicitent : « Ici, notre comité fait ce que l'État ne fait pas. Nous n'avons pas besoin de l'État. » D'autres protestent : « Le comité n'a aucun droit de lever une taxe ni d'expulser des familles. Seul l'État peut faire cela. »

\textbf{Problème philosophique posé à la classe :} le comité de quartier de Montée des Aigles est-il un État ? Peut-il durablement remplacer l'État ?

\courssub{Consignes}

Chaque groupe traite les quatre tâches suivantes, dans l'ordre, et consigne ses réponses par écrit.

\begin{enumerate}
\item \textbf{Le comité est-il un État ?} À l'aide de la grille des trois éléments constitutifs de l'État (population, territoire, pouvoir souverain), complétez un tableau à deux colonnes (« élément constitutif » / « le comité le possède-t-il ? justifiez »). Concluez.
\item \textbf{Sur quels fondements repose l'autorité du comité ?} Relevez dans le dossier au moins un indice pour chacun des quatre fondements du pouvoir étudiés (la force, la tradition, l'accord ou consentement, le droit). Lequel domine ? Ce fondement est-il suffisant pour rendre le pouvoir légitime ?
\item \textbf{Quels rôles l'État devrait-il assurer à Montée des Aigles ?} Construisez un tableau à trois colonnes : « rôle de l'État », « situation actuelle dans le quartier », « ce que l'État camerounais devrait faire ». Traitez au minimum : sécurité, justice, éducation, santé, voirie.
\item \textbf{Rédaction argumentée (demi-page, environ 20 lignes) :} « Le comité de quartier peut-il durablement remplacer l'État ? » Votre réponse doit : (a) prendre clairement position ; (b) mobiliser au moins deux philosophes étudiés dans la leçon (par exemple Hobbes, Locke, Rousseau) ; (c) s'appuyer sur les données chiffrées du dossier ; (d) se terminer par une conclusion justifiée.
\end{enumerate}

\textbf{Restitution :} chaque groupe dispose de 4 minutes pour présenter ses résultats. Le débat collectif porte sur la question 4 : les groupes en désaccord doivent répondre aux arguments adverses, pas seulement répéter les leurs.

\courssub{Ressources à mobiliser}

\begin{aretenir}[Boîte à outils de la leçon]
\begin{itemize}
\item \textbf{Définition :} l'État est une organisation politique qui réunit trois éléments : une \cle{population}, un \cle{territoire} délimité, et un \cle{pouvoir souverain} exerçant l'autorité suprême, reconnue comme légitime, sur ce territoire et cette population.
\item \textbf{Fondements du pouvoir :} la \cle{force} (obéir par crainte), la \cle{tradition} (obéir par habitude et héritage), l'\cle{accord} ou consentement (obéir parce qu'on a choisi ses gouvernants), le \cle{droit} (obéir à des règles légales, votées et applicables à tous).
\item \textbf{Rôles de l'État :} assurer la sécurité des personnes et des biens, rendre la justice, organiser l'éducation et la santé, construire et entretenir les infrastructures, prélever l'impôt pour financer le bien commun.
\item \textbf{Finalité de l'État :} garantir l'ordre et la paix (Hobbes), protéger les droits naturels des individus (Locke), réaliser la volonté générale et le bien commun (Rousseau).
\end{itemize}
\end{aretenir}

\begin{attention}[Pièges à éviter]
\begin{itemize}
\item Ne confondez pas \textbf{pouvoir de fait} (quelqu'un commande effectivement) et \textbf{pouvoir légitime} (on a de bonnes raisons d'obéir). Le comité commande : cela ne suffit pas à en faire un État.
\item La population et le territoire ne suffisent pas : c'est la \textbf{souveraineté reconnue} qui fait la différence. Un stade plein réunit une population sur un territoire, ce n'est pas un État.
\item Dans la rédaction, citez les philosophes pour \emph{appuyer} votre raisonnement, pas pour le remplacer : une citation sans explication ne vaut rien.
\end{itemize}
\end{attention}

\courssub{Démarche et corrigé commenté}

\begin{exoresolu}[Corrigé commenté de l'enquête de Montée des Aigles]
\textbf{Étape 1 — Le comité est-il un État ?} On applique la grille des trois éléments constitutifs.

\begin{center}
\begin{tabularx}{\linewidth}{|l|X|}\hline
\rowcolor{popblueL} \textbf{Élément constitutif} & \textbf{Le comité le possède-t-il ?} \\ \hline
Population & Partiellement : \num{12000} habitants identifiés, mais aucun statut juridique (pas d'état civil propre, pas de nationalité). \\ \hline
Territoire & Non : 3 km$^2$ « sans limites officielles reconnues » ; le sol reste soumis à la souveraineté de l'État camerounais. \\ \hline
Pouvoir souverain & Non : le comité exerce un pouvoir de fait (règles, taxe, sanctions), mais il n'est pas souverain : il ne crée pas le droit, il n'est reconnu par aucune loi, et il renvoie lui-même les litiges graves aux chefs traditionnels ou à l'administration. \\ \hline
\end{tabularx}
\end{center}

\textbf{Conclusion de l'étape 1 :} le comité n'est pas un État. C'est une \emph{association de fait} qui exerce localement quelques fonctions ressemblant à celles de l'État, sans en avoir la souveraineté ni la légalité.

\tcblower

\textbf{Étape 2 — Les fondements de l'autorité du comité.} On relève les indices dans le dossier :
\begin{itemize}
\item \textbf{Accord/consentement :} le président est « élu à main levée lors d'une assemblée générale ». C'est le fondement le plus fort : les habitants ont consenti à l'autorité du comité. On pense à Locke : le pouvoir légitime repose sur le consentement des gouvernés.
\item \textbf{Tradition :} les « chefs de bloc » sont « désignés par les anciens du quartier », et les litiges fonciers sont renvoyés aux chefs traditionnels. L'autorité s'appuie en partie sur la coutume.
\item \textbf{Force :} les expulsions forcées « par un groupe de jeunes » relèvent de la pure contrainte. Or un pouvoir fondé sur la force seule est fragile et illégitime : Hobbes lui-même montre que la crainte ne suffit pas, il faut une autorité reconnue.
\item \textbf{Droit :} c'est le fondement qui \emph{manque}. La « taxe » de 500 FCFA n'est pas un impôt légal (seule une loi peut créer un impôt) ; les « amendes » et « expulsions » n'ont aucune base juridique. Le comité est donc dans l'illégalité, même s'il est utile.
\end{itemize}
\textbf{Conclusion de l'étape 2 :} l'autorité du comité repose surtout sur le consentement et la tradition, mais l'absence de fondement juridique la rend précaire : que se passera-t-il le jour où une partie des habitants refusera de payer ou contestera une expulsion ?

\textbf{Étape 3 — Les rôles que l'État devrait assurer.}

\begin{center}
\begin{tabularx}{\linewidth}{|l|X|X|}\hline
\rowcolor{popblueL} \textbf{Rôle de l'État} & \textbf{Situation actuelle} & \textbf{Ce que l'État devrait faire} \\ \hline
Sécurité & 2 veilleurs payés par le comité ; commissariat à 14 km & Créer un poste de police ou une brigade de gendarmerie \\ \hline
Justice & Arbitrage informel ; expulsions arbitraires & Rendre la justice par des tribunaux compétents ; sécuriser les titres fonciers \\ \hline
Éducation & 62 \% des enfants scolarisés, écoles privées coûteuses & Construire une école publique et y affecter des enseignants \\ \hline
Santé & Hôpital à 17 km & Créer un centre de santé intégré \\ \hline
Voirie et eau & Pistes impraticables, forages privés & Rattacher le quartier à une commune, entretenir les voies, adduction d'eau \\ \hline
\end{tabularx}
\end{center}

\textbf{Étape 4 — Rédaction argumentée (modèle de copie).}

\emph{« Le comité de quartier peut-il durablement remplacer l'État ? »}

Le comité de Montée des Aigles rend de réels services : il entretient la piste grâce à 12,6 millions de FCFA par an, fait régner un certain ordre et arbitre les petits conflits. Son autorité repose en partie sur le consentement des habitants, ce qui, selon Locke, est la condition de tout pouvoir légitime. Peut-il pour autant remplacer l'État ? Nous soutenons que non, pour trois raisons.

D'abord, le comité ne possède pas la souveraineté : il ne crée pas le droit, sa « taxe » n'est pas un impôt légal, et ses sanctions (amendes, expulsions forcées) sont des actes de force, non de justice. Or, comme le montre Hobbes, seule une autorité souveraine, reconnue par tous et au-dessus de tous, peut garantir durablement la paix : sans elle, le jour où un groupe refuse les règles, le quartier retombe dans l'insécurité.

Ensuite, le comité ne peut pas assumer les rôles essentiels de l'État : il ne peut ni construire une école publique (38 \% des enfants ne sont pas scolarisés), ni créer un centre de santé, ni délivrer des titres fonciers alors que 38 litiges fonciers ont éclaté en 2024. Ces tâches exigent des moyens, des compétences et une légalité que seul l'État possède.

Enfin, la finalité même de l'État dépasse le comité : selon Rousseau, l'État doit réaliser la volonté générale, c'est-à-dire l'intérêt de tous et non celui d'un quartier isolé. Le comité défend les seuls habitants présents ; l'État garantit les mêmes droits à tous les citoyens, partout.

En conclusion, le comité est un palliatif utile, une initiative citoyenne à encourager, mais il ne peut pas durablement remplacer l'État : la solution n'est pas de se passer de l'État, mais d'exiger qu'il remplisse ses rôles à Montée des Aigles, en commençant par le rattachement du quartier à une commune.

\textbf{Commentaire de la copie :} notez la structure en quatre temps (thèse annoncée, trois arguments chacun appuyé sur un auteur et une donnée chiffrée, conclusion qui propose une solution). C'est exactement la démarche exigée au Probatoire / Baccalauréat technique.
\end{exoresolu}

\courssub{Bilan de l'activité}

Cette enquête a permis de mettre en évidence plusieurs acquis essentiels de la leçon :
\begin{itemize}
\item \textbf{Un pouvoir de fait n'est pas un État.} La grille des trois éléments constitutifs est un outil d'analyse rigoureux : le comité possède une population et un territoire approximatif, mais pas la souveraineté, qui est le critère décisif.
\item \textbf{La légitimité a plusieurs sources inégalement solides.} Le consentement et la tradition peuvent soutenir une autorité locale, mais sans le droit, elle reste fragile et peut basculer dans l'arbitraire (expulsions forcées).
\item \textbf{Les rôles de l'État sont irremplaçables.} Sécurité, justice, éducation, santé, infrastructures : aucune initiative privée ou communautaire ne peut les assumer durablement à grande échelle, faute de moyens et de légalité.
\item \textbf{La finalité de l'État éclaire le débat citoyen.} Le problème de Montée des Aigles n'est pas « l'État ou le comité », mais l'absence de l'État : la réponse philosophique consiste à réclamer la présence effective de l'État, garant de l'ordre (Hobbes), des droits (Locke) et du bien commun (Rousseau).
\item \textbf{Compétence transversale acquise :} les élèves ont rédigé et justifié une réponse argumentée, en articulant concepts, auteurs et données chiffrées — exactement la démarche de la dissertation philosophique attendue à l'examen.
\end{itemize}

\begin{savaistu}[Et au Cameroun ?]
Le rattachement des quartiers périurbains aux communes est un enjeu réel : les communes de Yaoundé et Douala s'étendent chaque année par intégration de nouveaux quartiers, ce qui permet l'arrivée de l'état civil, des écoles publiques et des centres de santé. Les comités de quartier, loin d'être des rivaux de l'État, deviennent alors des relais précieux entre les habitants et l'administration : la philosophie politique rejoint ici la vie quotidienne des citoyens camerounais.
\end{savaistu}

\newpage

\coursec{Méthodes et exercices résolus}

\coursec{L'État et le pouvoir}

\courssub{Savoir-faire 1 : Définir une notion et identifier ses éléments constitutifs}

\begin{methode}[Définir l'État et repérer ses éléments dans une situation]
\begin{enumerate}
\item \textbf{Rappeler la définition générale} : l'État est une organisation politique souveraine qui exerce un pouvoir légitime sur une population vivant sur un territoire délimité.
\item \textbf{Mémoriser les trois éléments constitutifs} (formule à retenir : \cle{Territoire -- Population -- Pouvoir souverain}). Sans l'un de ces trois éléments, il n'y a pas d'État.
\item \textbf{Distinguer l'État des notions voisines} : la \cle{nation} (communauté culturelle et historique), le \cle{gouvernement} (équipe dirigeante, temporaire), le \cle{régime politique} (forme d'organisation du pouvoir).
\item \textbf{Appliquer à la situation} : repérer dans le cas concret chacun des trois éléments (ex. : le Cameroun = territoire délimité par des frontières, population camerounaise, institutions souveraines : Présidence, Assemblée nationale, tribunaux).
\item \textbf{Conclure} en nommant précisément la notion et en justifiant la réponse par les éléments repérés.
\end{enumerate}
\end{methode}

\begin{exoresolu}[Identifier les éléments constitutifs de l'État]
\textbf{Énoncé.} À la suite d'une crise, un groupe armé occupe une région et proclame sur place un « nouvel État ». Il n'est reconnu par aucun pays, ne contrôle qu'une partie du territoire et n'a aucune administration stable. Ce groupe peut-il être considéré comme un État ? Justifiez votre réponse à partir de la définition de l'État.
\tcblower
\textbf{Solution détaillée.}

\textbf{Étape 1 : rappel de la définition.} Un État est une organisation politique souveraine reposant sur trois éléments constitutifs indissociables : un \cle{territoire} délimité, une \cle{population} et un \cle{pouvoir souverain} (autorité effective, reconnue et capable de faire appliquer le droit).

\textbf{Étape 2 : vérification élément par élément.}
\begin{itemize}
\item \emph{Territoire} : le groupe n'occupe qu'une partie d'une région, sans frontières stables ni reconnues. L'élément est donc incertain.
\item \emph{Population} : il n'existe pas de population constituée en corps politique acceptant cette autorité ; les habitants subissent l'occupation plutôt qu'ils ne reconnaissent un pouvoir.
\item \emph{Pouvoir souverain} : la souveraineté suppose une autorité effective et légitime, capable de faire des lois, de rendre la justice, de lever l'impôt. Or le groupe n'a aucune administration stable et n'est reconnu ni par la population ni par les autres États.
\end{itemize}

\textbf{Étape 3 : conclusion.} Aucun des trois éléments constitutifs n'est pleinement réuni : ce groupe armé n'est donc \textbf{pas un État}, mais une simple puissance de fait. Un État ne se proclame pas par la seule force : il suppose un territoire, une population et un pouvoir souverain effectivement exercé et reconnu.
\end{exoresolu}

\courssub{Savoir-faire 2 : Expliquer le fondement de l'obéissance au pouvoir}

\begin{methode}[Problématiser « pourquoi obéit-on à l'État ? »]
\begin{enumerate}
\item \textbf{Poser le paradoxe} : l'État commande et contraint ; pourquoi les hommes acceptent-ils d'obéir ? L'obéissance n'est pas naturelle.
\item \textbf{Mobiliser les trois grandes réponses} du programme :
\begin{itemize}
\item la \cle{force} (on obéit par peur : thèse de la domination) ;
\item le \cle{contrat social} (on obéit parce qu'on a consenti : Hobbes, Locke, Rousseau) ;
\item la \cle{légitimité} (on obéit parce que le pouvoir est reconnu comme juste et légal : Max Weber et ses trois types de domination légitime : traditionnelle, charismatique, rationnelle-légale).
\end{itemize}
\item \textbf{Comparer les thèses} : montrer ce que chacune explique et ce qu'elle laisse dans l'ombre (la force explique la soumission, pas la loyauté ; le contrat explique le consentement, mais il est une fiction).
\item \textbf{Illustrer} par un exemple concret (payer ses impôts, respecter le code de la route au Cameroun : crainte de la sanction ou acceptation de la règle ?).
\item \textbf{Conclure nuancée} : l'obéissance durable repose moins sur la peur que sur la légitimité.
\end{enumerate}
\end{methode}

\begin{exoresolu}[Dissertation guidée : obéit-on à l'État par peur ou par consentement ?]
\textbf{Énoncé.} « On n'obéit à l'État que parce qu'il est le plus fort. » Discutez cette affirmation en mobilisant les théories du fondement du pouvoir.
\tcblower
\textbf{Solution détaillée (plan et développement rédigé).}

\textbf{Introduction.} L'État dispose du monopole de la contrainte légitime : il peut punir celui qui viole la loi. Mais la force suffit-elle à fonder l'obéissance ? Si l'on n'obéissait que par peur, tout pouvoir serait instable. Il faut donc distinguer soumission et obéissance légitime.

\textbf{I. La thèse de la force : une explication partielle.}
Il est vrai que l'État possède des moyens de coercition (police, justice, prisons). La crainte de la sanction joue un rôle réel : beaucoup paient leurs impôts ou respectent les feux rouges par peur de l'amende. Mais un pouvoir fondé sur la seule force est un pouvoir fragile : dès que la force faiblit, l'obéissance cesse. La force produit la soumission, non la loyauté.

\textbf{II. Le contrat social : l'obéissance par consentement.}
Pour \cle{Hobbes} (\emph{Léviathan}), les hommes, pour échapper à la guerre de tous contre tous, abandonnent leur liberté naturelle à un souverain qui garantit la paix. Pour \cle{Rousseau} (\emph{Du contrat social}), obéir à la loi qu'on s'est prescrite, c'est obéir à la volonté générale et donc rester libre. L'obéissance est alors fondée sur le \cle{consentement} : j'obéis parce que la loi est aussi la mienne.

\textbf{III. La légitimité : la vraie force du pouvoir.}
\cle{Max Weber} montre que la domination durable repose sur la croyance en sa légitimité : traditionnelle (les chefs coutumiers), charismatique (le leader exceptionnel) ou rationnelle-légale (l'État de droit moderne, où l'on obéit à la loi et non aux personnes). Au Cameroun, l'obéissance aux institutions repose sur la Constitution et les procédures légales, non sur la seule contrainte.

\textbf{Conclusion.} L'affirmation est donc à rejeter : la force peut contraindre, mais seule la légitimité — issue du consentement et du droit — fonde une obéissance stable. L'État le plus durable n'est pas le plus fort, mais le plus reconnu.
\end{exoresolu}

\courssub{Savoir-faire 3 : Analyser une situation concrète à l'aide des rôles de l'État}

\begin{methode}[Appliquer les rôles de l'État à la vie courante]
\begin{enumerate}
\item \textbf{Lister les rôles de l'État} (formule à retenir : \cle{Sécurité -- Justice -- Services publics -- Économie -- Solidarité}) :
\begin{itemize}
\item assurer la \emph{sécurité} intérieure (police) et extérieure (armée) ;
\item rendre la \emph{justice} et garantir les droits ;
\item organiser les \emph{services publics} (écoles, hôpitaux, routes) ;
\item réguler l'\emph{économie} (impôts, monnaie, emploi) ;
\item assurer la \emph{solidarité} sociale (protection des plus démunis).
\end{itemize}
\item \textbf{Lire la situation} et y repérer l'action de l'État décrite.
\item \textbf{Nommer le rôle correspondant} et justifier : dire \emph{pourquoi} cette action relève de ce rôle.
\item \textbf{Distinguer} ce qui relève de l'État de ce qui relève des particuliers, des entreprises ou des associations.
\item \textbf{Rédiger la réponse} en une phrase complète : « Dans cette situation, l'État joue son rôle de... car... »
\end{enumerate}
\end{methode}

\begin{exoresolu}[Identifier les rôles de l'État dans des situations camerounaises]
\textbf{Énoncé.} Pour chacune des situations suivantes, identifiez le rôle de l'État mis en jeu et justifiez votre réponse :
(a) l'État construit un lycée technique à Bafoussam ; (b) un tribunal condamne un voleur à Yaoundé ; (c) l'armée protège les frontières de l'Extrême-Nord ; (d) l'État fixe le salaire minimum (SMIG) ; (e) un centre de santé public soigne gratuitement des enfants démunis.
\tcblower
\textbf{Solution détaillée.}

\textbf{(a) Construction d'un lycée technique à Bafoussam} : rôle de \cle{service public}. L'État organise l'éducation, service accessible à tous, car l'instruction est une condition de l'égalité des chances et du développement du pays.

\textbf{(b) Condamnation d'un voleur par un tribunal} : rôle \cle{judiciaire} (rendre la justice). L'État, par ses tribunaux, dit le droit, tranche les conflits et punit les infractions, garantissant ainsi l'ordre et les droits de chacun.

\textbf{(c) Protection des frontières par l'armée} : rôle de \cle{sécurité extérieure} (fonction régalienne). Défendre le territoire contre les menaces extérieures est la première mission de l'État, qui dispose seul de la force armée légitime.

\textbf{(d) Fixation du SMIG} : rôle de \cle{régulation économique}. L'État encadre les rapports entre employeurs et salariés pour éviter les abus et garantir un revenu minimal aux travailleurs.

\textbf{(e) Soins gratuits aux enfants démunis} : rôle de \cle{solidarité sociale}. L'État protège les plus faibles et corrige les inégalités, conformément à sa mission d'assurer la cohésion de la société.

\textbf{Conclusion.} Ces cinq situations montrent que l'État est présent dans tous les domaines de la vie quotidienne : sécurité, justice, éducation, économie et solidarité. Chaque action concrète peut être rattachée à l'un de ses rôles essentiels.
\end{exoresolu}

\courssub{Savoir-faire 4 : Rédiger et justifier une réponse argumentée sur la finalité de l'État}

\begin{methode}[Construire une réponse philosophique justifiée (méthode ARE)]
\begin{enumerate}
\item \textbf{Affirmer} (\emph{Affirmation}) : donner clairement sa thèse en une phrase (« La finalité de l'État est... »).
\item \textbf{Raisonner} (\emph{Raisonnement}) : justifier la thèse par un argument logique et un auteur (ex. : Rousseau — l'État vise le bien commun ; Locke — il existe pour protéger les droits naturels : vie, liberté, propriété ; Hegel — il réalise la liberté concrète).
\item \textbf{Exemplifier} (\emph{Exemple}) : appuyer l'argument sur un fait concret (une loi, une institution, une action de l'État camerounais).
\item \textbf{Envisager l'objection} : présenter brièvement la thèse adverse (l'État ne sert-il pas les intérêts des dominants ?) et y répondre.
\item \textbf{Conclure} en reformulant la thèse justifiée, sans introduire d'idée nouvelle.
\end{enumerate}
\end{methode}

\begin{exoresolu}[Rédiger un paragraphe argumenté sur la finalité de l'État]
\textbf{Énoncé.} En un paragraphe construit et justifié (thèse, argument, exemple, réponse à une objection), répondez à la question : l'État existe-t-il pour le bien commun ou pour les intérêts des gouvernants ?
\tcblower
\textbf{Solution détaillée.}

\textbf{Thèse.} La finalité véritable de l'État est le \cle{bien commun}, c'est-à-dire l'intérêt de tous les citoyens, et non l'intérêt particulier des gouvernants.

\textbf{Argument.} Selon \cle{Rousseau}, le pouvoir légitime tire son origine du contrat social : les individus s'unissent pour que soit réalisée la \emph{volonté générale}, qui vise l'intérêt de tous. Un pouvoir qui ne servirait que les gouvernants trahirait sa raison d'être et deviendrait illégitime. De même, \cle{Locke} soutient que l'État est institué pour protéger les droits naturels de tous — la vie, la liberté, la propriété — et que les gouvernants ne sont que des dépositaires révocables de ce pouvoir.

\textbf{Exemple.} Au Cameroun, la Constitution affirme que la souveraineté appartient au peuple ; les services publics (écoles, hôpitaux, routes) sont en principe destinés à tous les citoyens, quelle que soit leur région ou leur condition.

\textbf{Objection et réponse.} On objectera que, dans les faits, certains dirigeants utilisent l'État à leur profit (corruption, népotisme). Mais cet abus ne définit pas la finalité de l'État : il la trahit. C'est précisément parce que l'État a pour fin le bien commun que la corruption est condamnée par la loi et critiquée par l'opinion.

\textbf{Conclusion.} L'État existe donc par nature pour le bien commun ; lorsqu'il sert des intérêts particuliers, il se pervertit et perd sa légitimité.
\end{exoresolu}

\courssub{Savoir-faire 5 : Traiter un sujet de dissertation sur l'État (méthode complète)}

\begin{methode}[Réussir la dissertation au Probatoire/Baccalauréat technique]
\begin{enumerate}
\item \textbf{Analyser le sujet} (5 min) : souligner les mots-clés (État, pouvoir, obéir, liberté...), définir chacun, repérer la question posée.
\item \textbf{Dégager la problématique} : transformer le sujet en tension entre deux thèses (ex. : l'État protège la liberté / l'État menace la liberté).
\item \textbf{Construire un plan dialectique en deux ou trois parties} : thèse, antithèse, synthèse (ou dépassement).
\item \textbf{Rédiger l'introduction} en quatre temps : accroche, définition des termes, problématique, annonce du plan.
\item \textbf{Développer chaque partie} : une idée par paragraphe, un argument, un auteur, un exemple, une transition.
\item \textbf{Rédiger la conclusion} : réponse claire à la problématique, bilan, éventuelle ouverture courte.
\item \textbf{Relire} : orthographe, cohérence, présence des auteurs et exemples.
\end{enumerate}
\end{methode}

\begin{exoresolu}[Sujet type Bac : « L'État est-il un obstacle à la liberté ? »]
\textbf{Énoncé.} Traitez le sujet suivant sous forme de dissertation : « L'État est-il un obstacle à la liberté ? » (On donnera la problématique, le plan détaillé et les éléments de développement.)
\tcblower
\textbf{Solution détaillée.}

\textbf{Analyse du sujet.} Mots-clés : \cle{État} (pouvoir politique souverain), \cle{obstacle} (ce qui empêche), \cle{liberté} (pouvoir d'agir selon sa volonté). La question oppose apparemment l'autorité de l'État et la liberté individuelle.

\textbf{Problématique.} L'État, qui impose des lois contraignantes, semble limiter la liberté ; mais sans loi ni ordre, la liberté de chacun serait menacée par celle des autres. L'État détruit-il la liberté ou la rend-il possible ?

\textbf{Plan dialectique.}

\textbf{I. L'État comme obstacle à la liberté (thèse).}
Par ses lois, ses impôts, sa police, l'État contraint les individus : il interdit, oblige, punit. La force de l'État peut devenir tyrannique : les régimes dictatoriaux suppriment la liberté d'expression et d'opinion. On peut citer la critique anarchiste : toute autorité étatique est une oppression.

\textbf{II. L'État comme condition de la liberté (antithèse).}
Sans État, règne l'état de nature décrit par \cle{Hobbes} : « la guerre de tous contre tous », où la liberté du plus fort écrase celle du plus faible. Pour \cle{Rousseau}, « obéir à la loi qu'on s'est prescrite, c'est être libre » : la loi protège la liberté de tous en limitant celle de chacun. Exemple : le code pénal camerounais interdit la violence ; grâce à lui, chacun peut circuler, travailler, s'exprimer en sécurité.

\textbf{III. Synthèse : la liberté par le droit.}
L'État n'est un obstacle à la liberté que lorsqu'il est arbitraire et illégitime. L'\cle{État de droit}, au contraire, réalise la liberté : il garantit les droits de chacun par des lois égales pour tous. La vraie question n'est donc pas « État ou liberté ? », mais « quel État pour quelle liberté ? ».

\textbf{Conclusion.} L'État n'est pas par nature un obstacle à la liberté : il en est la condition lorsqu'il est fondé sur le droit et le consentement. Ce n'est que perverti en tyrannie qu'il la détruit.
\end{exoresolu}

\courssub{Savoir-faire 6 : Commenter une citation de philosophe sur le pouvoir}

\begin{methode}[Méthode du commentaire de texte ou de citation]
\begin{enumerate}
\item \textbf{Lire et situer} : identifier l'auteur et l'œuvre (ex. : Hobbes, \emph{Léviathan}).
\item \textbf{Dégager la thèse} : reformuler en une phrase ce que l'auteur affirme.
\item \textbf{Expliquer le raisonnement} : décomposer l'argumentation (d'où part l'auteur, que conclut-il, par quelles étapes ?).
\item \textbf{Définir les termes techniques} employés (souverain, contrat, volonté générale...).
\item \textbf{Illustrer} par un exemple contemporain ou camerounais.
\item \textbf{Discuter} : donner une limite ou une objection, puis une appréciation personnelle justifiée.
\end{enumerate}
\end{methode}

\begin{exoresolu}[Commenter une citation de Rousseau]
\textbf{Énoncé.} « Obéir à la loi qu'on s'est prescrite est la liberté. » (Rousseau, \emph{Du contrat social}.) Expliquez et discutez cette pensée.
\tcblower
\textbf{Solution détaillée.}

\textbf{Situation et thèse.} Dans le \emph{Contrat social} (1762), \cle{Rousseau} affirme que la liberté véritable n'est pas l'absence de lois, mais l'obéissance à des lois que le peuple s'est lui-même données.

\textbf{Explication du raisonnement.} Rousseau distingue deux libertés : la \emph{liberté naturelle} (faire tout ce qu'on peut, dans l'état de nature) et la \emph{liberté civile} (agir selon la loi commune). Par le \cle{contrat social}, chaque citoyen renonce à sa volonté particulière pour participer à la \cle{volonté générale}. La loi n'est alors plus un ordre étranger : elle est l'expression de ma propre volonté en tant que citoyen. Obéir à la loi, c'est donc obéir à soi-même — c'est cela, être libre.

\textbf{Illustration.} Le code de la route au Cameroun limite ma vitesse ; mais ces règles, votées par les représentants du peuple, protègent ma vie et celle des autres : en les respectant, je suis à la fois sujet et auteur de la loi.

\textbf{Discussion.} Cette thèse suppose que la loi exprime réellement la volonté générale. Or, si les lois sont faites par une minorité ou imposées sans consentement (dictature), obéir n'est plus être libre. La pensée de Rousseau est donc un idéal exigeant : elle fonde la démocratie, mais elle condamne par avance tout pouvoir qui ne consulte pas le peuple.

\textbf{Conclusion.} Pour Rousseau, la liberté ne s'oppose pas à la loi : elle s'accomplit par elle, à condition que la loi soit l'œuvre de tous. Cette idée reste le fondement de l'État démocratique moderne.
\end{exoresolu}

\begin{attention}[Les 5 erreurs qui coûtent des points au Bac]
\begin{enumerate}
\item \textbf{Confondre État, nation et gouvernement.} L'État est l'organisation politique permanente (territoire, population, pouvoir souverain) ; le gouvernement n'est que l'équipe dirigeante du moment ; la nation est une communauté culturelle. Employer l'un pour l'autre est une faute de définition éliminatoire.
\item \textbf{Affirmer que la force suffit à fonder l'État.} Erreur classique : « l'État, c'est le plus fort ». Un pouvoir fondé sur la seule force est instable ; il faut toujours distinguer \cle{puissance de fait} et \cle{pouvoir légitime} (Weber), et citer le contrat social (Hobbes, Rousseau).
\item \textbf{Citer des auteurs sans les expliquer.} Écrire « Rousseau dit que l'homme est bon » sans montrer le raisonnement ne rapporte rien. Chaque référence doit être suivie de l'argument et reliée à la problématique.
\item \textbf{Répondre hors sujet.} Beaucoup de copies récitent « tout le cours sur l'État » au lieu de répondre à la question posée. Il faut d'abord analyser les mots du sujet et construire une problématique ; toute partie qui ne répond pas à la question est sanctionnée.
\item \textbf{Négliger la justification et les exemples concrets.} Une réponse non justifiée (sans argument ni exemple) est considérée comme une opinion. Chaque affirmation doit être appuyée par un raisonnement et illustrée par un fait précis (institution, loi, situation camerounaise), puis conclue par une phrase de synthèse.
\end{enumerate}
\end{attention}

\newpage

\coursec{Exercices}

\coursec{L'État et le pouvoir : Exercices et applications}

\courssub{Je vérifie mes connaissances}

\begin{exoresolu}[QCM : les éléments constitutifs de l'État]
\textbf{Énoncé :} Pour chaque question, une seule réponse est correcte. Entoure la lettre correspondante.
\begin{enumerate}
\item Un État se définit par la réunion de trois éléments :
\begin{enumerate}
\item un roi, une armée, une religion ;
\item \textbf{un territoire, une population, un pouvoir souverain} ;
\item une capitale, un drapeau, une monnaie.
\end{enumerate}
\item La souveraineté de l'État signifie :
\begin{enumerate}
\item qu'il dépend d'une puissance étrangère ;
\item \textbf{qu'il détient le pouvoir suprême sur son territoire, sans reconnaître d'autorité au-dessus de lui} ;
\item que le chef de l'État est élu à vie.
\end{enumerate}
\item Selon Max Weber, l'État se caractérise par :
\begin{enumerate}
\item \textbf{le monopole de la violence physique légitime} ;
\item la possession du plus grand nombre de soldats ;
\item la richesse de son budget.
\end{enumerate}
\item Lequel de ces éléments ne relève pas des fonctions régaliennes de l'État ?
\begin{enumerate}
\item la justice ; \item la défense ; \item \textbf{l'organisation des loisirs des citoyens}.
\end{enumerate}
\end{enumerate}
\tcblower
\textbf{Solution détaillée :}
\begin{enumerate}
\item \textbf{b.} La définition classique (droit constitutionnel, sociologie politique) retient trois éléments constitutifs : un \cle{territoire} délimité, une \cle{population} stable, un \cle{pouvoir souverain} (gouvernement effectif). Les autres options confondent attributs (drapeau, capitale) ou forme de gouvernement (roi) avec la définition de l'État.
\item \textbf{b.} La \cle{souveraineté} (Bodin, droit international) est le pouvoir suprême, exclusif, inaliénable, imprescriptible, exercé sur un territoire sans supérieur hiérarchique externe.
\item \textbf{a.} Dans \emph{Le Savant et le Politique} (1919), Max Weber définit l'État comme « une communauté humaine qui, dans les limites d'un territoire déterminé, revendique avec succès le \cle{monopole de la violence physique légitime} ».
\item \textbf{c.} Les \cle{fonctions régaliennes} (justice, police, défense, diplomatie, monnaie) sont l'apanage exclusif de l'État. L'organisation des loisirs relève de la fonction sociale ou culturelle, souvent décentralisée ou associative.
\end{enumerate}
\end{exoresolu}

\begin{exoresolu}[Vrai ou faux justifié]
\textbf{Énoncé :} Indique si chaque affirmation est vraie ou fausse, puis justifie ta réponse en trois lignes maximum.
\begin{enumerate}
\item « Le gouvernement et l'État sont la même chose. »
\item « La force suffit à fonder un État durable. »
\item « L'État-providence se limite aux fonctions régaliennes (police, justice, armée, diplomatie). »
\item « La légitimité d'un pouvoir signifie qu'il est reconnu comme conforme à un droit ou à des valeurs acceptés par les citoyens. »
\end{enumerate}
\tcblower
\textbf{Solution détaillée :}
\begin{enumerate}
\item \textbf{Faux.} L'\cle{État} est une institution permanente (territoire, population, souveraineté), le \cle{gouvernement} en est l'organe exécutif temporaire et révocable (équipe ministérielle, majorité parlementaire). Confondre les deux conduit à l'absolutisme (« L'État, c'est moi »).
\item \textbf{Faux.} La \cle{force} (contrainte physique) peut imposer l'obéissance par la peur, mais elle ne crée pas de \cle{légitimité} (croyance en la justesse du pouvoir). Sans légitimité (traditionnelle, charismatique, rationnelle-légale chez Weber), le pouvoir est instable et coûteux à maintenir (répression permanente).
\item \textbf{Faux.} L'\cle{État-providence} (État social, post-1945) \emph{ajoute} aux fonctions régaliennes des \cle{fonctions sociales} (sécurité sociale, éducation, santé, retraites, allocation chômage) et \cle{fonctions économiques} (planification, services publics, régulation). Il vise la justice sociale et la correction des inégalités de marché.
\item \textbf{Vrai.} La \cle{légitimité} (Weber) est la « croyance en la légalité » ou la conformité à des principes supérieurs (droit, tradition, charisme). Elle transforme la contrainte en obligation morale, fondement de l'obéissance volontaire des citoyens.
\end{enumerate}
\end{exoresolu}

\begin{exoresolu}[Définitions]
\textbf{Énoncé :} Définis en trois à cinq lignes chacune des notions suivantes, en donnant pour chacune un exemple tiré de la vie camerounaise.
\begin{enumerate}
\item l'État ;
\item la souveraineté ;
\item la légitimité ;
\item l'État de droit.
\end{enumerate}
\tcblower
\textbf{Propositions de définitions (modèles) :}
\begin{enumerate}
\item \textbf{L'État :} Institution politique souveraine, dotée d'un territoire, d'une population et d'un pouvoir public monopolistique de la contrainte légitime (Weber), organisant la vie collective par le droit. \emph{Exemple :} La République du Cameroun, avec ses institutions (Présidence, Parlement, Gouvernement, Administration territoriale), son territoire (475 442 km²) et sa population (\textasciitilde 28 millions d'habitants).
\item \textbf{La souveraineté :} Pouvoir suprême, exclusif, inaliénable et imprescriptible, exercé par l'État sur son territoire et sa population, sans reconnaître de supérieur (indépendance externe) et s'imposant à tous en interne (Bodin, Constitution Art. 2). \emph{Exemple :} Le Cameroun décide seul de sa politique de défense, de sa monnaie (FCFA, bien que zone commune), de ses lois, sans ingérence étrangère directe.
\item \textbf{La légitimité :} Caractère de ce qui est fondé en droit, en justice ou en acceptation sociale ; pour Weber, croyance des gouvernés en la validité de l'ordre établi (légalité rationnelle, tradition, charisme), transformant la force en autorité. \emph{Exemple :} L'élection présidentielle de 2018 au Cameroun confère une légitimité rationnelle-légale au Président élu, reconnue par le Conseil constitutionnel, même si elle est contestée par une partie de l'opposition.
\item \textbf{L'État de droit :} État dont l'action est entièrement soumise au droit : séparation des pouvoirs, hiérarchie des normes (Constitution au sommet), garantie des libertés fondamentales, justice indépendante, administration contrôlée par le juge. \emph{Exemple :} Au Cameroun, le Conseil constitutionnel (créé 1996, réformé 2023) peut censurer une loi votée par le Parlement si elle viole la Constitution ; les justiciables peuvent saisir le Tribunal administratif contre l'État.
\end{enumerate}
\end{exoresolu}

\begin{exoresolu}[Repérage des rôles de l'État]
\textbf{Énoncé :} Pour chacune des situations suivantes, indique le rôle de l'État concerné (fonction régalienne, fonction économique, fonction sociale).
\begin{enumerate}
\item L'armée camerounaise sécurise la frontière dans l'Extrême-Nord.
\item L'État construit un lycée technique à Bafoussam.
\item Un tribunal de Yaoundé juge un litige entre deux commerçants.
\item L'État subventionne le prix des intrants agricoles pour les producteurs de maïs de l'Adamaoua.
\item La vaccination gratuite des enfants est organisée dans les centres de santé publics.
\end{enumerate}
\tcblower
\textbf{Solution détaillée :}
\begin{enumerate}
\item \textbf{Fonction régalienne} (Défense / Sécurité extérieure) : Monopole de la force armée pour protéger l'intégrité territoriale.
\item \textbf{Fonction sociale} (Éducation / Formation) : Service public d'enseignement, investissement dans le capital humain, réduction des inégalités d'accès au savoir.
\item \textbf{Fonction régalienne} (Justice) : Monopole de la contrainte juridictionnelle, règlement pacifique des conflits, application du droit.
\item \textbf{Fonction économique} (Agriculture / Soutien à la production) : Interventionnisme étatique pour corriger les défaillances du marché, soutenir un secteur stratégique, assurer la sécurité alimentaire.
\item \textbf{Fonction sociale} (Santé publique / Protection sociale) : Accès universel aux soins préventifs, politique de santé publique, solidarité nationale.
\end{enumerate}
\end{exoresolu}

\courssub{Je m'entraîne}

\begin{exercice}[Étude de cas : deux pouvoirs, deux légitimités]
Dans un village de la région de l'Ouest, un chef traditionnel impose une amende à un habitant qui a enfreint une coutume. Le même mois, à Douala, un commissaire de police dresse une contravention contre un automobiliste.
\begin{enumerate}
\item Compare la nature des deux pouvoirs en présence (origine, étendue, fondement juridique).
\item Sur quels fondements repose la légitimité de chacun (légitimité traditionnelle, légitimité rationnelle-légale) selon la typologie de Max Weber ?
\item Montre, en t'appuyant sur Max Weber, pourquoi seul l'État détient le monopole de la contrainte légitime sur l'ensemble du territoire national, et quelle est la place du chef traditionnel dans l'ordre juridique camerounais.
\end{enumerate}
\end{exercice}

\begin{exoresolu}[Analyse de données : le budget de l'État]
\textbf{Énoncé :} Voici la répartition simplifiée (fictive) des dépenses de l'État du Cameroun pour une année, en pourcentage du budget total :

\begin{center}
\begin{tabularx}{\linewidth}{|l|X|}\hline
\rowcolor{popblueL} \textbf{Secteur} & \textbf{Part du budget (\%)} \\ \hline
Éducation et formation & 20 \\ \hline
Santé & 10 \\ \hline
Défense et sécurité & 15 \\ \hline
Infrastructures et équipement & 25 \\ \hline
Dette et charges de l'État & 20 \\ \hline
Agriculture et développement rural & 10 \\ \hline
\end{tabularx}
\end{center}

\begin{enumerate}
\item Construis un diagramme circulaire représentant cette répartition (calcule l'angle de chaque secteur : rappel, 1\,\% correspond à 3,6 degrés).
\item Quels sont les deux postes de dépenses les plus importants ?
\item Rédige cinq lignes de conclusion : quelles priorités de l'État camerounais ce budget révèle-t-il, et quelle finalité de l'État illustre-t-il ?
\end{enumerate}
\tcblower
\textbf{Solution détaillée :}
\begin{enumerate}
\item \textbf{Calcul des angles (arrondis au degré) :}
\begin{itemize}
\item Éducation : $20 \times 3,6 = 72^\circ$
\item Santé : $10 \times 3,6 = 36^\circ$
\item Défense : $15 \times 3,6 = 54^\circ$
\item Infrastructures : $25 \times 3,6 = 90^\circ$
\item Dette : $20 \times 3,6 = 72^\circ$
\item Agriculture : $10 \times 3,6 = 36^\circ$
\item \textbf{Total : } $360^\circ$.
\end{itemize}
\textbf{Diagramme circulaire (TikZ) :}
\begin{popfigure}
\begin{tikzpicture}
% Couleurs
\definecolor{colEdu}{HTML}{1F77B4} % popblue
\definecolor{colHealth}{HTML}{2CA02C} % popgreen
\definecolor{colDef}{HTML}{D62728} % poporange/red
\definecolor{colInfra}{HTML}{FF7F0E} % poporange
\definecolor{colDebt}{HTML}{9467BD} % poppurple
\definecolor{colAgri}{HTML}{8C564B} % brown

% Secteurs (sens horaire, départ 90° = nord)
% 1. Education 72° (90 -> 18)
\fill[colEdu] (0,0) -- (90:3) arc (90:18:3) -- cycle;
% 2. Santé 36° (18 -> -18 / 342)
\fill[colHealth] (0,0) -- (18:3) arc (18:-18:3) -- cycle;
% 3. Défense 54° (-18 -> -72 / 288)
\fill[colDef] (0,0) -- (-18:3) arc (-18:-72:3) -- cycle;
% 4. Infrastructures 90° (-72 -> -162 / 198)
\fill[colInfra] (0,0) -- (-72:3) arc (-72:-162:3) -- cycle;
% 5. Dette 72° (-162 -> -234 / 126)
\fill[colDebt] (0,0) -- (-162:3) arc (-162:-234:3) -- cycle;
% 6. Agriculture 36° (-234 -> -270 / 90)
\fill[colAgri] (0,0) -- (-234:3) arc (-234:-270:3) -- cycle;

% Légende
\begin{scope}[xshift=4cm, yshift=1.5cm]
\node[anchor=west] at (0,1.5) {\textbf{Légende :}};
\node[anchor=west, fill=colEdu] at (0,1.0) {\rule{0.5cm}{0.3cm} Éducation (20\,\%)};
\node[anchor=west, fill=colHealth] at (0,0.5) {\rule{0.5cm}{0.3cm} Santé (10\,\%)};
\node[anchor=west, fill=colDef] at (0,0.0) {\rule{0.5cm}{0.3cm} Défense (15\,\%)};
\node[anchor=west, fill=colInfra] at (0,-0.5) {\rule{0.5cm}{0.3cm} Infrastructures (25\,\%)};
\node[anchor=west, fill=colDebt] at (0,-1.0) {\rule{0.5cm}{0.3cm} Dette (20\,\%)};
\node[anchor=west, fill=colAgri] at (0,-1.5) {\rule{0.5cm}{0.3cm} Agriculture (10\,\%)};
\end{scope}
\end{tikzpicture}
\legende{Répartition fictive du budget de l'État du Cameroun (en \%). Source : Données d'exercice.}
\end{popfigure}

\item \textbf{Les deux postes les plus importants :} \textbf{Infrastructures et équipement (25\,\%)} et \textbf{Éducation et formation (20\,\%)} (ex-æquo avec Dette et charges de l'État à 20\,\%).

\item \textbf{Conclusion :} Ce budget révèle une priorité forte pour l'\textbf{investissement physique} (infrastructures : routes, barrages, bâtiments administratifs) et le \textbf{capital humain} (éducation), piliers de la stratégie « Cameroun émergent à l'horizon 2035 ». Le poids de la \textbf{dette (20\,\%)} marque la contrainte financière limitant les marges de manœuvre. L'État illustre ici sa \cle{finalité de développement} (fonction économique) et de \cle{promotion sociale} (éducation, santé), au-delà de ses seules fonctions régaliennes (défense 15\,\%), confirmant sa nature d'État-providence en construction.
\end{enumerate}
\end{exoresolu}

\begin{exercice}[Question à argumentation construite : la Constitution camerounaise]
Lis attentivement cet extrait simplifié de la Constitution de la République du Cameroun (Loi n° 96-06 du 18 janvier 1996) : le préambule proclame l'attachement du peuple camerounais aux libertés fondamentales inscrites dans la Déclaration universelle des droits de l'homme de 1948 et la Charte africaine des droits de l'homme et des peuples de 1981 ; l'article premier affirme que le Cameroun est « une République décentralisée, unie et indivisible, laïque, démocratique et sociale », dont la devise est « Paix -- Travail -- Patrie », et que la souveraineté nationale appartient au peuple qui l'exerce par ses représentants ou par voie de référendum.
\begin{enumerate}
\item Dégage la conception de l'État qui s'exprime dans ce texte (État de droit, souveraineté populaire, laïcité, décentralisation).
\item Confronte cette conception à la thèse de Hobbes (\emph{Léviathan}), pour qui la paix civile exige que les hommes abandonnent leur droit naturel à un souverain tout-puissant, non responsable devant les sujets.
\item En six à huit lignes, indique quelle conception te semble la plus proche de la vie politique camerounaise actuelle, en justifiant ta réponse par des faits précis (institutions, pratique du pouvoir, droits de l'homme).
\end{enumerate}
\end{exercice}

\begin{exercice}[Application des notions à une situation concrète]
Lors d'une élection présidentielle au Cameroun, une partie significative de la population conteste les résultats proclamés par le Conseil constitutionnel et refuse de reconnaître le vainqueur.
\begin{enumerate}
\item Distingue, dans cette situation, la \emph{légalité} (conformité aux procédures juridiques en vigueur) et la \emph{légitimité} (acceptation morale, adhésion des citoyens) du pouvoir contesté.
\item Explique pourquoi, selon Rousseau (\emph{Du contrat social}), un pouvoir ne peut durer que s'il repose sur la \cle{volonté générale} et non sur la seule force ou la légalité formelle. Cite la formule : « Le plus fort n'est jamais assez fort pour être toujours le maître, s'il ne transforme sa force en droit et l'obéissance en devoir. »
\item Propose, en cinq lignes, une conclusion personnelle et justifiée sur les conditions d'un pouvoir pleinement accepté par les citoyens (transparence, alternance, respect des libertés, justice sociale).
\end{enumerate}
\end{exercice}

\courssub{Je me prépare au Bac}

\begin{exercice}[Dissertation : « L'État est-il un mal nécessaire ? »]
\textbf{Sujet :} L'État est-il un mal nécessaire ?

\textbf{Travail demandé :}
\begin{enumerate}
\item \textbf{Analyse le sujet :} dégage le sens des termes « État » (institution souveraine), « mal » (atteinte à la liberté, contrainte, violence, inégalité), « nécessaire » (indispensable à la survie, condition de possibilité de la vie commune), puis formule la \cle{problématique} (ex. : La contrainte étatique est-elle une fatalité réductrice de liberté ou le fondement même de la liberté civile ?).
\item \textbf{Rédige l'introduction complète} de la dissertation (accroche philosophique ou historique, définition rigoureuse des termes, problématique, annonce du plan dialectique en trois temps).
\item \textbf{Construis le plan détaillé en trois parties :}
\begin{itemize}
\item \textbf{Thèse 1 : L'État est un mal.} Il limite la liberté naturelle, use de la contrainte (police, impôts, justice), peut devenir tyrannique ou totalitaire. Appuie-toi sur la critique libérale (Locke : pouvoir limité, droit de propriété antérieur à l'État) ou anarchiste (Proudhon, Bakounine : « L'État, c'est la négation de la liberté », Kropotkine : entraide volontaire vs contrainte étatique).
\item \textbf{Thèse 2 : L'État est nécessaire.} Sans lui, c'est l'état de nature : insécurité, « guerre de tous contre tous » (Hobbes, \emph{Léviathan}), impossibilité de l'industrie, du commerce, de la culture. Le contrat social fonde la sécurité et la paix comme conditions de l'existence même. L'État est le \cle{mal nécessaire} (Hobbes) ou le \cle{passage de la liberté naturelle à la liberté civile} (Rousseau).
\item \textbf{Synthèse : L'État est un bien \emph{s'il est limité par le droit et légitime}.} Ni mal absolu, ni bien absolu : l'\cle{État de droit} (Montesquieu, Kelsen, Constitution camerounaise) soumet le pouvoir à la loi, sépare les pouvoirs, garantit les droits fondamentaux. La finalité de l'État devient le \cle{bien commun} et la \cle{protection de la liberté} (Rousseau : « L'obéissance à la loi qu'on s'est prescrite, c'est la liberté »).
\end{itemize}
\item \textbf{Illustre chaque partie par un exemple concret tiré de la vie camerounaise :} contrôles de police routiers (contrainte/mal), paiement de l'impôt (contrainte/nécessaire), écoles publiques gratuites (service/bien), protection des populations contre Boko Haram (sécurité/nécessaire), indépendance relative du judiciaire (État de droit/bien).
\end{enumerate}
\end{exercice}

\begin{exercice}[Explication de texte : Rousseau, \emph{Du contrat social}]
\textbf{Texte :} « L'homme est né libre, et partout il est dans les fers. Tel se croit le maître des autres, qui ne laisse pas d'être plus esclave qu'eux. Comment ce changement s'est-il fait ? Je l'ignore. Qu'est-ce qui peut le rendre légitime ? Je crois pouvoir résoudre cette question. » (Rousseau, \emph{Du contrat social}, Livre I, Chapitre 1, 1762)

\textbf{Travail demandé :}
\begin{enumerate}
\item Présente l'auteur et l'œuvre en deux lignes (contexte : Lumières, philosophie politique, contratualisme).
\item Explique le sens de la formule « né libre » (liberté naturelle, indépendance, absence de subordination) et celui de l'expression « dans les fers » (dépendance sociale, inégalités factices, aliénation, servitude volontaire).
\item Dégage la thèse de Rousseau : la liberté est naturelle à l'homme, mais la vie en société telle qu'elle existe historiquement l'enchaîne ; le problème philosophique central est de savoir si cette sujétion peut devenir \cle{légitime} (juste, consentie).
\item Montre comment cette question annonce la solution du \cle{contrat social} : « Trouver une forme d'association qui défende et protège de toute la force commune la personne et les biens de chaque associé, et par laquelle chacun, s'unissant à tous, n'obéisse pourtant qu'à lui-même et reste aussi libre qu'auparavant. » (Livre I, Ch. 6). Obéir à des lois que le peuple s'est données (volonté générale), c'est rester libre (autonomie = auto-nomie).
\item \textbf{Discussion :} La soumission aux lois de l'État camerounais (Code pénal, Code du travail, fiscalité, Code électoral) est-elle compatible avec la liberté des citoyens au sens de Rousseau ? Justifie ta réponse en cinq à huit lignes en distinguant \cle{liberté naturelle} (illimitée, destructrice), \cle{liberté civile} (garantie par la loi), et en évaluant la participation réelle du peuple camerounais à l'élaboration de la loi (représentation, référendum, société civile).
\end{enumerate}
\end{exercice}

\begin{exercice}[Sujet de synthèse : pour quoi l'État existe-t-il ?]
Un élève de Première déclare : « L'État ne sert qu'à nous commander et à prélever des impôts ; on vivrait mieux sans lui. »

\textbf{Travail demandé :}
\begin{enumerate}
\item Reformule la position de cet élève et identifie la conception de l'État qu'elle suppose (vision anarchiste, libérale radicale ou simple mécontentement citoyen ; confusion entre État et gouvernement, oubli des services rendus).
\item Rappelle les \textbf{trois grands rôles de l'État} (fonctions régaliennes, économiques, sociales) en illustrant \textbf{chacun par un exemple camerounais précis et daté si possible} (ex. : Régalienne -- Opération \emph{Épervier} contre la corruption / Sécurité Extrême-Nord ; Économique -- Projet barrage de Nachtigal, Zones économiques spéciales ; Sociale -- Gratuité de l'enseignement primaire public, Couverture santé universelle CSU).
\item Explique, en t'appuyant sur Hobbes (\emph{Léviathan}, Ch. 13-17), ce que serait une société sans État (\cle{état de nature} : droit de tout sur tout, insécurité permanente, « vie solitaire, pauvre, nasty, brutish, and short » / « guerre de tous contre tous »), et pourquoi la peur de la mort violente pousse les hommes à instituer un pouvoir absolu.
\item Montre, en t'appuyant sur Rousseau (\emph{Du contrat social}, Livre I, Ch. 6-9), que la finalité de l'État n'est pas la domination mais la \cle{protection de la liberté et du bien commun} : transformation de la liberté naturelle en liberté civile, de la force en droit, de l'instinct en justice ; l'État est l'artificiel qui rend l'homme vraiment humain et moral.
\item Rédige une \textbf{conclusion personnelle argumentée (huit à dix lignes)} répondant à la question : l'État existe-t-il pour dominer ou pour servir les citoyens ? Nuance : l'État \emph{peut} dominer (dérive autoritaire, corruption), mais sa \cle{finalité normative} (sa raison d'être philosophique et constitutionnelle) est le service du bien commun ; la vigilance citoyenne (contrôle, vote, presse libre, société civile) est la condition pour qu'il reste serviteur et non maître.
\end{enumerate}
\end{exercice}

\newpage

\coursec{Corrigés des exercices}

\begin{corrige}[Exercice 5 — Étude de cas : deux pouvoirs, deux légitimités]

\textbf{1. Comparaison des deux pouvoirs.}

\begin{center}
\begin{tabularx}{\linewidth}{|l|X|X|}\hline
\rowcolor{popblueL} \textbf{Critère} & \textbf{Chef traditionnel} & \textbf{Commissaire de police} \\ \hline
Origine du pouvoir & Coutume, hérédité, sacralité (investiture traditionnelle) & Délégation de l'État, nomination administrative, concours \\ \hline
Étendue & Limitée au village et aux affaires coutumières (statut personnel, terres, rites) & S'étend à tout le territoire national, à tous les citoyens \\ \hline
Fondement juridique & Coutume reconnue (droit coutumier toléré si conforme à la loi) & Loi, décrets, Code de procédure pénale, Code de la route \\ \hline
\end{tabularx}
\end{center}

\textbf{2. Fondements de la légitimité (typologie de Weber).} Le chef traditionnel relève de la \cle{légitimité traditionnelle} : on lui obéit parce que « cela a toujours été ainsi », par attachement à la coutume et au lignage. Le commissaire relève de la \cle{légitimité rationnelle-légale} : on lui obéit non à sa personne, mais à la fonction définie par des règles impersonnelles (la loi). Aucun des deux ne relève ici de la légitimité charismatique (fondée sur les qualités exceptionnelles d'un leader).

\textbf{3. Le monopole étatique et la place du chef traditionnel.} Selon Weber, l'État est la seule instance qui « revendique avec succès le monopole de la violence physique légitime » sur l'ensemble d'un territoire. Le chef traditionnel ne peut donc exercer une contrainte que dans les limites que l'État lui concède : au Cameroun, les chefferies sont reconnues par l'administration (décret de 1977 sur l'organisation des chefferies traditionnelles), mais leurs décisions ne peuvent contrevenir ni à la Constitution ni aux lois ; une amende coutumière abusive peut être contestée devant les tribunaux de l'État. Le pouvoir coutumier est ainsi \emph{subordonné} au pouvoir étatique : il survit comme relais administratif et culturel, non comme souverain concurrent.

\textbf{Points de vigilance :} ne pas opposer les deux pouvoirs comme « archaïque » contre « moderne » (jugement de valeur à éviter) ; citer explicitement Weber ; bien montrer que la légitimité traditionnelle est réelle mais \emph{locale et déléguée}.

\end{corrige}

\begin{corrige}[Exercice 7 — Question à argumentation construite : la Constitution camerounaise]

\textbf{1. La conception de l'État exprimée.} Le texte définit un \cle{État de droit} : le pouvoir est limité par des normes supérieures (droits de l'homme de 1948, Charte africaine de 1981). La \cle{souveraineté populaire} (le peuple exerce le pouvoir par ses représentants ou par référendum) s'oppose à la souveraineté monarchique. La \cle{laïcité} sépare l'État des religions ; le caractère « social » inscrit une finalité de justice sociale ; la \cle{décentralisation} répartit le pouvoir entre l'État central et les collectivités territoriales.

\textbf{2. Confrontation avec Hobbes.} Pour Hobbes (\emph{Léviathan}, 1651), la paix exige que chacun transfère son droit naturel à un souverain \emph{absolu et non responsable} : le contrat lie les sujets entre eux, mais le souverain n'est jamais lié. La Constitution camerounaise repose sur la logique inverse : le pouvoir \emph{émane du peuple}, est \emph{limité} par les droits fondamentaux et se \emph{partage} (décentralisation). On passe du contrat d'aliénation (Hobbes) au contrat de délégation (Locke, Rousseau).

\textbf{3. Exemple de développement personnel (6-8 lignes).} La vie politique camerounaise actuelle se situe entre les deux modèles. Formellement, elle relève de l'État de droit : Constitution écrite, Conseil constitutionnel, élections pluralistes, presse privée, collectivités décentralisées (régions instituées en 2019). Mais la pratique conserve des traits hobbesiens : forte concentration du pouvoir exécutif, longévité du chef de l'État au pouvoir, administration centralisée. La conception constitutionnelle est donc un \emph{idéal normatif} que la réalité n'atteint qu'imparfaitement : l'État de droit est moins un acquis qu'un processus à consolider par la vigilance citoyenne et l'indépendance effective des institutions.

\textbf{Points de vigilance :} citer le texte (devoir ancré sur le document) ; ne pas confondre souveraineté \emph{nationale} et \emph{populaire} ; éviter le pamphlet : argumenter avec des faits institutionnels précis.

\end{corrige}

\begin{corrige}[Exercice 8 — Application : légalité et légitimité d'une élection contestée]

\textbf{1. Distinction légalité / légitimité.} Le vainqueur proclamé par le Conseil constitutionnel possède la \cle{légalité} : sa victoire est conforme aux procédures juridiques en vigueur (proclamation officielle, délais de recours respectés). Mais si une partie significative de la population estime le scrutin entaché d'irrégularités, sa \cle{légitimité} est fragilisée : l'adhésion morale fait défaut. Un pouvoir peut donc être légal sans être pleinement légitime ; inversement, un opposant peut jouir d'une légitimité populaire sans légalité.

\textbf{2. L'analyse rousseauiste.} Rousseau écrit : « Le plus fort n'est jamais assez fort pour être toujours le maître, s'il ne transforme sa force en droit et l'obéissance en devoir. » La force nue (contrôle de l'appareil d'État, de l'armée) ne fonde rien de durable : elle n'engendre que la soumission par la peur, que la révolte démentira dès que possible. Seule la \cle{volonté générale} — la loi voulue par le peuple souverain — transforme la contrainte en obligation morale. Un pouvoir contesté dans sa légitimité doit donc la reconstruire par le consentement, non par la répression.

\textbf{3. Exemple de conclusion personnelle (5 lignes).} Un pouvoir pleinement accepté suppose la transparence électoral (publication des procès-verbaux, observateurs indépendants), l'alternance possible, le respect effectif des libertés d'expression et de manifestation, et une justice sociale qui donne à chacun des raisons d'adhérer à l'ordre commun. La légalité est la condition minimale ; la légitimité, qui se mérite chaque jour, en est le couronnement.

\textbf{Points de vigilance :} ne jamais confondre les deux notions dans la copie (erreur la plus fréquente) ; la citation de Rousseau doit être exacte et \emph{expliquée}, pas simplement posée ; rester dans l'analyse philosophique sans prendre parti politiquement.

\end{corrige}

\begin{corrige}[Exercice 9 — Dissertation : « L'État est-il un mal nécessaire ? »]

\textbf{1. Analyse du sujet.} \cle{État} : institution souveraine détenant le monopole de la contrainte légitime sur un territoire (Weber). \cle{Mal} : ce qui nuit — ici la contrainte, la limitation des libertés, la violence possible. \cle{Nécessaire} : ce dont on ne peut se passer. Le sujet contient une tension : un « mal » devrait être supprimé, un « nécessaire » doit être conservé. \textbf{Problématique :} la contrainte étatique est-elle une atteinte regrettable à la liberté, ou la condition même de la liberté civile ?

\textbf{2. Exemple d'introduction rédigée.} En 2011, l'effondrement de l'État somalien a montré ce qu'est une société sans pouvoir central : milices, piraterie, insécurité permanente. Faut-il pour autant bénir toute autorité étatique ? L'État, institution souveraine détenant le monopole de la violence légitime, apparaît comme un « mal » dès lors qu'il contraint, taxe, punit ; mais comme « nécessaire » dès lors qu'aucune vie commune ordonnée ne semble possible sans lui. La contrainte étatique est-elle alors une fatalité qui réduit la liberté, ou le fondement paradoxal de la liberté civile ? Nous verrons d'abord que l'État peut être pensé comme un mal (I), puis qu'il est une nécessité pour sortir de l'état de nature (II), enfin qu'il devient un bien lorsqu'il est limité par le droit (III).

\textbf{3. Plan détaillé.}
\begin{itemize}
\item \textbf{I. L'État comme mal.} La contrainte (impôts, police, justice) limite la liberté naturelle ; le pouvoir tend à s'étendre et peut devenir tyrannique. Locke : le pouvoir doit être strictement limité, la propriété est antérieure à l'État. Proudhon : « Être gouverné, c'est être noté, enregistré, imposé, commandé ». \emph{Exemple camerounais :} contrôles routiers multiples, lourdeurs fiscales perçues par les commerçants.
\item \textbf{II. L'État comme nécessité.} Hobbes : à l'état de nature, « guerre de tous contre tous », vie « solitaire, pauvre, sordide, abrutie et brève » ; seul le Léviathan garantit la paix. Rousseau : le contrat transforme la liberté naturelle en liberté civile. \emph{Exemples :} sécurisation de l'Extrême-Nord contre Boko Haram ; l'impôt finance les écoles et les hôpitaux.
\item \textbf{III. Synthèse : un bien s'il est limité et légitime.} L'État de droit (Montesquieu : séparation des pouvoirs ; Kelsen : hiérarchie des normes) soumet l'État à sa propre loi. La finalité de l'État est le bien commun. Rousseau : « L'obéissance à la loi qu'on s'est prescrite est liberté. » \emph{Exemple :} gratuité de l'école primaire publique, recours possibles devant le juge administratif.
\end{itemize}

\textbf{Barème indicatif (sur 20) :} analyse du sujet et problématique : 4 pts ; introduction : 3 pts ; cohérence dialectique du plan : 5 pts ; exactitude des références (Hobbes, Locke, Rousseau) : 4 pts ; exemples camerounais pertinents : 2 pts ; langue et présentation : 2 pts.

\textbf{Points de vigilance :} ne pas traiter « mal » et « nécessaire » comme deux questions séparées mais comme une tension unique ; éviter le catalogue d'auteurs sans argumentation ; la synthèse ne doit pas être un simple compromis mou mais dépasser l'opposition (l'État de droit comme solution).

\end{corrige}

\begin{corrige}[Exercice 10 — Explication de texte : Rousseau, \emph{Du contrat social}]

\textbf{1. Présentation.} Jean-Jacques Rousseau (1712-1778), philosophe des Lumières, publie \emph{Du contrat social} en 1762 : cet ouvrage de philosophie politique cherche à fonder la légitimité de l'obéissance politique sur le consentement, dans la tradition contractualiste.

\textbf{2. Explication des formules.} « Né libre » : par nature, l'homme n'est subordonné à personne ; la liberté est son état originaire, aucune autorité n'a de fondement naturel (ni droit divin, ni force). « Partout dans les fers » : dans les sociétés historiques réelles, les hommes sont enchaînés par la dépendance, les inégalités factices, les conventions injustes — y compris les maîtres, esclaves de leurs esclaves.

\textbf{3. La thèse.} Rousseau distingue deux questions : le fait de la sujétion (« comment ce changement s'est-il fait ? Je l'ignore ») et son droit (« qu'est-ce qui peut le rendre légitime ? »). Le problème central est donc normatif : à quelle condition l'obéissance à un pouvoir peut-elle être \emph{juste} ? Ni la force ni la tradition ne suffisent : seul le consentement fonde l'obligation.

\textbf{4. Vers le contrat social.} La solution (Livre I, ch. 6) : « trouver une forme d'association [...] par laquelle chacun, s'unissant à tous, n'obéisse pourtant qu'à lui-même et reste aussi libre qu'auparavant ». En obéissant à la \cle{volonté générale} — la loi que le peuple se donne —, le citoyen n'obéit qu'à lui-même : c'est l'\cle{autonomie}. La liberté civile remplace avantageusement la liberté naturelle.

\textbf{5. Discussion (5-8 lignes).} Au sens de Rousseau, obéir aux lois camerounaises est compatible avec la liberté à condition que ces lois expriment réellement la volonté générale : lois votées par un Parlement élu, participation électorale effective, possibilité de référendum, liberté de presse et de la société civile. La liberté naturelle (faire tout ce qu'on veut) est illimitée et destructrice ; la liberté civile (obéir à la loi commune) est la seule liberté réelle en société. Mais si la loi n'exprime que la volonté d'un pouvoir sans participation populaire authentique, elle redevient des « fers » : la formule de Rousseau fonctionne alors comme un idéal critique permettant de juger la réalité institutionnelle.

\textbf{Barème indicatif (sur 20) :} présentation auteur/œuvre : 2 pts ; explicitation des formules : 4 pts ; dégagement de la thèse : 4 pts ; lien avec le contrat social : 4 pts ; discussion personnelle argumentée : 4 pts ; langue : 2 pts.

\textbf{Points de vigilance :} suivre l'ordre du texte (méthode de l'explication) ; ne pas paraphraser ; la discussion doit mobiliser la distinction liberté naturelle / liberté civile et rester philosophique, non polémique.

\end{corrige}

\begin{corrige}[Exercice 11 — Sujet de synthèse : pour quoi l'État existe-t-il ?]

\textbf{1. Reformulation.} L'élève tient un discours de type \emph{anarchisant} ou libéral radical : l'État serait pure contrainte (commandement) et pure ponction (impôt), sans contrepartie. Cette position repose sur deux confusions : entre l'État (institution permanente) et le gouvernement (équipe dirigeante contestable), et entre le coût visible de l'État (impôts, contrôles) et ses services moins visibles (sécurité, écoles, routes, santé).

\textbf{2. Les trois rôles de l'État illustrés.}
\begin{itemize}
\item \textbf{Fonctions régaliennes :} l'armée sécurise l'Extrême-Nord face à Boko Haram ; la justice poursuit les infractions (opérations anticorruption) ; la police assure l'ordre public.
\item \textbf{Fonctions économiques :} construction du barrage hydroélectrique de Nachtigal (mis en service en 2024), routes nationales, soutien à la production agricole.
\item \textbf{Fonctions sociales :} gratuité de l'enseignement primaire public (depuis 2000), subvention des examens officiels, campagnes de vaccination gratuites, projet de couverture santé universelle.
\end{itemize}

\textbf{3. La société sans État selon Hobbes.} À l'\cle{état de nature} (\emph{Léviathan}, ch. 13), chacun a droit sur tout, personne n'est assez fort pour dominer durablement : règne la « guerre de tous contre tous », et la vie est « solitaire, pauvre, sordide, abrutie et brève ». Sans pouvoir commun qui inspire la crainte, ni industrie, ni agriculture, ni savoir ne sont possibles. C'est la peur de la mort violente et le désir de conservation qui poussent les hommes à instituer par pacte un souverain doté d'un pouvoir suffisant pour imposer la paix.

\textbf{4. La finalité de l'État selon Rousseau.} Pour Rousseau, le pacte légitime ne fait pas de l'État un maître mais un serviteur du \cle{bien commun} : l'homme y échange la liberté naturelle (illimitée, fondée sur la force) contre la liberté civile (garantie par la loi) et la propriété légitime ; « l'obéissance à la loi qu'on s'est prescrite est liberté ». La finalité de l'État est donc morale : transformer l'instinct en justice, faire de l'individu un citoyen.

\textbf{5. Exemple de conclusion personnelle (8-10 lignes).} L'État existe en droit pour servir, même s'il peut en fait dominer. Son histoire et sa pratique montrent des dérives réelles : autoritarisme, corruption, contraintes excessives — et le discours de l'élève exprime un mécontentement compréhensible. Mais l'absence d'État, loin d'être la liberté, serait selon Hobbes l'insécurité généralisée où la liberté de chacun s'annule. La solution n'est donc ni la suppression de l'État ni sa soumission aveugle : c'est l'État de droit, où le pouvoir est limité par la Constitution, contrôlé par le juge et par les citoyens (vote, presse libre, société civile), et orienté vers le bien commun. L'État est un serviteur exigeant : il ne reste serviteur que si les citoyens exercent leur vigilance. Sa finalité normative — la sécurité, la justice, la liberté civile — justifie son existence ; sa vigilante limitation justifie notre exigence.

\textbf{Barème indicatif (sur 20) :} reformulation et diagnostic : 3 pts ; trois rôles illustrés d'exemples camerounais précis : 4 pts ; analyse de Hobbes : 4 pts ; analyse de Rousseau : 4 pts ; conclusion personnelle nuancée et argumentée : 3 pts ; langue et organisation : 2 pts.

\textbf{Points de vigilance :} exiger des exemples \emph{précis} (noms, dates) et non des généralités ; la conclusion doit trancher avec nuance (distinction fait/droit, domination possible/finalité de service) ; vérifier que Hobbes et Rousseau ne sont pas confondus (souverain absolu chez l'un, volonté générale chez l'autre).

\end{corrige}

\newpage

\coursec{Fiche bilan}

\begin{popfigure}
\begin{center}
\begin{tikzpicture}[mindmap, grow cyclic, every node/.style={concept, font=\small},
  root concept/.append style={concept color=popblue, font=\bfseries\small, minimum size=3.2cm},
  level 1/.append style={level distance=4.1cm, sibling angle=60},
  level 1 concept/.append style={font=\footnotesize, minimum size=2.4cm}]
\node[root concept]{L'\'Etat et le pouvoir}
  child[concept color=poporange]{ node {D\'efinition : territoire, population, pouvoir souverain} }
  child[concept color=popgreen]{ node {Fondements : force, droit, contrat social, l\'egitimit\'e} }
  child[concept color=popgold]{ node {Ob\'eissance : contrainte ou consentement ?} }
  child[concept color=poppurple]{ node {R\^oles : s\'ecurit\'e, justice, \'education, \'economie, social} }
  child[concept color=poppink]{ node {Finalit\'e : ordre, bien commun, libert\'e, justice} }
  child[concept color=popblueL!80!popblue]{ node {Auteurs : Hobbes, Locke, Rousseau, Weber} };
\end{tikzpicture}
\end{center}
\legende{Carte mentale de la le\c{c}on : autour de la notion centrale d'\'Etat, six branches essentielles (d\'efinition, fondements, probl\`eme de l'ob\'eissance, r\^oles, finalit\'e, auteurs de r\'ef\'erence).}
\end{popfigure}

\begin{bilan}
\textbf{D\'efinitions cl\'es}
\begin{itemize}
\item \cle{\'Etat} : forme d'organisation politique d'une soci\'et\'e, reposant sur trois \'el\'ements constitutifs : un \cle{territoire} d\'elimit\'e par des fronti\`eres, une \cle{population} (la nation) et un \cle{pouvoir souverain} qui s'exerce par des institutions.
\item \cle{Souverainet\'e} : puissance supr\^eme de l'\'Etat, qui ne reconna\^it aucune autorit\'e au-dessus de lui (interne : lois valables pour tous ; externe : ind\'ependance face aux autres \'Etats).
\item \cle{Pouvoir} : capacit\'e d'imposer des d\'ecisions ; il est \cle{l\'egitime} lorsqu'il est reconnu comme conforme au droit et accept\'e par les gouvern\'es.
\item \cle{Monopole de la violence l\'egitime} (M. Weber) : l'\'Etat est la seule instance autoris\'ee \`a user l\'egalement de la contrainte physique (police, arm\'ee, justice).
\item \cle{Contrat social} : convention par laquelle les individus renoncent \`a une partie de leur libert\'e naturelle en \'echange de la s\'ecurit\'e et des droits garantis par l'\'Etat.
\item \cle{Bien commun} : ce qui profite \`a tous les membres de la soci\'et\'e et que nul ne peut atteindre seul (paix, justice, services publics).
\end{itemize}

\textbf{Th\`eses des auteurs \`a conna\^itre par c\oe ur}
\begin{center}
\begin{tabularx}{\linewidth}{|l|X|X|}
\hline
\rowcolor{popblueL} \textbf{Auteur} & \textbf{Th\`ese} & \textbf{Formule rep\`ere} \\
\hline
Hobbes & L'\'Etat fort (L\'eviathan) met fin \`a la guerre de tous contre tous & \og L'homme est un loup pour l'homme \fg \\
\hline
Locke & L'\'Etat prot\`ege les droits naturels (vie, libert\'e, propri\'et\'e) ; pouvoir limit\'e & Droit de r\'esistance \`a la tyrannie \\
\hline
Rousseau & La souverainet\'e appartient au peuple ; la loi exprime la volont\'e g\'en\'erale & \og Ob\'eir \`a la loi qu'on s'est prescrite est libert\'e \fg \\
\hline
Weber & L'\'Etat d\'etient le monopole de la violence physique l\'egitime & Trois l\'egitimit\'es : traditionnelle, charismatique, rationnelle-l\'egale \\
\hline
\end{tabularx}
\end{center}

\textbf{M\'ethodes express}
\begin{itemize}
\item \emph{Dissertation} : d\'efinir l'\'Etat d\`es l'introduction, probl\'ematiser (l'\'Etat est-il au service de l'homme ou l'homme au service de l'\'Etat ?), confronter deux th\`eses (Hobbes contre Locke/Rousseau), conclure avec nuance.
\item \emph{Explication de texte} : situer l'auteur, d\'egager la th\`ese, suivre l'argumentation pas \`a pas, discuter.
\item \emph{Application concr\`ete} : relier chaque notion \`a une situation camerounaise (police et gendarmerie = monopole de la force ; \'ecoles et h\^opitaux publics = r\^ole social ; \'elections = l\'egitimit\'e rationnelle-l\'egale).
\end{itemize}
\end{bilan}

\begin{aretenir}[Checklist avant le Bac]
\begin{itemize}
\item[$\square$] Je sais d\'efinir l'\'Etat et citer ses trois \'el\'ements constitutifs.
\item[$\square$] Je distingue clairement \emph{pouvoir de fait} et \emph{pouvoir l\'egitime}.
\item[$\square$] Je connais la d\'efinition de la souverainet\'e (interne et externe).
\item[$\square$] Je peux expliquer la th\`ese de Hobbes (L\'eviathan, \'etat de nature).
\item[$\square$] Je peux expliquer la th\`ese de Locke (droits naturels, pouvoir limit\'e).
\item[$\square$] Je peux expliquer la th\`ese de Rousseau (volont\'e g\'en\'erale, souverainet\'e populaire).
\item[$\square$] Je sais d\'efinir le monopole de la violence l\'egitime selon Weber.
\item[$\square$] Je sais r\'epondre \`a la question : \og pourquoi ob\'eit-on au pouvoir ? \fg (contrainte et consentement).
\item[$\square$] Je peux \'enum\'erer les r\^oles de l'\'Etat (s\'ecurit\'e, justice, \'education, \'economie, protection sociale).
\item[$\square$] Je peux discuter la finalit\'e de l'\'Etat : ordre, bien commun, libert\'e, justice.
\item[$\square$] Je sais illustrer chaque notion par un exemple camerounais pr\'ecis.
\item[$\square$] Je sais construire une probl\'ematique et r\'ediger une introduction de dissertation.
\end{itemize}
\end{aretenir}

\findecours

\end{document}
